% 高一98_市西中学_10月月考.tex
%   —— 高一上数学“2026-2027 年市西中学高一上 10 月月考”（讲评版）
% 试卷版：xelatex 高一98_市西中学_10月月考.tex
% 答案版：xelatex "\def\isanswer{1}% 高一98_市西中学_10月月考.tex
%   —— 高一上数学“2026-2027 年市西中学高一上 10 月月考”（讲评版）
% 试卷版：xelatex 高一98_市西中学_10月月考.tex
% 答案版：xelatex "\def\isanswer{1}% 高一98_市西中学_10月月考.tex
%   —— 高一上数学“2026-2027 年市西中学高一上 10 月月考”（讲评版）
% 试卷版：xelatex 高一98_市西中学_10月月考.tex
% 答案版：xelatex "\def\isanswer{1}% 高一98_市西中学_10月月考.tex
%   —— 高一上数学“2026-2027 年市西中学高一上 10 月月考”（讲评版）
% 试卷版：xelatex 高一98_市西中学_10月月考.tex
% 答案版：xelatex "\def\isanswer{1}\input{高一98_市西中学_10月月考.tex}"
% 紧凑版：xelatex "\def\iscompact{1}\input{高一98_市西中学_10月月考.tex}"
%
% 原始材料：微信公众号“魔都数学加油站”文章，连载“27学年高一98”；
% 原文标题“27学年高一98——2026-2027年上海市市西中学高一上10月月考”；
% 原文链接：https://mp.weixin.qq.com/s/duK4wz-HT3W9tjOeYQaVfg。
% 正文为 8 张 PNG 页（`原图/高一98_市西中学_10月月考/01..08.png`），
% 尺寸混排（第 1—4、6、7 页为 4762×6735，第 5、8 页为 2560×3620），按页序读取。
%
% 篇幅与题号：原卷自第 3 题起编号（第 1、2 题不在原文内）。填空题 3—12（10 题）、
% 选择题 13—16（4 题）、解答题 17—21（5 题），共 19 题。
%
% 讲评版：全部 19 题均印【解析】，其中第 13—16 题的答案在解析末句
% “故选 $B$／$D$／$A$／$D$”（题干括号内未印答案）。全卷无插图。
%
% 副标题：原卷未印副标题。本稿按内容标作“集合与常用逻辑用语、不等式”。
%
% 原卷笔误/存疑（均照录未改，正文对应处另有 % 注）：
%   ① 第 21 题 (1) 求证行、解析“假设”行与“故”行三处把 $x_3$ 误印作 $x_2$
%      （“$x_1$、$x_2$、$x_2$、$x_4$、$x_5$”）；解析同一行后半句
%      “这与 $x_1+x_2+x_3+x_4+x_5=5$ 矛盾”却用 $x_3$，显系笔误，照录；
%   ② 第 6 题解析的判别式行印作“$\Delta=(-4)^2-4a\geqslant0$”，与
%      $f(1)>0$ 并列，照录。
%
% 与原卷的排版差异：
%   ① 原卷每题下直接印【解析】，本稿按全稿惯例将答案与解析置于答案版；
%   ② 原卷未印副标题，本稿补出；填空横线统一排 \blankline[4.4em]；
%   ③ 原卷数集记号排斜体 $R$、$\Z$，本稿按全稿惯例改用 $\R$、$\Z$；
%   ④ 解答题按惯例加 \ansspace；选择题括号排 \mbox{（\quad）}；
%   ⑤ 第 12 题定义的运算 $\odot$ 原卷排实心小圆，本稿排 $\odot$。
%
% 行数核对：逐题数 \step/\stepm 条数，合计 82 条（Q3：1；Q4：2；Q5：1；Q6：4；Q7：3；
% Q8：2；Q9：5；Q10：2；Q11：5；Q12：4；Q13：3；Q14：2；Q15：2；Q16：9；Q17：2；
% Q18：11；Q19：6；Q20：11；Q21：7）。全卷 19 题均印【解析】。
% 原卷解析的换行按语意合并，未逐行强拆。
%
\documentclass[a4paper,11pt]{article}
\usepackage{common}
\begin{document}

\examtitlest[7]{2026--2027 年市西中学高一上 10 月月考}{集合与常用逻辑用语、不等式}

\bigsec{一、填空题}

\begin{enumerate}
\setcounter{enumi}{2}

\item 集合 $A=\{x\mid y=\sqrt{1-x^2}\},B=\{y\mid y=x^2+1,x\in A\}$，则 $A\cap\overline B=\blankline[4.4em]$.

\ifanswer
\solbegin
\step{$1-x^2\geqslant0\Rightarrow A=[-1,1]$，从而 $B=[1,2]$，所以 $A\cap\overline B=[-1,1)$.}
\solend
\fi

\item 已知 $b,c\in R$，关于 $x$ 的不等式 $x^2-bx+c<0$ 的解集为 $(-3,2)$，则 $b+c=\blankline[4.4em]$.

\ifanswer
\solbegin
\step{由题意得 $x=-3$ 与 $x=2$ 是 $x^2-bx+c=0$ 的两根，}
\step{所以 $\begin{cases}-3+2=b\\ -3\times2=c\end{cases}$，解得 $\begin{cases}b=-1\\ c=-6\end{cases}$，所以 $b+c=-7$.}
\solend
\fi

\item 若 $\alpha:\lvert x-m\rvert<1$ 是 $\beta:x^2-4x<0$ 的充分条件，则实数 $m$ 的取值范围是\blankline[4.4em].

\ifanswer
\solbegin
\step{$(m-1,m+1)\subseteq(0,4)\Leftrightarrow m-1\geqslant0$ 且 $m+1\leqslant4\Leftrightarrow m\in[1,3]$.}
\solend
\fi

\item 已知方程 $x^2-4x+a=0$ 的两根都大于 $1$，则 $a$ 的取值范围为\blankline[4.4em].

\ifanswer
\solbegin
\step{设函数 $f(x)=x^2-4x+a$，对称轴为 $x=2$，}
\step{因为方程 $x^2-4x+a=0$ 的两根都大于 $1$，}
\step{所以 $\begin{cases}\Delta=(-4)^2-4a\geqslant0\\ f(1)>0\end{cases}$，即 $\begin{cases}16-4a\geqslant0\\ 1-4+a>0\end{cases}$，解得 $3<a\leqslant4$，}
\step{即 $a$ 的取值范围为 $(3,4]$.}
\solend
\fi

\item 若不等式 $(a-2)x^2+2(a-2)x-4<0$ 对 $x\in R$ 恒成立，则实数 $a$ 的取值范围为\blankline[4.4em].

\ifanswer
\solbegin
\step{①当 $a-2=0$，即 $a=2$ 时，不等式为 $-4<0$ 恒成立，符合题意；}
\step{②当 $a\neq2$ 时，则 $\begin{cases}a-2<0\\ \Delta=4(a-2)^2+16(a-2)<0\end{cases}$，解得 $-2<a<2$.}
\step{综上所述，实数 $a$ 的取值范围是 $(-2,2]$.}
\solend
\fi

\item 已知 $-1\leqslant a+b\leqslant4,2\leqslant a-b\leqslant3$，则 $3a-2b$ 的取值范围为\blankline[4.4em].

\ifanswer
\solbegin
\step{设 $3a-2b=m(a+b)+n(a-b)$，则 $m+n=3,m-n=-2$，}
\step{所以 $m=\dfrac12,n=\dfrac52$，所以 $3a-2b=\dfrac12(a+b)+\dfrac52(a-b)\in\left[\dfrac92,\dfrac{19}2\right]$.}
\solend
\fi

\item 设 $x_1$、$x_2$ 是关于 $x$ 的方程 $2x^2-4mx+2m^2+3m-2=0$ 的两个实数根，则 $\lvert x_1-x_2\rvert$ 的最小值为\blankline[4.4em].

\ifanswer
\solbegin
\step{$x_1$、$x_2$ 是方程 $2x^2-4mx+2m^2+3m-2=0$ 的两个实数根，}
\step{所以 $\Delta=(4m)^2-8(2m^2+3m-2)\geqslant0$，解得 $m\leqslant\dfrac23$，}
\step{由韦达定理得 $x_1+x_2=2m,x_1\cdot x_2=\dfrac{2m^2+3m-2}2$，}
\step{则 $\lvert x_1-x_2\rvert=\sqrt{(x_1+x_2)^2-4x_1\cdot x_2}=\sqrt{4-6m}$，}
\step{当 $m=\dfrac23$ 时，$\lvert x_1-x_2\rvert$ 的最小值为 $0$.}
\solend
\fi

\item 设 $x\in R$，$[x]$ 表示不大于 $x$ 的最大整数，如 $[\pi]=3$，$[-1.2]=-2$，$\left[\dfrac12\right]=0$，则使 $\lvert\lvert x^2-1\rvert\rvert=3$ 成立的 $x$ 的取值范围是\blankline[4.4em].

\ifanswer
\solbegin
\step{若 $\lvert\lvert x^2-1\rvert\rvert=3$，则 $3\leqslant\lvert x^2-1\rvert<4$，即 $4\leqslant x^2<5$，}
\step{解得 $2\leqslant x<\sqrt5$ 或 $-\sqrt5<x\leqslant-2$，故 $x$ 的取值范围是 $(-\sqrt5,-2]\cup[2,\sqrt5)$.}
\solend
\fi

\item 设 $a$ 为实数，若关于 $x$ 的一元一次不等式组 $\begin{cases}2x+a>0\\ 3x-6a<0\end{cases}$ 的解集中有且仅有 $4$ 个整数，则 $a$ 的取值范围是\blankline[4.4em].

\ifanswer
\solbegin
\step{一元一次不等式组 $\begin{cases}2x+a>0\\ 3x-6a<0\end{cases}\Leftrightarrow-\dfrac a2<x<2a$，且 $a>0$，}
\step{关于 $x$ 的一元一次不等式组 $\begin{cases}2x+a>0\\ 3x-6a<0\end{cases}$ 的解集中有且仅有 $4$ 个整数，}
\step{①当 $4$ 个整数为 $0$，$1$，$2$，$3$ 时，则 $\begin{cases}-1\leqslant-\dfrac a2<0\\ 3<2a\leqslant4\end{cases}$，所以 $\dfrac32<a\leqslant2$，}
\step{②当 $4$ 个整数为 $-1$，$0$，$1$，$2$ 时，则 $\begin{cases}-2\leqslant-\dfrac a2<-1\\ 2<2a\leqslant3\end{cases}$，无解，}
\step{综上，$a$ 的取值范围是 $\left(\dfrac32,2\right]$.}
\solend
\fi

\item 集合 $M$、$N$、$S$ 都是非空集合，现规定如下运算：

$M\odot N\odot S=\{x\mid x\in(M\cap N)\cup(N\cap S)\cup(S\cap M)$ 且 $x\notin M\cap N\cap S\}$．假设集合 $A=\{x\mid a<x<b\}$，$B=\{x\mid c<x<d\}$，$C=\{x\mid e<x<f\}$，其中实数 $a$、$b$、$c$、$d$、$e$、$f$ 满足：$a<c<e<0<b<d<f$．计算 $A\odot B\odot C=\blankline[4.4em]$.

\ifanswer
\solbegin
\step{因为集合 $A=\{x\mid a<x<b\}$，$B=\{x\mid c<x<d\}$，$C=\{x\mid e<x<f\}$，}
\step{又因为 $a<c<e<0<b<d<f$，}
\step{所以 $A\cap B=\{x\mid c<x<b\}$，$B\cap C=\{x\mid e<x<d\}$，$C\cap A=\{x\mid e<x<b\}$，}
\step{$A\cap B\cap C=\{x\mid e<x<b\}$，故 $A\odot B\odot C=\{x\mid c<x\leqslant e$ 或 $b\leqslant x<d\}$.}
\solend
\fi

\end{enumerate}

\bigsec{二、选择题}

\begin{enumerate}
\setcounter{enumi}{12}

\item “$x\neq1$”是“$x^2-1\neq0$”的\mbox{（\quad）}条件.

\keepwithprev
A. 充分非必要\qquad B. 必要非充分

C. 充要\qquad D. 既非充分也非必要

\ifanswer
\solbegin
\step{$x^2-1\neq0\Leftrightarrow x\neq1$ 且 $x\neq-1\Rightarrow x\neq1$.}
\step{反之不成立，如 $x=-1$ 时，$x\neq1$，但 $x^2-1=0$.}
\step{因此“$x\neq1$”是“$x^2-1\neq0$”的必要非充分条件，故选 $B$.}
\solend
\fi

\item 若 $\dfrac1a<\dfrac1b<0$，则下列结论中不正确的是\mbox{（\quad）}

\keepwithprev
A. $a+b<0$\qquad B. $a^2<b^2$\qquad C. $ab<b^2$\qquad D. $ab<a^2$

\ifanswer
\solbegin
\step{由 $\dfrac1a<\dfrac1b<0$ 得 $b<a<0$，因此 $a+b<0$，且 $a^2<ab<b^2$，}
\step{所以 A、B、C 正确，D 不正确．故选 $D$.}
\solend
\fi

\item 非空集合 $X=\{x\mid a+1\leqslant x\leqslant3a-5\}$，$Y=\{x\mid1\leqslant x\leqslant16\}$，使得 $X\subseteq(X\cap Y)$ 成立的所有 $a$ 的集合是\mbox{（\quad）}

\keepwithprev
A. $\{a\mid3\leqslant a\leqslant7\}$\qquad B. $\{a\mid0\leqslant a\leqslant7\}$\qquad C. $\{a\mid3<a\leqslant7\}$\qquad D. $\{a\mid a\leqslant7\}$

\ifanswer
\solbegin
\step{因为 $X\subseteq(X\cap Y)$ 成立，所以 $X\subseteq Y$ 成立.}
\step{$X\neq\varnothing$，要使 $X\subseteq Y$，则 $\begin{cases}a\geqslant3\\ a+1\geqslant1\\ 3a-5\leqslant16\end{cases}$，即 $\begin{cases}a\geqslant3\\ a\geqslant0\\ a\leqslant7\end{cases}$，所以 $3\leqslant a\leqslant7$，故选 $A$.}
\solend
\fi

\item 已知 $M$ 是满足下列条件的集合：①$0\in M$，$1\in M$；②若 $x$、$y\in M$，则 $x-y\in M$；③若 $x\in M$ 且 $x\neq0$，则 $\dfrac1x\in M$．则下列说法正确的个数为\mbox{（\quad）}

（1）$\dfrac12\in M$；（2）$x+y\in M$；（3）$x^2\in M$．

\keepwithprev
A. $0$\qquad B. $1$\qquad C. $2$\qquad D. $3$

\ifanswer
\solbegin
\step{（1）$\dfrac12\in M$ 正确，理由如下：}
\step{由①得 $0\in M$，$1\in M$，由②得 $0-1=-1\in M$，$1-(-1)=2\in M$，}
\step{由③得 $\dfrac12\in M$；}
\step{（2）由①得 $0\in M$，由题意 $y\in M$，}
\step{所以由②得 $0-y=-y\in M$，又 $x\in M$，所以 $x-(-y)=x+y\in M$；}
\step{（3）$x\in M$，由②得 $x-1\in M$，由③得 $\dfrac1x\in M$，$\dfrac1{x-1}\in M$，}
\step{所以 $\dfrac1x-\dfrac1{x-1}\in M$，即 $\dfrac1{x(x-1)}\in M$，所以 $x(x-1)\in M$，}
\step{由（2）结论得 $x(x-1)+x\in M$，即 $x^2\in M$，}
\step{所以说法正确的个数为 $3$．故选 $D$.}
\solend
\fi

\end{enumerate}

\bigsec{三、解答题}

\begin{enumerate}
\setcounter{enumi}{16}

\item 已知 $a>b$，用比较法证明：$a^3>b^3$.

\ansspace*[8]

\ifanswer
\solbegin
\step{$a>b$，$a^3-b^3=(a-b)(a^2+ab+b^2)=(a-b)\left[\left(a+\dfrac b2\right)^2+\dfrac34b^2\right]>0$，}
\step{故 $a^3>b^3$.}
\solend
\fi

\item 求下列关于 $x$ 的不等式的解集：

（1）$\dfrac{x^2(x+3)(x^2-x+1)}{x^2-5x+4}>0$；

（2）$\dfrac{ax-1}{x-2}\geqslant0\left(a<\dfrac12\right)$.

\ansspace*[12]

\ifanswer
\solbegin
\stepm{（1）}{分母 $x^2-5x+4=(x-1)(x-4)$，要求 $x\neq1,4$；}
\step{又 $x=0$ 时原不等式不成立，且 $x^2-x+1=\left(x-\dfrac12\right)^2+\dfrac34>0$，}
\step{因此原不等式等价于 $\dfrac{x+3}{(x-1)(x-4)}>0,x\neq0$.}
\step{在数轴上标出 $-3,1,4$，分段判号得 $-3<x<1$ 或 $x>4$，}
\step{再排除 $x=0$，解集为 $(-3,0)\cup(0,1)\cup(4,+\infty)$.}
\stepm{（2）}{原不等式要求 $x\neq2$，按 $a$ 的符号分类讨论.}
\step{①当 $a<0$ 时，$\dfrac1a<2$，原不等式等价于 $\dfrac{x-\dfrac1a}{x-2}\leqslant0\Leftrightarrow\dfrac1a\leqslant x<2$，}
\step{解集为 $\left[\dfrac1a,2\right)$.}
\step{②当 $a=0$ 时，原不等式为 $\dfrac{-1}{x-2}\geqslant0\Leftrightarrow x<2$，解集为 $(-\infty,2)$.}
\step{③当 $0<a<\dfrac12$ 时，$\dfrac1a>2$，}
\step{原不等式等价于 $\dfrac{x-\dfrac1a}{x-2}\geqslant0\Leftrightarrow x<2$ 或 $x\geqslant\dfrac1a$，解集为 $(-\infty,2)\cup\left[\dfrac1a,+\infty\right)$.}
\solend
\fi

\item 集合 $A=\{x\mid-2\leqslant x\leqslant5\}$，$B=\{x\mid m+1\leqslant x\leqslant2m-1\}$.

（1）若 $B\subseteq A$，求实数 $m$ 的取值范围；

（2）若 $A\cap B\neq\varnothing$，求实数 $m$ 的取值范围.

\ansspace*[10]

\ifanswer
\solbegin
\stepm{（1）}{当 $B=\varnothing$ 时，$m+1>2m-1$，解得 $m<2$，}
\step{当 $B\neq\varnothing$ 时，$\begin{cases}m+1\leqslant2m-1\\ m+1\geqslant-2\\ 2m-1\leqslant5\end{cases}$，解得 $2\leqslant m\leqslant3$，}
\step{综上，实数 $m$ 的取值范围是 $(-\infty,3]$；}
\stepm{（2）}{当 $B=\varnothing$ 时，$m+1>2m-1$，解得 $m<2$，}
\step{当 $B\neq\varnothing$ 时，$\begin{cases}m+1\leqslant2m-1\\ 2m-1<-2\end{cases}$ 或 $\begin{cases}m+1\leqslant2m-1\\ m+1>5\end{cases}$，解得 $m\in\varnothing$ 或 $m>4$.}
\step{所以当 $A\cap B\neq\varnothing$ 时，实数 $m$ 的取值范围是 $[2,4]$.}
\solend
\fi

\item 在解决实际问题时，往往会有不同的思路和方法，这些方法有些正确，有些错误；有些简洁，有些复杂.

问题①：设 $a\in R$，集合 $M=\{x\mid(x-1)(x-a)=0\}$，若 $x\in M$ 是 $x\in\{1,2,3\}$ 的充分条件，求：$a$ 的取值集合.

问题②：设 $x\in Z$，$a=x+3$，$b=18-x$，若 $c\in Z$，求证：$c-a$ 和 $c-b$ 至少有一个数是奇数

（1）小明在解决问题①，他认为原问题等价于 $\{1,a\}\subseteq\{1,2,3\}$，解得 $a$ 的取值集合为 $\{2,3\}$，张老师判断小明解题错误，请计算正确的 $a$ 的取值集合；

（2）小红认为既然 $c\in Z$，只需根据 $c$ 是奇数还是偶数，分类讨论即可；小华则认为可以使用反证法解决问题，请你选一种你认为更好的方法并证明.

\ansspace*[12]

\ifanswer
\solbegin
\stepm{（1）}{当 $a=1$ 时，集合 $M=\{x\mid(x-1)^2=0\}=\{1\}$，}
\step{此时 $x\in M\Rightarrow x\in\{1,2,3\}$，即 $x\in M$ 是 $x\in\{1,2,3\}$ 的充分条件.}
\step{当 $a\neq1$ 时，集合 $M=\{x\mid(x-1)(x-a)=0\}=\{1,a\}$，}
\step{由 $x\in M$ 是 $x\in\{1,2,3\}$ 的充分条件，得 $\{1,a\}\subseteq\{1,2,3\}$，}
\step{所以 $a=2$ 或 $a=3$.}
\step{综上，实数 $a$ 的取值集合为 $\{1,2,3\}$；}
\stepm{（2）}{利用反证法证明．证明如下：}
\step{假设 $c-a$、$c-b$ 均为偶数，则 $(c-b)-(c-a)=a-b$ 为偶数.}
\step{因为 $a-b=x+3-18+x=2x-15$ 为奇数，}
\step{与 $b-a$ 为偶数相矛盾，所以原假设不成立.}
\step{故 $c-a$、$c-b$ 至少有一个数是奇数.}
\solend
\fi

\item 已知实数 $x_1$、$x_2$、$x_3$、$x_4$、$x_5$，满足 $x_1+x_2+x_3+x_4+x_5=5$.

（1）求证：$x_1$、$x_2$、$x_2$、$x_4$、$x_5$ 中至少有一个实数不小于 $1$；% 原卷把 x_3 误印作 x_2，照录

（2）设 $x_1$、$x_2$、$x_3$、$x_4$、$x_5$ 这五个实数两两不等，集合 $A=\{x_1,x_2,x_3,x_4,x_5\}$，若 $B\neq\varnothing$ 且 $B\subseteq A$，记 $G(B)$ 是 $B$ 中所有元素之和，对所有的 $B$，求 $G(B)$ 的平均值.

\ansspace*[13]

\ifanswer
\solbegin
\stepm{（1）}{假设 $x_1$、$x_2$、$x_2$、$x_4$、$x_5$ 中的每一个数都小于 $1$，}% 原卷 x_3 误印作 x_2，照录
\step{得 $x_1+x_2+x_2+x_4+x_5<5$，这与 $x_1+x_2+x_3+x_4+x_5=5$ 矛盾，}% 前半句原卷 x_3 误印作 x_2，照录
\step{故 $x_1$、$x_2$、$x_2$、$x_4$、$x_5$ 中至少有一个实数不小于 $1$；}% 原卷 x_3 误印作 x_2，照录
\stepm{（2）}{集合 $A=\{x_1,x_2,x_3,x_4,x_5\}$ 的非空子集个数为 $32-1=31$，}
\step{由于 $G(B)$ 是 $B$ 中所有元素之和，}
\step{得 $G(B)=16(x_1+x_2+x_3+x_4+x_5)=16\times5=80$，}
\step{则 $G(B)$ 的平均值为 $\dfrac{80}{31}$.}
\solend
\fi

\end{enumerate}

\end{document}
"
% 紧凑版：xelatex "\def\iscompact{1}% 高一98_市西中学_10月月考.tex
%   —— 高一上数学“2026-2027 年市西中学高一上 10 月月考”（讲评版）
% 试卷版：xelatex 高一98_市西中学_10月月考.tex
% 答案版：xelatex "\def\isanswer{1}\input{高一98_市西中学_10月月考.tex}"
% 紧凑版：xelatex "\def\iscompact{1}\input{高一98_市西中学_10月月考.tex}"
%
% 原始材料：微信公众号“魔都数学加油站”文章，连载“27学年高一98”；
% 原文标题“27学年高一98——2026-2027年上海市市西中学高一上10月月考”；
% 原文链接：https://mp.weixin.qq.com/s/duK4wz-HT3W9tjOeYQaVfg。
% 正文为 8 张 PNG 页（`原图/高一98_市西中学_10月月考/01..08.png`），
% 尺寸混排（第 1—4、6、7 页为 4762×6735，第 5、8 页为 2560×3620），按页序读取。
%
% 篇幅与题号：原卷自第 3 题起编号（第 1、2 题不在原文内）。填空题 3—12（10 题）、
% 选择题 13—16（4 题）、解答题 17—21（5 题），共 19 题。
%
% 讲评版：全部 19 题均印【解析】，其中第 13—16 题的答案在解析末句
% “故选 $B$／$D$／$A$／$D$”（题干括号内未印答案）。全卷无插图。
%
% 副标题：原卷未印副标题。本稿按内容标作“集合与常用逻辑用语、不等式”。
%
% 原卷笔误/存疑（均照录未改，正文对应处另有 % 注）：
%   ① 第 21 题 (1) 求证行、解析“假设”行与“故”行三处把 $x_3$ 误印作 $x_2$
%      （“$x_1$、$x_2$、$x_2$、$x_4$、$x_5$”）；解析同一行后半句
%      “这与 $x_1+x_2+x_3+x_4+x_5=5$ 矛盾”却用 $x_3$，显系笔误，照录；
%   ② 第 6 题解析的判别式行印作“$\Delta=(-4)^2-4a\geqslant0$”，与
%      $f(1)>0$ 并列，照录。
%
% 与原卷的排版差异：
%   ① 原卷每题下直接印【解析】，本稿按全稿惯例将答案与解析置于答案版；
%   ② 原卷未印副标题，本稿补出；填空横线统一排 \blankline[4.4em]；
%   ③ 原卷数集记号排斜体 $R$、$\Z$，本稿按全稿惯例改用 $\R$、$\Z$；
%   ④ 解答题按惯例加 \ansspace；选择题括号排 \mbox{（\quad）}；
%   ⑤ 第 12 题定义的运算 $\odot$ 原卷排实心小圆，本稿排 $\odot$。
%
% 行数核对：逐题数 \step/\stepm 条数，合计 82 条（Q3：1；Q4：2；Q5：1；Q6：4；Q7：3；
% Q8：2；Q9：5；Q10：2；Q11：5；Q12：4；Q13：3；Q14：2；Q15：2；Q16：9；Q17：2；
% Q18：11；Q19：6；Q20：11；Q21：7）。全卷 19 题均印【解析】。
% 原卷解析的换行按语意合并，未逐行强拆。
%
\documentclass[a4paper,11pt]{article}
\usepackage{common}
\begin{document}

\examtitlest[7]{2026--2027 年市西中学高一上 10 月月考}{集合与常用逻辑用语、不等式}

\bigsec{一、填空题}

\begin{enumerate}
\setcounter{enumi}{2}

\item 集合 $A=\{x\mid y=\sqrt{1-x^2}\},B=\{y\mid y=x^2+1,x\in A\}$，则 $A\cap\overline B=\blankline[4.4em]$.

\ifanswer
\solbegin
\step{$1-x^2\geqslant0\Rightarrow A=[-1,1]$，从而 $B=[1,2]$，所以 $A\cap\overline B=[-1,1)$.}
\solend
\fi

\item 已知 $b,c\in R$，关于 $x$ 的不等式 $x^2-bx+c<0$ 的解集为 $(-3,2)$，则 $b+c=\blankline[4.4em]$.

\ifanswer
\solbegin
\step{由题意得 $x=-3$ 与 $x=2$ 是 $x^2-bx+c=0$ 的两根，}
\step{所以 $\begin{cases}-3+2=b\\ -3\times2=c\end{cases}$，解得 $\begin{cases}b=-1\\ c=-6\end{cases}$，所以 $b+c=-7$.}
\solend
\fi

\item 若 $\alpha:\lvert x-m\rvert<1$ 是 $\beta:x^2-4x<0$ 的充分条件，则实数 $m$ 的取值范围是\blankline[4.4em].

\ifanswer
\solbegin
\step{$(m-1,m+1)\subseteq(0,4)\Leftrightarrow m-1\geqslant0$ 且 $m+1\leqslant4\Leftrightarrow m\in[1,3]$.}
\solend
\fi

\item 已知方程 $x^2-4x+a=0$ 的两根都大于 $1$，则 $a$ 的取值范围为\blankline[4.4em].

\ifanswer
\solbegin
\step{设函数 $f(x)=x^2-4x+a$，对称轴为 $x=2$，}
\step{因为方程 $x^2-4x+a=0$ 的两根都大于 $1$，}
\step{所以 $\begin{cases}\Delta=(-4)^2-4a\geqslant0\\ f(1)>0\end{cases}$，即 $\begin{cases}16-4a\geqslant0\\ 1-4+a>0\end{cases}$，解得 $3<a\leqslant4$，}
\step{即 $a$ 的取值范围为 $(3,4]$.}
\solend
\fi

\item 若不等式 $(a-2)x^2+2(a-2)x-4<0$ 对 $x\in R$ 恒成立，则实数 $a$ 的取值范围为\blankline[4.4em].

\ifanswer
\solbegin
\step{①当 $a-2=0$，即 $a=2$ 时，不等式为 $-4<0$ 恒成立，符合题意；}
\step{②当 $a\neq2$ 时，则 $\begin{cases}a-2<0\\ \Delta=4(a-2)^2+16(a-2)<0\end{cases}$，解得 $-2<a<2$.}
\step{综上所述，实数 $a$ 的取值范围是 $(-2,2]$.}
\solend
\fi

\item 已知 $-1\leqslant a+b\leqslant4,2\leqslant a-b\leqslant3$，则 $3a-2b$ 的取值范围为\blankline[4.4em].

\ifanswer
\solbegin
\step{设 $3a-2b=m(a+b)+n(a-b)$，则 $m+n=3,m-n=-2$，}
\step{所以 $m=\dfrac12,n=\dfrac52$，所以 $3a-2b=\dfrac12(a+b)+\dfrac52(a-b)\in\left[\dfrac92,\dfrac{19}2\right]$.}
\solend
\fi

\item 设 $x_1$、$x_2$ 是关于 $x$ 的方程 $2x^2-4mx+2m^2+3m-2=0$ 的两个实数根，则 $\lvert x_1-x_2\rvert$ 的最小值为\blankline[4.4em].

\ifanswer
\solbegin
\step{$x_1$、$x_2$ 是方程 $2x^2-4mx+2m^2+3m-2=0$ 的两个实数根，}
\step{所以 $\Delta=(4m)^2-8(2m^2+3m-2)\geqslant0$，解得 $m\leqslant\dfrac23$，}
\step{由韦达定理得 $x_1+x_2=2m,x_1\cdot x_2=\dfrac{2m^2+3m-2}2$，}
\step{则 $\lvert x_1-x_2\rvert=\sqrt{(x_1+x_2)^2-4x_1\cdot x_2}=\sqrt{4-6m}$，}
\step{当 $m=\dfrac23$ 时，$\lvert x_1-x_2\rvert$ 的最小值为 $0$.}
\solend
\fi

\item 设 $x\in R$，$[x]$ 表示不大于 $x$ 的最大整数，如 $[\pi]=3$，$[-1.2]=-2$，$\left[\dfrac12\right]=0$，则使 $\lvert\lvert x^2-1\rvert\rvert=3$ 成立的 $x$ 的取值范围是\blankline[4.4em].

\ifanswer
\solbegin
\step{若 $\lvert\lvert x^2-1\rvert\rvert=3$，则 $3\leqslant\lvert x^2-1\rvert<4$，即 $4\leqslant x^2<5$，}
\step{解得 $2\leqslant x<\sqrt5$ 或 $-\sqrt5<x\leqslant-2$，故 $x$ 的取值范围是 $(-\sqrt5,-2]\cup[2,\sqrt5)$.}
\solend
\fi

\item 设 $a$ 为实数，若关于 $x$ 的一元一次不等式组 $\begin{cases}2x+a>0\\ 3x-6a<0\end{cases}$ 的解集中有且仅有 $4$ 个整数，则 $a$ 的取值范围是\blankline[4.4em].

\ifanswer
\solbegin
\step{一元一次不等式组 $\begin{cases}2x+a>0\\ 3x-6a<0\end{cases}\Leftrightarrow-\dfrac a2<x<2a$，且 $a>0$，}
\step{关于 $x$ 的一元一次不等式组 $\begin{cases}2x+a>0\\ 3x-6a<0\end{cases}$ 的解集中有且仅有 $4$ 个整数，}
\step{①当 $4$ 个整数为 $0$，$1$，$2$，$3$ 时，则 $\begin{cases}-1\leqslant-\dfrac a2<0\\ 3<2a\leqslant4\end{cases}$，所以 $\dfrac32<a\leqslant2$，}
\step{②当 $4$ 个整数为 $-1$，$0$，$1$，$2$ 时，则 $\begin{cases}-2\leqslant-\dfrac a2<-1\\ 2<2a\leqslant3\end{cases}$，无解，}
\step{综上，$a$ 的取值范围是 $\left(\dfrac32,2\right]$.}
\solend
\fi

\item 集合 $M$、$N$、$S$ 都是非空集合，现规定如下运算：

$M\odot N\odot S=\{x\mid x\in(M\cap N)\cup(N\cap S)\cup(S\cap M)$ 且 $x\notin M\cap N\cap S\}$．假设集合 $A=\{x\mid a<x<b\}$，$B=\{x\mid c<x<d\}$，$C=\{x\mid e<x<f\}$，其中实数 $a$、$b$、$c$、$d$、$e$、$f$ 满足：$a<c<e<0<b<d<f$．计算 $A\odot B\odot C=\blankline[4.4em]$.

\ifanswer
\solbegin
\step{因为集合 $A=\{x\mid a<x<b\}$，$B=\{x\mid c<x<d\}$，$C=\{x\mid e<x<f\}$，}
\step{又因为 $a<c<e<0<b<d<f$，}
\step{所以 $A\cap B=\{x\mid c<x<b\}$，$B\cap C=\{x\mid e<x<d\}$，$C\cap A=\{x\mid e<x<b\}$，}
\step{$A\cap B\cap C=\{x\mid e<x<b\}$，故 $A\odot B\odot C=\{x\mid c<x\leqslant e$ 或 $b\leqslant x<d\}$.}
\solend
\fi

\end{enumerate}

\bigsec{二、选择题}

\begin{enumerate}
\setcounter{enumi}{12}

\item “$x\neq1$”是“$x^2-1\neq0$”的\mbox{（\quad）}条件.

\keepwithprev
A. 充分非必要\qquad B. 必要非充分

C. 充要\qquad D. 既非充分也非必要

\ifanswer
\solbegin
\step{$x^2-1\neq0\Leftrightarrow x\neq1$ 且 $x\neq-1\Rightarrow x\neq1$.}
\step{反之不成立，如 $x=-1$ 时，$x\neq1$，但 $x^2-1=0$.}
\step{因此“$x\neq1$”是“$x^2-1\neq0$”的必要非充分条件，故选 $B$.}
\solend
\fi

\item 若 $\dfrac1a<\dfrac1b<0$，则下列结论中不正确的是\mbox{（\quad）}

\keepwithprev
A. $a+b<0$\qquad B. $a^2<b^2$\qquad C. $ab<b^2$\qquad D. $ab<a^2$

\ifanswer
\solbegin
\step{由 $\dfrac1a<\dfrac1b<0$ 得 $b<a<0$，因此 $a+b<0$，且 $a^2<ab<b^2$，}
\step{所以 A、B、C 正确，D 不正确．故选 $D$.}
\solend
\fi

\item 非空集合 $X=\{x\mid a+1\leqslant x\leqslant3a-5\}$，$Y=\{x\mid1\leqslant x\leqslant16\}$，使得 $X\subseteq(X\cap Y)$ 成立的所有 $a$ 的集合是\mbox{（\quad）}

\keepwithprev
A. $\{a\mid3\leqslant a\leqslant7\}$\qquad B. $\{a\mid0\leqslant a\leqslant7\}$\qquad C. $\{a\mid3<a\leqslant7\}$\qquad D. $\{a\mid a\leqslant7\}$

\ifanswer
\solbegin
\step{因为 $X\subseteq(X\cap Y)$ 成立，所以 $X\subseteq Y$ 成立.}
\step{$X\neq\varnothing$，要使 $X\subseteq Y$，则 $\begin{cases}a\geqslant3\\ a+1\geqslant1\\ 3a-5\leqslant16\end{cases}$，即 $\begin{cases}a\geqslant3\\ a\geqslant0\\ a\leqslant7\end{cases}$，所以 $3\leqslant a\leqslant7$，故选 $A$.}
\solend
\fi

\item 已知 $M$ 是满足下列条件的集合：①$0\in M$，$1\in M$；②若 $x$、$y\in M$，则 $x-y\in M$；③若 $x\in M$ 且 $x\neq0$，则 $\dfrac1x\in M$．则下列说法正确的个数为\mbox{（\quad）}

（1）$\dfrac12\in M$；（2）$x+y\in M$；（3）$x^2\in M$．

\keepwithprev
A. $0$\qquad B. $1$\qquad C. $2$\qquad D. $3$

\ifanswer
\solbegin
\step{（1）$\dfrac12\in M$ 正确，理由如下：}
\step{由①得 $0\in M$，$1\in M$，由②得 $0-1=-1\in M$，$1-(-1)=2\in M$，}
\step{由③得 $\dfrac12\in M$；}
\step{（2）由①得 $0\in M$，由题意 $y\in M$，}
\step{所以由②得 $0-y=-y\in M$，又 $x\in M$，所以 $x-(-y)=x+y\in M$；}
\step{（3）$x\in M$，由②得 $x-1\in M$，由③得 $\dfrac1x\in M$，$\dfrac1{x-1}\in M$，}
\step{所以 $\dfrac1x-\dfrac1{x-1}\in M$，即 $\dfrac1{x(x-1)}\in M$，所以 $x(x-1)\in M$，}
\step{由（2）结论得 $x(x-1)+x\in M$，即 $x^2\in M$，}
\step{所以说法正确的个数为 $3$．故选 $D$.}
\solend
\fi

\end{enumerate}

\bigsec{三、解答题}

\begin{enumerate}
\setcounter{enumi}{16}

\item 已知 $a>b$，用比较法证明：$a^3>b^3$.

\ansspace*[8]

\ifanswer
\solbegin
\step{$a>b$，$a^3-b^3=(a-b)(a^2+ab+b^2)=(a-b)\left[\left(a+\dfrac b2\right)^2+\dfrac34b^2\right]>0$，}
\step{故 $a^3>b^3$.}
\solend
\fi

\item 求下列关于 $x$ 的不等式的解集：

（1）$\dfrac{x^2(x+3)(x^2-x+1)}{x^2-5x+4}>0$；

（2）$\dfrac{ax-1}{x-2}\geqslant0\left(a<\dfrac12\right)$.

\ansspace*[12]

\ifanswer
\solbegin
\stepm{（1）}{分母 $x^2-5x+4=(x-1)(x-4)$，要求 $x\neq1,4$；}
\step{又 $x=0$ 时原不等式不成立，且 $x^2-x+1=\left(x-\dfrac12\right)^2+\dfrac34>0$，}
\step{因此原不等式等价于 $\dfrac{x+3}{(x-1)(x-4)}>0,x\neq0$.}
\step{在数轴上标出 $-3,1,4$，分段判号得 $-3<x<1$ 或 $x>4$，}
\step{再排除 $x=0$，解集为 $(-3,0)\cup(0,1)\cup(4,+\infty)$.}
\stepm{（2）}{原不等式要求 $x\neq2$，按 $a$ 的符号分类讨论.}
\step{①当 $a<0$ 时，$\dfrac1a<2$，原不等式等价于 $\dfrac{x-\dfrac1a}{x-2}\leqslant0\Leftrightarrow\dfrac1a\leqslant x<2$，}
\step{解集为 $\left[\dfrac1a,2\right)$.}
\step{②当 $a=0$ 时，原不等式为 $\dfrac{-1}{x-2}\geqslant0\Leftrightarrow x<2$，解集为 $(-\infty,2)$.}
\step{③当 $0<a<\dfrac12$ 时，$\dfrac1a>2$，}
\step{原不等式等价于 $\dfrac{x-\dfrac1a}{x-2}\geqslant0\Leftrightarrow x<2$ 或 $x\geqslant\dfrac1a$，解集为 $(-\infty,2)\cup\left[\dfrac1a,+\infty\right)$.}
\solend
\fi

\item 集合 $A=\{x\mid-2\leqslant x\leqslant5\}$，$B=\{x\mid m+1\leqslant x\leqslant2m-1\}$.

（1）若 $B\subseteq A$，求实数 $m$ 的取值范围；

（2）若 $A\cap B\neq\varnothing$，求实数 $m$ 的取值范围.

\ansspace*[10]

\ifanswer
\solbegin
\stepm{（1）}{当 $B=\varnothing$ 时，$m+1>2m-1$，解得 $m<2$，}
\step{当 $B\neq\varnothing$ 时，$\begin{cases}m+1\leqslant2m-1\\ m+1\geqslant-2\\ 2m-1\leqslant5\end{cases}$，解得 $2\leqslant m\leqslant3$，}
\step{综上，实数 $m$ 的取值范围是 $(-\infty,3]$；}
\stepm{（2）}{当 $B=\varnothing$ 时，$m+1>2m-1$，解得 $m<2$，}
\step{当 $B\neq\varnothing$ 时，$\begin{cases}m+1\leqslant2m-1\\ 2m-1<-2\end{cases}$ 或 $\begin{cases}m+1\leqslant2m-1\\ m+1>5\end{cases}$，解得 $m\in\varnothing$ 或 $m>4$.}
\step{所以当 $A\cap B\neq\varnothing$ 时，实数 $m$ 的取值范围是 $[2,4]$.}
\solend
\fi

\item 在解决实际问题时，往往会有不同的思路和方法，这些方法有些正确，有些错误；有些简洁，有些复杂.

问题①：设 $a\in R$，集合 $M=\{x\mid(x-1)(x-a)=0\}$，若 $x\in M$ 是 $x\in\{1,2,3\}$ 的充分条件，求：$a$ 的取值集合.

问题②：设 $x\in Z$，$a=x+3$，$b=18-x$，若 $c\in Z$，求证：$c-a$ 和 $c-b$ 至少有一个数是奇数

（1）小明在解决问题①，他认为原问题等价于 $\{1,a\}\subseteq\{1,2,3\}$，解得 $a$ 的取值集合为 $\{2,3\}$，张老师判断小明解题错误，请计算正确的 $a$ 的取值集合；

（2）小红认为既然 $c\in Z$，只需根据 $c$ 是奇数还是偶数，分类讨论即可；小华则认为可以使用反证法解决问题，请你选一种你认为更好的方法并证明.

\ansspace*[12]

\ifanswer
\solbegin
\stepm{（1）}{当 $a=1$ 时，集合 $M=\{x\mid(x-1)^2=0\}=\{1\}$，}
\step{此时 $x\in M\Rightarrow x\in\{1,2,3\}$，即 $x\in M$ 是 $x\in\{1,2,3\}$ 的充分条件.}
\step{当 $a\neq1$ 时，集合 $M=\{x\mid(x-1)(x-a)=0\}=\{1,a\}$，}
\step{由 $x\in M$ 是 $x\in\{1,2,3\}$ 的充分条件，得 $\{1,a\}\subseteq\{1,2,3\}$，}
\step{所以 $a=2$ 或 $a=3$.}
\step{综上，实数 $a$ 的取值集合为 $\{1,2,3\}$；}
\stepm{（2）}{利用反证法证明．证明如下：}
\step{假设 $c-a$、$c-b$ 均为偶数，则 $(c-b)-(c-a)=a-b$ 为偶数.}
\step{因为 $a-b=x+3-18+x=2x-15$ 为奇数，}
\step{与 $b-a$ 为偶数相矛盾，所以原假设不成立.}
\step{故 $c-a$、$c-b$ 至少有一个数是奇数.}
\solend
\fi

\item 已知实数 $x_1$、$x_2$、$x_3$、$x_4$、$x_5$，满足 $x_1+x_2+x_3+x_4+x_5=5$.

（1）求证：$x_1$、$x_2$、$x_2$、$x_4$、$x_5$ 中至少有一个实数不小于 $1$；% 原卷把 x_3 误印作 x_2，照录

（2）设 $x_1$、$x_2$、$x_3$、$x_4$、$x_5$ 这五个实数两两不等，集合 $A=\{x_1,x_2,x_3,x_4,x_5\}$，若 $B\neq\varnothing$ 且 $B\subseteq A$，记 $G(B)$ 是 $B$ 中所有元素之和，对所有的 $B$，求 $G(B)$ 的平均值.

\ansspace*[13]

\ifanswer
\solbegin
\stepm{（1）}{假设 $x_1$、$x_2$、$x_2$、$x_4$、$x_5$ 中的每一个数都小于 $1$，}% 原卷 x_3 误印作 x_2，照录
\step{得 $x_1+x_2+x_2+x_4+x_5<5$，这与 $x_1+x_2+x_3+x_4+x_5=5$ 矛盾，}% 前半句原卷 x_3 误印作 x_2，照录
\step{故 $x_1$、$x_2$、$x_2$、$x_4$、$x_5$ 中至少有一个实数不小于 $1$；}% 原卷 x_3 误印作 x_2，照录
\stepm{（2）}{集合 $A=\{x_1,x_2,x_3,x_4,x_5\}$ 的非空子集个数为 $32-1=31$，}
\step{由于 $G(B)$ 是 $B$ 中所有元素之和，}
\step{得 $G(B)=16(x_1+x_2+x_3+x_4+x_5)=16\times5=80$，}
\step{则 $G(B)$ 的平均值为 $\dfrac{80}{31}$.}
\solend
\fi

\end{enumerate}

\end{document}
"
%
% 原始材料：微信公众号“魔都数学加油站”文章，连载“27学年高一98”；
% 原文标题“27学年高一98——2026-2027年上海市市西中学高一上10月月考”；
% 原文链接：https://mp.weixin.qq.com/s/duK4wz-HT3W9tjOeYQaVfg。
% 正文为 8 张 PNG 页（`原图/高一98_市西中学_10月月考/01..08.png`），
% 尺寸混排（第 1—4、6、7 页为 4762×6735，第 5、8 页为 2560×3620），按页序读取。
%
% 篇幅与题号：原卷自第 3 题起编号（第 1、2 题不在原文内）。填空题 3—12（10 题）、
% 选择题 13—16（4 题）、解答题 17—21（5 题），共 19 题。
%
% 讲评版：全部 19 题均印【解析】，其中第 13—16 题的答案在解析末句
% “故选 $B$／$D$／$A$／$D$”（题干括号内未印答案）。全卷无插图。
%
% 副标题：原卷未印副标题。本稿按内容标作“集合与常用逻辑用语、不等式”。
%
% 原卷笔误/存疑（均照录未改，正文对应处另有 % 注）：
%   ① 第 21 题 (1) 求证行、解析“假设”行与“故”行三处把 $x_3$ 误印作 $x_2$
%      （“$x_1$、$x_2$、$x_2$、$x_4$、$x_5$”）；解析同一行后半句
%      “这与 $x_1+x_2+x_3+x_4+x_5=5$ 矛盾”却用 $x_3$，显系笔误，照录；
%   ② 第 6 题解析的判别式行印作“$\Delta=(-4)^2-4a\geqslant0$”，与
%      $f(1)>0$ 并列，照录。
%
% 与原卷的排版差异：
%   ① 原卷每题下直接印【解析】，本稿按全稿惯例将答案与解析置于答案版；
%   ② 原卷未印副标题，本稿补出；填空横线统一排 \blankline[4.4em]；
%   ③ 原卷数集记号排斜体 $R$、$\Z$，本稿按全稿惯例改用 $\R$、$\Z$；
%   ④ 解答题按惯例加 \ansspace；选择题括号排 \mbox{（\quad）}；
%   ⑤ 第 12 题定义的运算 $\odot$ 原卷排实心小圆，本稿排 $\odot$。
%
% 行数核对：逐题数 \step/\stepm 条数，合计 82 条（Q3：1；Q4：2；Q5：1；Q6：4；Q7：3；
% Q8：2；Q9：5；Q10：2；Q11：5；Q12：4；Q13：3；Q14：2；Q15：2；Q16：9；Q17：2；
% Q18：11；Q19：6；Q20：11；Q21：7）。全卷 19 题均印【解析】。
% 原卷解析的换行按语意合并，未逐行强拆。
%
\documentclass[a4paper,11pt]{article}
\usepackage{common}
\begin{document}

\examtitlest[7]{2026--2027 年市西中学高一上 10 月月考}{集合与常用逻辑用语、不等式}

\bigsec{一、填空题}

\begin{enumerate}
\setcounter{enumi}{2}

\item 集合 $A=\{x\mid y=\sqrt{1-x^2}\},B=\{y\mid y=x^2+1,x\in A\}$，则 $A\cap\overline B=\blankline[4.4em]$.

\ifanswer
\solbegin
\step{$1-x^2\geqslant0\Rightarrow A=[-1,1]$，从而 $B=[1,2]$，所以 $A\cap\overline B=[-1,1)$.}
\solend
\fi

\item 已知 $b,c\in R$，关于 $x$ 的不等式 $x^2-bx+c<0$ 的解集为 $(-3,2)$，则 $b+c=\blankline[4.4em]$.

\ifanswer
\solbegin
\step{由题意得 $x=-3$ 与 $x=2$ 是 $x^2-bx+c=0$ 的两根，}
\step{所以 $\begin{cases}-3+2=b\\ -3\times2=c\end{cases}$，解得 $\begin{cases}b=-1\\ c=-6\end{cases}$，所以 $b+c=-7$.}
\solend
\fi

\item 若 $\alpha:\lvert x-m\rvert<1$ 是 $\beta:x^2-4x<0$ 的充分条件，则实数 $m$ 的取值范围是\blankline[4.4em].

\ifanswer
\solbegin
\step{$(m-1,m+1)\subseteq(0,4)\Leftrightarrow m-1\geqslant0$ 且 $m+1\leqslant4\Leftrightarrow m\in[1,3]$.}
\solend
\fi

\item 已知方程 $x^2-4x+a=0$ 的两根都大于 $1$，则 $a$ 的取值范围为\blankline[4.4em].

\ifanswer
\solbegin
\step{设函数 $f(x)=x^2-4x+a$，对称轴为 $x=2$，}
\step{因为方程 $x^2-4x+a=0$ 的两根都大于 $1$，}
\step{所以 $\begin{cases}\Delta=(-4)^2-4a\geqslant0\\ f(1)>0\end{cases}$，即 $\begin{cases}16-4a\geqslant0\\ 1-4+a>0\end{cases}$，解得 $3<a\leqslant4$，}
\step{即 $a$ 的取值范围为 $(3,4]$.}
\solend
\fi

\item 若不等式 $(a-2)x^2+2(a-2)x-4<0$ 对 $x\in R$ 恒成立，则实数 $a$ 的取值范围为\blankline[4.4em].

\ifanswer
\solbegin
\step{①当 $a-2=0$，即 $a=2$ 时，不等式为 $-4<0$ 恒成立，符合题意；}
\step{②当 $a\neq2$ 时，则 $\begin{cases}a-2<0\\ \Delta=4(a-2)^2+16(a-2)<0\end{cases}$，解得 $-2<a<2$.}
\step{综上所述，实数 $a$ 的取值范围是 $(-2,2]$.}
\solend
\fi

\item 已知 $-1\leqslant a+b\leqslant4,2\leqslant a-b\leqslant3$，则 $3a-2b$ 的取值范围为\blankline[4.4em].

\ifanswer
\solbegin
\step{设 $3a-2b=m(a+b)+n(a-b)$，则 $m+n=3,m-n=-2$，}
\step{所以 $m=\dfrac12,n=\dfrac52$，所以 $3a-2b=\dfrac12(a+b)+\dfrac52(a-b)\in\left[\dfrac92,\dfrac{19}2\right]$.}
\solend
\fi

\item 设 $x_1$、$x_2$ 是关于 $x$ 的方程 $2x^2-4mx+2m^2+3m-2=0$ 的两个实数根，则 $\lvert x_1-x_2\rvert$ 的最小值为\blankline[4.4em].

\ifanswer
\solbegin
\step{$x_1$、$x_2$ 是方程 $2x^2-4mx+2m^2+3m-2=0$ 的两个实数根，}
\step{所以 $\Delta=(4m)^2-8(2m^2+3m-2)\geqslant0$，解得 $m\leqslant\dfrac23$，}
\step{由韦达定理得 $x_1+x_2=2m,x_1\cdot x_2=\dfrac{2m^2+3m-2}2$，}
\step{则 $\lvert x_1-x_2\rvert=\sqrt{(x_1+x_2)^2-4x_1\cdot x_2}=\sqrt{4-6m}$，}
\step{当 $m=\dfrac23$ 时，$\lvert x_1-x_2\rvert$ 的最小值为 $0$.}
\solend
\fi

\item 设 $x\in R$，$[x]$ 表示不大于 $x$ 的最大整数，如 $[\pi]=3$，$[-1.2]=-2$，$\left[\dfrac12\right]=0$，则使 $\lvert\lvert x^2-1\rvert\rvert=3$ 成立的 $x$ 的取值范围是\blankline[4.4em].

\ifanswer
\solbegin
\step{若 $\lvert\lvert x^2-1\rvert\rvert=3$，则 $3\leqslant\lvert x^2-1\rvert<4$，即 $4\leqslant x^2<5$，}
\step{解得 $2\leqslant x<\sqrt5$ 或 $-\sqrt5<x\leqslant-2$，故 $x$ 的取值范围是 $(-\sqrt5,-2]\cup[2,\sqrt5)$.}
\solend
\fi

\item 设 $a$ 为实数，若关于 $x$ 的一元一次不等式组 $\begin{cases}2x+a>0\\ 3x-6a<0\end{cases}$ 的解集中有且仅有 $4$ 个整数，则 $a$ 的取值范围是\blankline[4.4em].

\ifanswer
\solbegin
\step{一元一次不等式组 $\begin{cases}2x+a>0\\ 3x-6a<0\end{cases}\Leftrightarrow-\dfrac a2<x<2a$，且 $a>0$，}
\step{关于 $x$ 的一元一次不等式组 $\begin{cases}2x+a>0\\ 3x-6a<0\end{cases}$ 的解集中有且仅有 $4$ 个整数，}
\step{①当 $4$ 个整数为 $0$，$1$，$2$，$3$ 时，则 $\begin{cases}-1\leqslant-\dfrac a2<0\\ 3<2a\leqslant4\end{cases}$，所以 $\dfrac32<a\leqslant2$，}
\step{②当 $4$ 个整数为 $-1$，$0$，$1$，$2$ 时，则 $\begin{cases}-2\leqslant-\dfrac a2<-1\\ 2<2a\leqslant3\end{cases}$，无解，}
\step{综上，$a$ 的取值范围是 $\left(\dfrac32,2\right]$.}
\solend
\fi

\item 集合 $M$、$N$、$S$ 都是非空集合，现规定如下运算：

$M\odot N\odot S=\{x\mid x\in(M\cap N)\cup(N\cap S)\cup(S\cap M)$ 且 $x\notin M\cap N\cap S\}$．假设集合 $A=\{x\mid a<x<b\}$，$B=\{x\mid c<x<d\}$，$C=\{x\mid e<x<f\}$，其中实数 $a$、$b$、$c$、$d$、$e$、$f$ 满足：$a<c<e<0<b<d<f$．计算 $A\odot B\odot C=\blankline[4.4em]$.

\ifanswer
\solbegin
\step{因为集合 $A=\{x\mid a<x<b\}$，$B=\{x\mid c<x<d\}$，$C=\{x\mid e<x<f\}$，}
\step{又因为 $a<c<e<0<b<d<f$，}
\step{所以 $A\cap B=\{x\mid c<x<b\}$，$B\cap C=\{x\mid e<x<d\}$，$C\cap A=\{x\mid e<x<b\}$，}
\step{$A\cap B\cap C=\{x\mid e<x<b\}$，故 $A\odot B\odot C=\{x\mid c<x\leqslant e$ 或 $b\leqslant x<d\}$.}
\solend
\fi

\end{enumerate}

\bigsec{二、选择题}

\begin{enumerate}
\setcounter{enumi}{12}

\item “$x\neq1$”是“$x^2-1\neq0$”的\mbox{（\quad）}条件.

\keepwithprev
A. 充分非必要\qquad B. 必要非充分

C. 充要\qquad D. 既非充分也非必要

\ifanswer
\solbegin
\step{$x^2-1\neq0\Leftrightarrow x\neq1$ 且 $x\neq-1\Rightarrow x\neq1$.}
\step{反之不成立，如 $x=-1$ 时，$x\neq1$，但 $x^2-1=0$.}
\step{因此“$x\neq1$”是“$x^2-1\neq0$”的必要非充分条件，故选 $B$.}
\solend
\fi

\item 若 $\dfrac1a<\dfrac1b<0$，则下列结论中不正确的是\mbox{（\quad）}

\keepwithprev
A. $a+b<0$\qquad B. $a^2<b^2$\qquad C. $ab<b^2$\qquad D. $ab<a^2$

\ifanswer
\solbegin
\step{由 $\dfrac1a<\dfrac1b<0$ 得 $b<a<0$，因此 $a+b<0$，且 $a^2<ab<b^2$，}
\step{所以 A、B、C 正确，D 不正确．故选 $D$.}
\solend
\fi

\item 非空集合 $X=\{x\mid a+1\leqslant x\leqslant3a-5\}$，$Y=\{x\mid1\leqslant x\leqslant16\}$，使得 $X\subseteq(X\cap Y)$ 成立的所有 $a$ 的集合是\mbox{（\quad）}

\keepwithprev
A. $\{a\mid3\leqslant a\leqslant7\}$\qquad B. $\{a\mid0\leqslant a\leqslant7\}$\qquad C. $\{a\mid3<a\leqslant7\}$\qquad D. $\{a\mid a\leqslant7\}$

\ifanswer
\solbegin
\step{因为 $X\subseteq(X\cap Y)$ 成立，所以 $X\subseteq Y$ 成立.}
\step{$X\neq\varnothing$，要使 $X\subseteq Y$，则 $\begin{cases}a\geqslant3\\ a+1\geqslant1\\ 3a-5\leqslant16\end{cases}$，即 $\begin{cases}a\geqslant3\\ a\geqslant0\\ a\leqslant7\end{cases}$，所以 $3\leqslant a\leqslant7$，故选 $A$.}
\solend
\fi

\item 已知 $M$ 是满足下列条件的集合：①$0\in M$，$1\in M$；②若 $x$、$y\in M$，则 $x-y\in M$；③若 $x\in M$ 且 $x\neq0$，则 $\dfrac1x\in M$．则下列说法正确的个数为\mbox{（\quad）}

（1）$\dfrac12\in M$；（2）$x+y\in M$；（3）$x^2\in M$．

\keepwithprev
A. $0$\qquad B. $1$\qquad C. $2$\qquad D. $3$

\ifanswer
\solbegin
\step{（1）$\dfrac12\in M$ 正确，理由如下：}
\step{由①得 $0\in M$，$1\in M$，由②得 $0-1=-1\in M$，$1-(-1)=2\in M$，}
\step{由③得 $\dfrac12\in M$；}
\step{（2）由①得 $0\in M$，由题意 $y\in M$，}
\step{所以由②得 $0-y=-y\in M$，又 $x\in M$，所以 $x-(-y)=x+y\in M$；}
\step{（3）$x\in M$，由②得 $x-1\in M$，由③得 $\dfrac1x\in M$，$\dfrac1{x-1}\in M$，}
\step{所以 $\dfrac1x-\dfrac1{x-1}\in M$，即 $\dfrac1{x(x-1)}\in M$，所以 $x(x-1)\in M$，}
\step{由（2）结论得 $x(x-1)+x\in M$，即 $x^2\in M$，}
\step{所以说法正确的个数为 $3$．故选 $D$.}
\solend
\fi

\end{enumerate}

\bigsec{三、解答题}

\begin{enumerate}
\setcounter{enumi}{16}

\item 已知 $a>b$，用比较法证明：$a^3>b^3$.

\ansspace*[8]

\ifanswer
\solbegin
\step{$a>b$，$a^3-b^3=(a-b)(a^2+ab+b^2)=(a-b)\left[\left(a+\dfrac b2\right)^2+\dfrac34b^2\right]>0$，}
\step{故 $a^3>b^3$.}
\solend
\fi

\item 求下列关于 $x$ 的不等式的解集：

（1）$\dfrac{x^2(x+3)(x^2-x+1)}{x^2-5x+4}>0$；

（2）$\dfrac{ax-1}{x-2}\geqslant0\left(a<\dfrac12\right)$.

\ansspace*[12]

\ifanswer
\solbegin
\stepm{（1）}{分母 $x^2-5x+4=(x-1)(x-4)$，要求 $x\neq1,4$；}
\step{又 $x=0$ 时原不等式不成立，且 $x^2-x+1=\left(x-\dfrac12\right)^2+\dfrac34>0$，}
\step{因此原不等式等价于 $\dfrac{x+3}{(x-1)(x-4)}>0,x\neq0$.}
\step{在数轴上标出 $-3,1,4$，分段判号得 $-3<x<1$ 或 $x>4$，}
\step{再排除 $x=0$，解集为 $(-3,0)\cup(0,1)\cup(4,+\infty)$.}
\stepm{（2）}{原不等式要求 $x\neq2$，按 $a$ 的符号分类讨论.}
\step{①当 $a<0$ 时，$\dfrac1a<2$，原不等式等价于 $\dfrac{x-\dfrac1a}{x-2}\leqslant0\Leftrightarrow\dfrac1a\leqslant x<2$，}
\step{解集为 $\left[\dfrac1a,2\right)$.}
\step{②当 $a=0$ 时，原不等式为 $\dfrac{-1}{x-2}\geqslant0\Leftrightarrow x<2$，解集为 $(-\infty,2)$.}
\step{③当 $0<a<\dfrac12$ 时，$\dfrac1a>2$，}
\step{原不等式等价于 $\dfrac{x-\dfrac1a}{x-2}\geqslant0\Leftrightarrow x<2$ 或 $x\geqslant\dfrac1a$，解集为 $(-\infty,2)\cup\left[\dfrac1a,+\infty\right)$.}
\solend
\fi

\item 集合 $A=\{x\mid-2\leqslant x\leqslant5\}$，$B=\{x\mid m+1\leqslant x\leqslant2m-1\}$.

（1）若 $B\subseteq A$，求实数 $m$ 的取值范围；

（2）若 $A\cap B\neq\varnothing$，求实数 $m$ 的取值范围.

\ansspace*[10]

\ifanswer
\solbegin
\stepm{（1）}{当 $B=\varnothing$ 时，$m+1>2m-1$，解得 $m<2$，}
\step{当 $B\neq\varnothing$ 时，$\begin{cases}m+1\leqslant2m-1\\ m+1\geqslant-2\\ 2m-1\leqslant5\end{cases}$，解得 $2\leqslant m\leqslant3$，}
\step{综上，实数 $m$ 的取值范围是 $(-\infty,3]$；}
\stepm{（2）}{当 $B=\varnothing$ 时，$m+1>2m-1$，解得 $m<2$，}
\step{当 $B\neq\varnothing$ 时，$\begin{cases}m+1\leqslant2m-1\\ 2m-1<-2\end{cases}$ 或 $\begin{cases}m+1\leqslant2m-1\\ m+1>5\end{cases}$，解得 $m\in\varnothing$ 或 $m>4$.}
\step{所以当 $A\cap B\neq\varnothing$ 时，实数 $m$ 的取值范围是 $[2,4]$.}
\solend
\fi

\item 在解决实际问题时，往往会有不同的思路和方法，这些方法有些正确，有些错误；有些简洁，有些复杂.

问题①：设 $a\in R$，集合 $M=\{x\mid(x-1)(x-a)=0\}$，若 $x\in M$ 是 $x\in\{1,2,3\}$ 的充分条件，求：$a$ 的取值集合.

问题②：设 $x\in Z$，$a=x+3$，$b=18-x$，若 $c\in Z$，求证：$c-a$ 和 $c-b$ 至少有一个数是奇数

（1）小明在解决问题①，他认为原问题等价于 $\{1,a\}\subseteq\{1,2,3\}$，解得 $a$ 的取值集合为 $\{2,3\}$，张老师判断小明解题错误，请计算正确的 $a$ 的取值集合；

（2）小红认为既然 $c\in Z$，只需根据 $c$ 是奇数还是偶数，分类讨论即可；小华则认为可以使用反证法解决问题，请你选一种你认为更好的方法并证明.

\ansspace*[12]

\ifanswer
\solbegin
\stepm{（1）}{当 $a=1$ 时，集合 $M=\{x\mid(x-1)^2=0\}=\{1\}$，}
\step{此时 $x\in M\Rightarrow x\in\{1,2,3\}$，即 $x\in M$ 是 $x\in\{1,2,3\}$ 的充分条件.}
\step{当 $a\neq1$ 时，集合 $M=\{x\mid(x-1)(x-a)=0\}=\{1,a\}$，}
\step{由 $x\in M$ 是 $x\in\{1,2,3\}$ 的充分条件，得 $\{1,a\}\subseteq\{1,2,3\}$，}
\step{所以 $a=2$ 或 $a=3$.}
\step{综上，实数 $a$ 的取值集合为 $\{1,2,3\}$；}
\stepm{（2）}{利用反证法证明．证明如下：}
\step{假设 $c-a$、$c-b$ 均为偶数，则 $(c-b)-(c-a)=a-b$ 为偶数.}
\step{因为 $a-b=x+3-18+x=2x-15$ 为奇数，}
\step{与 $b-a$ 为偶数相矛盾，所以原假设不成立.}
\step{故 $c-a$、$c-b$ 至少有一个数是奇数.}
\solend
\fi

\item 已知实数 $x_1$、$x_2$、$x_3$、$x_4$、$x_5$，满足 $x_1+x_2+x_3+x_4+x_5=5$.

（1）求证：$x_1$、$x_2$、$x_2$、$x_4$、$x_5$ 中至少有一个实数不小于 $1$；% 原卷把 x_3 误印作 x_2，照录

（2）设 $x_1$、$x_2$、$x_3$、$x_4$、$x_5$ 这五个实数两两不等，集合 $A=\{x_1,x_2,x_3,x_4,x_5\}$，若 $B\neq\varnothing$ 且 $B\subseteq A$，记 $G(B)$ 是 $B$ 中所有元素之和，对所有的 $B$，求 $G(B)$ 的平均值.

\ansspace*[13]

\ifanswer
\solbegin
\stepm{（1）}{假设 $x_1$、$x_2$、$x_2$、$x_4$、$x_5$ 中的每一个数都小于 $1$，}% 原卷 x_3 误印作 x_2，照录
\step{得 $x_1+x_2+x_2+x_4+x_5<5$，这与 $x_1+x_2+x_3+x_4+x_5=5$ 矛盾，}% 前半句原卷 x_3 误印作 x_2，照录
\step{故 $x_1$、$x_2$、$x_2$、$x_4$、$x_5$ 中至少有一个实数不小于 $1$；}% 原卷 x_3 误印作 x_2，照录
\stepm{（2）}{集合 $A=\{x_1,x_2,x_3,x_4,x_5\}$ 的非空子集个数为 $32-1=31$，}
\step{由于 $G(B)$ 是 $B$ 中所有元素之和，}
\step{得 $G(B)=16(x_1+x_2+x_3+x_4+x_5)=16\times5=80$，}
\step{则 $G(B)$ 的平均值为 $\dfrac{80}{31}$.}
\solend
\fi

\end{enumerate}

\end{document}
"
% 紧凑版：xelatex "\def\iscompact{1}% 高一98_市西中学_10月月考.tex
%   —— 高一上数学“2026-2027 年市西中学高一上 10 月月考”（讲评版）
% 试卷版：xelatex 高一98_市西中学_10月月考.tex
% 答案版：xelatex "\def\isanswer{1}% 高一98_市西中学_10月月考.tex
%   —— 高一上数学“2026-2027 年市西中学高一上 10 月月考”（讲评版）
% 试卷版：xelatex 高一98_市西中学_10月月考.tex
% 答案版：xelatex "\def\isanswer{1}\input{高一98_市西中学_10月月考.tex}"
% 紧凑版：xelatex "\def\iscompact{1}\input{高一98_市西中学_10月月考.tex}"
%
% 原始材料：微信公众号“魔都数学加油站”文章，连载“27学年高一98”；
% 原文标题“27学年高一98——2026-2027年上海市市西中学高一上10月月考”；
% 原文链接：https://mp.weixin.qq.com/s/duK4wz-HT3W9tjOeYQaVfg。
% 正文为 8 张 PNG 页（`原图/高一98_市西中学_10月月考/01..08.png`），
% 尺寸混排（第 1—4、6、7 页为 4762×6735，第 5、8 页为 2560×3620），按页序读取。
%
% 篇幅与题号：原卷自第 3 题起编号（第 1、2 题不在原文内）。填空题 3—12（10 题）、
% 选择题 13—16（4 题）、解答题 17—21（5 题），共 19 题。
%
% 讲评版：全部 19 题均印【解析】，其中第 13—16 题的答案在解析末句
% “故选 $B$／$D$／$A$／$D$”（题干括号内未印答案）。全卷无插图。
%
% 副标题：原卷未印副标题。本稿按内容标作“集合与常用逻辑用语、不等式”。
%
% 原卷笔误/存疑（均照录未改，正文对应处另有 % 注）：
%   ① 第 21 题 (1) 求证行、解析“假设”行与“故”行三处把 $x_3$ 误印作 $x_2$
%      （“$x_1$、$x_2$、$x_2$、$x_4$、$x_5$”）；解析同一行后半句
%      “这与 $x_1+x_2+x_3+x_4+x_5=5$ 矛盾”却用 $x_3$，显系笔误，照录；
%   ② 第 6 题解析的判别式行印作“$\Delta=(-4)^2-4a\geqslant0$”，与
%      $f(1)>0$ 并列，照录。
%
% 与原卷的排版差异：
%   ① 原卷每题下直接印【解析】，本稿按全稿惯例将答案与解析置于答案版；
%   ② 原卷未印副标题，本稿补出；填空横线统一排 \blankline[4.4em]；
%   ③ 原卷数集记号排斜体 $R$、$\Z$，本稿按全稿惯例改用 $\R$、$\Z$；
%   ④ 解答题按惯例加 \ansspace；选择题括号排 \mbox{（\quad）}；
%   ⑤ 第 12 题定义的运算 $\odot$ 原卷排实心小圆，本稿排 $\odot$。
%
% 行数核对：逐题数 \step/\stepm 条数，合计 82 条（Q3：1；Q4：2；Q5：1；Q6：4；Q7：3；
% Q8：2；Q9：5；Q10：2；Q11：5；Q12：4；Q13：3；Q14：2；Q15：2；Q16：9；Q17：2；
% Q18：11；Q19：6；Q20：11；Q21：7）。全卷 19 题均印【解析】。
% 原卷解析的换行按语意合并，未逐行强拆。
%
\documentclass[a4paper,11pt]{article}
\usepackage{common}
\begin{document}

\examtitlest[7]{2026--2027 年市西中学高一上 10 月月考}{集合与常用逻辑用语、不等式}

\bigsec{一、填空题}

\begin{enumerate}
\setcounter{enumi}{2}

\item 集合 $A=\{x\mid y=\sqrt{1-x^2}\},B=\{y\mid y=x^2+1,x\in A\}$，则 $A\cap\overline B=\blankline[4.4em]$.

\ifanswer
\solbegin
\step{$1-x^2\geqslant0\Rightarrow A=[-1,1]$，从而 $B=[1,2]$，所以 $A\cap\overline B=[-1,1)$.}
\solend
\fi

\item 已知 $b,c\in R$，关于 $x$ 的不等式 $x^2-bx+c<0$ 的解集为 $(-3,2)$，则 $b+c=\blankline[4.4em]$.

\ifanswer
\solbegin
\step{由题意得 $x=-3$ 与 $x=2$ 是 $x^2-bx+c=0$ 的两根，}
\step{所以 $\begin{cases}-3+2=b\\ -3\times2=c\end{cases}$，解得 $\begin{cases}b=-1\\ c=-6\end{cases}$，所以 $b+c=-7$.}
\solend
\fi

\item 若 $\alpha:\lvert x-m\rvert<1$ 是 $\beta:x^2-4x<0$ 的充分条件，则实数 $m$ 的取值范围是\blankline[4.4em].

\ifanswer
\solbegin
\step{$(m-1,m+1)\subseteq(0,4)\Leftrightarrow m-1\geqslant0$ 且 $m+1\leqslant4\Leftrightarrow m\in[1,3]$.}
\solend
\fi

\item 已知方程 $x^2-4x+a=0$ 的两根都大于 $1$，则 $a$ 的取值范围为\blankline[4.4em].

\ifanswer
\solbegin
\step{设函数 $f(x)=x^2-4x+a$，对称轴为 $x=2$，}
\step{因为方程 $x^2-4x+a=0$ 的两根都大于 $1$，}
\step{所以 $\begin{cases}\Delta=(-4)^2-4a\geqslant0\\ f(1)>0\end{cases}$，即 $\begin{cases}16-4a\geqslant0\\ 1-4+a>0\end{cases}$，解得 $3<a\leqslant4$，}
\step{即 $a$ 的取值范围为 $(3,4]$.}
\solend
\fi

\item 若不等式 $(a-2)x^2+2(a-2)x-4<0$ 对 $x\in R$ 恒成立，则实数 $a$ 的取值范围为\blankline[4.4em].

\ifanswer
\solbegin
\step{①当 $a-2=0$，即 $a=2$ 时，不等式为 $-4<0$ 恒成立，符合题意；}
\step{②当 $a\neq2$ 时，则 $\begin{cases}a-2<0\\ \Delta=4(a-2)^2+16(a-2)<0\end{cases}$，解得 $-2<a<2$.}
\step{综上所述，实数 $a$ 的取值范围是 $(-2,2]$.}
\solend
\fi

\item 已知 $-1\leqslant a+b\leqslant4,2\leqslant a-b\leqslant3$，则 $3a-2b$ 的取值范围为\blankline[4.4em].

\ifanswer
\solbegin
\step{设 $3a-2b=m(a+b)+n(a-b)$，则 $m+n=3,m-n=-2$，}
\step{所以 $m=\dfrac12,n=\dfrac52$，所以 $3a-2b=\dfrac12(a+b)+\dfrac52(a-b)\in\left[\dfrac92,\dfrac{19}2\right]$.}
\solend
\fi

\item 设 $x_1$、$x_2$ 是关于 $x$ 的方程 $2x^2-4mx+2m^2+3m-2=0$ 的两个实数根，则 $\lvert x_1-x_2\rvert$ 的最小值为\blankline[4.4em].

\ifanswer
\solbegin
\step{$x_1$、$x_2$ 是方程 $2x^2-4mx+2m^2+3m-2=0$ 的两个实数根，}
\step{所以 $\Delta=(4m)^2-8(2m^2+3m-2)\geqslant0$，解得 $m\leqslant\dfrac23$，}
\step{由韦达定理得 $x_1+x_2=2m,x_1\cdot x_2=\dfrac{2m^2+3m-2}2$，}
\step{则 $\lvert x_1-x_2\rvert=\sqrt{(x_1+x_2)^2-4x_1\cdot x_2}=\sqrt{4-6m}$，}
\step{当 $m=\dfrac23$ 时，$\lvert x_1-x_2\rvert$ 的最小值为 $0$.}
\solend
\fi

\item 设 $x\in R$，$[x]$ 表示不大于 $x$ 的最大整数，如 $[\pi]=3$，$[-1.2]=-2$，$\left[\dfrac12\right]=0$，则使 $\lvert\lvert x^2-1\rvert\rvert=3$ 成立的 $x$ 的取值范围是\blankline[4.4em].

\ifanswer
\solbegin
\step{若 $\lvert\lvert x^2-1\rvert\rvert=3$，则 $3\leqslant\lvert x^2-1\rvert<4$，即 $4\leqslant x^2<5$，}
\step{解得 $2\leqslant x<\sqrt5$ 或 $-\sqrt5<x\leqslant-2$，故 $x$ 的取值范围是 $(-\sqrt5,-2]\cup[2,\sqrt5)$.}
\solend
\fi

\item 设 $a$ 为实数，若关于 $x$ 的一元一次不等式组 $\begin{cases}2x+a>0\\ 3x-6a<0\end{cases}$ 的解集中有且仅有 $4$ 个整数，则 $a$ 的取值范围是\blankline[4.4em].

\ifanswer
\solbegin
\step{一元一次不等式组 $\begin{cases}2x+a>0\\ 3x-6a<0\end{cases}\Leftrightarrow-\dfrac a2<x<2a$，且 $a>0$，}
\step{关于 $x$ 的一元一次不等式组 $\begin{cases}2x+a>0\\ 3x-6a<0\end{cases}$ 的解集中有且仅有 $4$ 个整数，}
\step{①当 $4$ 个整数为 $0$，$1$，$2$，$3$ 时，则 $\begin{cases}-1\leqslant-\dfrac a2<0\\ 3<2a\leqslant4\end{cases}$，所以 $\dfrac32<a\leqslant2$，}
\step{②当 $4$ 个整数为 $-1$，$0$，$1$，$2$ 时，则 $\begin{cases}-2\leqslant-\dfrac a2<-1\\ 2<2a\leqslant3\end{cases}$，无解，}
\step{综上，$a$ 的取值范围是 $\left(\dfrac32,2\right]$.}
\solend
\fi

\item 集合 $M$、$N$、$S$ 都是非空集合，现规定如下运算：

$M\odot N\odot S=\{x\mid x\in(M\cap N)\cup(N\cap S)\cup(S\cap M)$ 且 $x\notin M\cap N\cap S\}$．假设集合 $A=\{x\mid a<x<b\}$，$B=\{x\mid c<x<d\}$，$C=\{x\mid e<x<f\}$，其中实数 $a$、$b$、$c$、$d$、$e$、$f$ 满足：$a<c<e<0<b<d<f$．计算 $A\odot B\odot C=\blankline[4.4em]$.

\ifanswer
\solbegin
\step{因为集合 $A=\{x\mid a<x<b\}$，$B=\{x\mid c<x<d\}$，$C=\{x\mid e<x<f\}$，}
\step{又因为 $a<c<e<0<b<d<f$，}
\step{所以 $A\cap B=\{x\mid c<x<b\}$，$B\cap C=\{x\mid e<x<d\}$，$C\cap A=\{x\mid e<x<b\}$，}
\step{$A\cap B\cap C=\{x\mid e<x<b\}$，故 $A\odot B\odot C=\{x\mid c<x\leqslant e$ 或 $b\leqslant x<d\}$.}
\solend
\fi

\end{enumerate}

\bigsec{二、选择题}

\begin{enumerate}
\setcounter{enumi}{12}

\item “$x\neq1$”是“$x^2-1\neq0$”的\mbox{（\quad）}条件.

\keepwithprev
A. 充分非必要\qquad B. 必要非充分

C. 充要\qquad D. 既非充分也非必要

\ifanswer
\solbegin
\step{$x^2-1\neq0\Leftrightarrow x\neq1$ 且 $x\neq-1\Rightarrow x\neq1$.}
\step{反之不成立，如 $x=-1$ 时，$x\neq1$，但 $x^2-1=0$.}
\step{因此“$x\neq1$”是“$x^2-1\neq0$”的必要非充分条件，故选 $B$.}
\solend
\fi

\item 若 $\dfrac1a<\dfrac1b<0$，则下列结论中不正确的是\mbox{（\quad）}

\keepwithprev
A. $a+b<0$\qquad B. $a^2<b^2$\qquad C. $ab<b^2$\qquad D. $ab<a^2$

\ifanswer
\solbegin
\step{由 $\dfrac1a<\dfrac1b<0$ 得 $b<a<0$，因此 $a+b<0$，且 $a^2<ab<b^2$，}
\step{所以 A、B、C 正确，D 不正确．故选 $D$.}
\solend
\fi

\item 非空集合 $X=\{x\mid a+1\leqslant x\leqslant3a-5\}$，$Y=\{x\mid1\leqslant x\leqslant16\}$，使得 $X\subseteq(X\cap Y)$ 成立的所有 $a$ 的集合是\mbox{（\quad）}

\keepwithprev
A. $\{a\mid3\leqslant a\leqslant7\}$\qquad B. $\{a\mid0\leqslant a\leqslant7\}$\qquad C. $\{a\mid3<a\leqslant7\}$\qquad D. $\{a\mid a\leqslant7\}$

\ifanswer
\solbegin
\step{因为 $X\subseteq(X\cap Y)$ 成立，所以 $X\subseteq Y$ 成立.}
\step{$X\neq\varnothing$，要使 $X\subseteq Y$，则 $\begin{cases}a\geqslant3\\ a+1\geqslant1\\ 3a-5\leqslant16\end{cases}$，即 $\begin{cases}a\geqslant3\\ a\geqslant0\\ a\leqslant7\end{cases}$，所以 $3\leqslant a\leqslant7$，故选 $A$.}
\solend
\fi

\item 已知 $M$ 是满足下列条件的集合：①$0\in M$，$1\in M$；②若 $x$、$y\in M$，则 $x-y\in M$；③若 $x\in M$ 且 $x\neq0$，则 $\dfrac1x\in M$．则下列说法正确的个数为\mbox{（\quad）}

（1）$\dfrac12\in M$；（2）$x+y\in M$；（3）$x^2\in M$．

\keepwithprev
A. $0$\qquad B. $1$\qquad C. $2$\qquad D. $3$

\ifanswer
\solbegin
\step{（1）$\dfrac12\in M$ 正确，理由如下：}
\step{由①得 $0\in M$，$1\in M$，由②得 $0-1=-1\in M$，$1-(-1)=2\in M$，}
\step{由③得 $\dfrac12\in M$；}
\step{（2）由①得 $0\in M$，由题意 $y\in M$，}
\step{所以由②得 $0-y=-y\in M$，又 $x\in M$，所以 $x-(-y)=x+y\in M$；}
\step{（3）$x\in M$，由②得 $x-1\in M$，由③得 $\dfrac1x\in M$，$\dfrac1{x-1}\in M$，}
\step{所以 $\dfrac1x-\dfrac1{x-1}\in M$，即 $\dfrac1{x(x-1)}\in M$，所以 $x(x-1)\in M$，}
\step{由（2）结论得 $x(x-1)+x\in M$，即 $x^2\in M$，}
\step{所以说法正确的个数为 $3$．故选 $D$.}
\solend
\fi

\end{enumerate}

\bigsec{三、解答题}

\begin{enumerate}
\setcounter{enumi}{16}

\item 已知 $a>b$，用比较法证明：$a^3>b^3$.

\ansspace*[8]

\ifanswer
\solbegin
\step{$a>b$，$a^3-b^3=(a-b)(a^2+ab+b^2)=(a-b)\left[\left(a+\dfrac b2\right)^2+\dfrac34b^2\right]>0$，}
\step{故 $a^3>b^3$.}
\solend
\fi

\item 求下列关于 $x$ 的不等式的解集：

（1）$\dfrac{x^2(x+3)(x^2-x+1)}{x^2-5x+4}>0$；

（2）$\dfrac{ax-1}{x-2}\geqslant0\left(a<\dfrac12\right)$.

\ansspace*[12]

\ifanswer
\solbegin
\stepm{（1）}{分母 $x^2-5x+4=(x-1)(x-4)$，要求 $x\neq1,4$；}
\step{又 $x=0$ 时原不等式不成立，且 $x^2-x+1=\left(x-\dfrac12\right)^2+\dfrac34>0$，}
\step{因此原不等式等价于 $\dfrac{x+3}{(x-1)(x-4)}>0,x\neq0$.}
\step{在数轴上标出 $-3,1,4$，分段判号得 $-3<x<1$ 或 $x>4$，}
\step{再排除 $x=0$，解集为 $(-3,0)\cup(0,1)\cup(4,+\infty)$.}
\stepm{（2）}{原不等式要求 $x\neq2$，按 $a$ 的符号分类讨论.}
\step{①当 $a<0$ 时，$\dfrac1a<2$，原不等式等价于 $\dfrac{x-\dfrac1a}{x-2}\leqslant0\Leftrightarrow\dfrac1a\leqslant x<2$，}
\step{解集为 $\left[\dfrac1a,2\right)$.}
\step{②当 $a=0$ 时，原不等式为 $\dfrac{-1}{x-2}\geqslant0\Leftrightarrow x<2$，解集为 $(-\infty,2)$.}
\step{③当 $0<a<\dfrac12$ 时，$\dfrac1a>2$，}
\step{原不等式等价于 $\dfrac{x-\dfrac1a}{x-2}\geqslant0\Leftrightarrow x<2$ 或 $x\geqslant\dfrac1a$，解集为 $(-\infty,2)\cup\left[\dfrac1a,+\infty\right)$.}
\solend
\fi

\item 集合 $A=\{x\mid-2\leqslant x\leqslant5\}$，$B=\{x\mid m+1\leqslant x\leqslant2m-1\}$.

（1）若 $B\subseteq A$，求实数 $m$ 的取值范围；

（2）若 $A\cap B\neq\varnothing$，求实数 $m$ 的取值范围.

\ansspace*[10]

\ifanswer
\solbegin
\stepm{（1）}{当 $B=\varnothing$ 时，$m+1>2m-1$，解得 $m<2$，}
\step{当 $B\neq\varnothing$ 时，$\begin{cases}m+1\leqslant2m-1\\ m+1\geqslant-2\\ 2m-1\leqslant5\end{cases}$，解得 $2\leqslant m\leqslant3$，}
\step{综上，实数 $m$ 的取值范围是 $(-\infty,3]$；}
\stepm{（2）}{当 $B=\varnothing$ 时，$m+1>2m-1$，解得 $m<2$，}
\step{当 $B\neq\varnothing$ 时，$\begin{cases}m+1\leqslant2m-1\\ 2m-1<-2\end{cases}$ 或 $\begin{cases}m+1\leqslant2m-1\\ m+1>5\end{cases}$，解得 $m\in\varnothing$ 或 $m>4$.}
\step{所以当 $A\cap B\neq\varnothing$ 时，实数 $m$ 的取值范围是 $[2,4]$.}
\solend
\fi

\item 在解决实际问题时，往往会有不同的思路和方法，这些方法有些正确，有些错误；有些简洁，有些复杂.

问题①：设 $a\in R$，集合 $M=\{x\mid(x-1)(x-a)=0\}$，若 $x\in M$ 是 $x\in\{1,2,3\}$ 的充分条件，求：$a$ 的取值集合.

问题②：设 $x\in Z$，$a=x+3$，$b=18-x$，若 $c\in Z$，求证：$c-a$ 和 $c-b$ 至少有一个数是奇数

（1）小明在解决问题①，他认为原问题等价于 $\{1,a\}\subseteq\{1,2,3\}$，解得 $a$ 的取值集合为 $\{2,3\}$，张老师判断小明解题错误，请计算正确的 $a$ 的取值集合；

（2）小红认为既然 $c\in Z$，只需根据 $c$ 是奇数还是偶数，分类讨论即可；小华则认为可以使用反证法解决问题，请你选一种你认为更好的方法并证明.

\ansspace*[12]

\ifanswer
\solbegin
\stepm{（1）}{当 $a=1$ 时，集合 $M=\{x\mid(x-1)^2=0\}=\{1\}$，}
\step{此时 $x\in M\Rightarrow x\in\{1,2,3\}$，即 $x\in M$ 是 $x\in\{1,2,3\}$ 的充分条件.}
\step{当 $a\neq1$ 时，集合 $M=\{x\mid(x-1)(x-a)=0\}=\{1,a\}$，}
\step{由 $x\in M$ 是 $x\in\{1,2,3\}$ 的充分条件，得 $\{1,a\}\subseteq\{1,2,3\}$，}
\step{所以 $a=2$ 或 $a=3$.}
\step{综上，实数 $a$ 的取值集合为 $\{1,2,3\}$；}
\stepm{（2）}{利用反证法证明．证明如下：}
\step{假设 $c-a$、$c-b$ 均为偶数，则 $(c-b)-(c-a)=a-b$ 为偶数.}
\step{因为 $a-b=x+3-18+x=2x-15$ 为奇数，}
\step{与 $b-a$ 为偶数相矛盾，所以原假设不成立.}
\step{故 $c-a$、$c-b$ 至少有一个数是奇数.}
\solend
\fi

\item 已知实数 $x_1$、$x_2$、$x_3$、$x_4$、$x_5$，满足 $x_1+x_2+x_3+x_4+x_5=5$.

（1）求证：$x_1$、$x_2$、$x_2$、$x_4$、$x_5$ 中至少有一个实数不小于 $1$；% 原卷把 x_3 误印作 x_2，照录

（2）设 $x_1$、$x_2$、$x_3$、$x_4$、$x_5$ 这五个实数两两不等，集合 $A=\{x_1,x_2,x_3,x_4,x_5\}$，若 $B\neq\varnothing$ 且 $B\subseteq A$，记 $G(B)$ 是 $B$ 中所有元素之和，对所有的 $B$，求 $G(B)$ 的平均值.

\ansspace*[13]

\ifanswer
\solbegin
\stepm{（1）}{假设 $x_1$、$x_2$、$x_2$、$x_4$、$x_5$ 中的每一个数都小于 $1$，}% 原卷 x_3 误印作 x_2，照录
\step{得 $x_1+x_2+x_2+x_4+x_5<5$，这与 $x_1+x_2+x_3+x_4+x_5=5$ 矛盾，}% 前半句原卷 x_3 误印作 x_2，照录
\step{故 $x_1$、$x_2$、$x_2$、$x_4$、$x_5$ 中至少有一个实数不小于 $1$；}% 原卷 x_3 误印作 x_2，照录
\stepm{（2）}{集合 $A=\{x_1,x_2,x_3,x_4,x_5\}$ 的非空子集个数为 $32-1=31$，}
\step{由于 $G(B)$ 是 $B$ 中所有元素之和，}
\step{得 $G(B)=16(x_1+x_2+x_3+x_4+x_5)=16\times5=80$，}
\step{则 $G(B)$ 的平均值为 $\dfrac{80}{31}$.}
\solend
\fi

\end{enumerate}

\end{document}
"
% 紧凑版：xelatex "\def\iscompact{1}% 高一98_市西中学_10月月考.tex
%   —— 高一上数学“2026-2027 年市西中学高一上 10 月月考”（讲评版）
% 试卷版：xelatex 高一98_市西中学_10月月考.tex
% 答案版：xelatex "\def\isanswer{1}\input{高一98_市西中学_10月月考.tex}"
% 紧凑版：xelatex "\def\iscompact{1}\input{高一98_市西中学_10月月考.tex}"
%
% 原始材料：微信公众号“魔都数学加油站”文章，连载“27学年高一98”；
% 原文标题“27学年高一98——2026-2027年上海市市西中学高一上10月月考”；
% 原文链接：https://mp.weixin.qq.com/s/duK4wz-HT3W9tjOeYQaVfg。
% 正文为 8 张 PNG 页（`原图/高一98_市西中学_10月月考/01..08.png`），
% 尺寸混排（第 1—4、6、7 页为 4762×6735，第 5、8 页为 2560×3620），按页序读取。
%
% 篇幅与题号：原卷自第 3 题起编号（第 1、2 题不在原文内）。填空题 3—12（10 题）、
% 选择题 13—16（4 题）、解答题 17—21（5 题），共 19 题。
%
% 讲评版：全部 19 题均印【解析】，其中第 13—16 题的答案在解析末句
% “故选 $B$／$D$／$A$／$D$”（题干括号内未印答案）。全卷无插图。
%
% 副标题：原卷未印副标题。本稿按内容标作“集合与常用逻辑用语、不等式”。
%
% 原卷笔误/存疑（均照录未改，正文对应处另有 % 注）：
%   ① 第 21 题 (1) 求证行、解析“假设”行与“故”行三处把 $x_3$ 误印作 $x_2$
%      （“$x_1$、$x_2$、$x_2$、$x_4$、$x_5$”）；解析同一行后半句
%      “这与 $x_1+x_2+x_3+x_4+x_5=5$ 矛盾”却用 $x_3$，显系笔误，照录；
%   ② 第 6 题解析的判别式行印作“$\Delta=(-4)^2-4a\geqslant0$”，与
%      $f(1)>0$ 并列，照录。
%
% 与原卷的排版差异：
%   ① 原卷每题下直接印【解析】，本稿按全稿惯例将答案与解析置于答案版；
%   ② 原卷未印副标题，本稿补出；填空横线统一排 \blankline[4.4em]；
%   ③ 原卷数集记号排斜体 $R$、$\Z$，本稿按全稿惯例改用 $\R$、$\Z$；
%   ④ 解答题按惯例加 \ansspace；选择题括号排 \mbox{（\quad）}；
%   ⑤ 第 12 题定义的运算 $\odot$ 原卷排实心小圆，本稿排 $\odot$。
%
% 行数核对：逐题数 \step/\stepm 条数，合计 82 条（Q3：1；Q4：2；Q5：1；Q6：4；Q7：3；
% Q8：2；Q9：5；Q10：2；Q11：5；Q12：4；Q13：3；Q14：2；Q15：2；Q16：9；Q17：2；
% Q18：11；Q19：6；Q20：11；Q21：7）。全卷 19 题均印【解析】。
% 原卷解析的换行按语意合并，未逐行强拆。
%
\documentclass[a4paper,11pt]{article}
\usepackage{common}
\begin{document}

\examtitlest[7]{2026--2027 年市西中学高一上 10 月月考}{集合与常用逻辑用语、不等式}

\bigsec{一、填空题}

\begin{enumerate}
\setcounter{enumi}{2}

\item 集合 $A=\{x\mid y=\sqrt{1-x^2}\},B=\{y\mid y=x^2+1,x\in A\}$，则 $A\cap\overline B=\blankline[4.4em]$.

\ifanswer
\solbegin
\step{$1-x^2\geqslant0\Rightarrow A=[-1,1]$，从而 $B=[1,2]$，所以 $A\cap\overline B=[-1,1)$.}
\solend
\fi

\item 已知 $b,c\in R$，关于 $x$ 的不等式 $x^2-bx+c<0$ 的解集为 $(-3,2)$，则 $b+c=\blankline[4.4em]$.

\ifanswer
\solbegin
\step{由题意得 $x=-3$ 与 $x=2$ 是 $x^2-bx+c=0$ 的两根，}
\step{所以 $\begin{cases}-3+2=b\\ -3\times2=c\end{cases}$，解得 $\begin{cases}b=-1\\ c=-6\end{cases}$，所以 $b+c=-7$.}
\solend
\fi

\item 若 $\alpha:\lvert x-m\rvert<1$ 是 $\beta:x^2-4x<0$ 的充分条件，则实数 $m$ 的取值范围是\blankline[4.4em].

\ifanswer
\solbegin
\step{$(m-1,m+1)\subseteq(0,4)\Leftrightarrow m-1\geqslant0$ 且 $m+1\leqslant4\Leftrightarrow m\in[1,3]$.}
\solend
\fi

\item 已知方程 $x^2-4x+a=0$ 的两根都大于 $1$，则 $a$ 的取值范围为\blankline[4.4em].

\ifanswer
\solbegin
\step{设函数 $f(x)=x^2-4x+a$，对称轴为 $x=2$，}
\step{因为方程 $x^2-4x+a=0$ 的两根都大于 $1$，}
\step{所以 $\begin{cases}\Delta=(-4)^2-4a\geqslant0\\ f(1)>0\end{cases}$，即 $\begin{cases}16-4a\geqslant0\\ 1-4+a>0\end{cases}$，解得 $3<a\leqslant4$，}
\step{即 $a$ 的取值范围为 $(3,4]$.}
\solend
\fi

\item 若不等式 $(a-2)x^2+2(a-2)x-4<0$ 对 $x\in R$ 恒成立，则实数 $a$ 的取值范围为\blankline[4.4em].

\ifanswer
\solbegin
\step{①当 $a-2=0$，即 $a=2$ 时，不等式为 $-4<0$ 恒成立，符合题意；}
\step{②当 $a\neq2$ 时，则 $\begin{cases}a-2<0\\ \Delta=4(a-2)^2+16(a-2)<0\end{cases}$，解得 $-2<a<2$.}
\step{综上所述，实数 $a$ 的取值范围是 $(-2,2]$.}
\solend
\fi

\item 已知 $-1\leqslant a+b\leqslant4,2\leqslant a-b\leqslant3$，则 $3a-2b$ 的取值范围为\blankline[4.4em].

\ifanswer
\solbegin
\step{设 $3a-2b=m(a+b)+n(a-b)$，则 $m+n=3,m-n=-2$，}
\step{所以 $m=\dfrac12,n=\dfrac52$，所以 $3a-2b=\dfrac12(a+b)+\dfrac52(a-b)\in\left[\dfrac92,\dfrac{19}2\right]$.}
\solend
\fi

\item 设 $x_1$、$x_2$ 是关于 $x$ 的方程 $2x^2-4mx+2m^2+3m-2=0$ 的两个实数根，则 $\lvert x_1-x_2\rvert$ 的最小值为\blankline[4.4em].

\ifanswer
\solbegin
\step{$x_1$、$x_2$ 是方程 $2x^2-4mx+2m^2+3m-2=0$ 的两个实数根，}
\step{所以 $\Delta=(4m)^2-8(2m^2+3m-2)\geqslant0$，解得 $m\leqslant\dfrac23$，}
\step{由韦达定理得 $x_1+x_2=2m,x_1\cdot x_2=\dfrac{2m^2+3m-2}2$，}
\step{则 $\lvert x_1-x_2\rvert=\sqrt{(x_1+x_2)^2-4x_1\cdot x_2}=\sqrt{4-6m}$，}
\step{当 $m=\dfrac23$ 时，$\lvert x_1-x_2\rvert$ 的最小值为 $0$.}
\solend
\fi

\item 设 $x\in R$，$[x]$ 表示不大于 $x$ 的最大整数，如 $[\pi]=3$，$[-1.2]=-2$，$\left[\dfrac12\right]=0$，则使 $\lvert\lvert x^2-1\rvert\rvert=3$ 成立的 $x$ 的取值范围是\blankline[4.4em].

\ifanswer
\solbegin
\step{若 $\lvert\lvert x^2-1\rvert\rvert=3$，则 $3\leqslant\lvert x^2-1\rvert<4$，即 $4\leqslant x^2<5$，}
\step{解得 $2\leqslant x<\sqrt5$ 或 $-\sqrt5<x\leqslant-2$，故 $x$ 的取值范围是 $(-\sqrt5,-2]\cup[2,\sqrt5)$.}
\solend
\fi

\item 设 $a$ 为实数，若关于 $x$ 的一元一次不等式组 $\begin{cases}2x+a>0\\ 3x-6a<0\end{cases}$ 的解集中有且仅有 $4$ 个整数，则 $a$ 的取值范围是\blankline[4.4em].

\ifanswer
\solbegin
\step{一元一次不等式组 $\begin{cases}2x+a>0\\ 3x-6a<0\end{cases}\Leftrightarrow-\dfrac a2<x<2a$，且 $a>0$，}
\step{关于 $x$ 的一元一次不等式组 $\begin{cases}2x+a>0\\ 3x-6a<0\end{cases}$ 的解集中有且仅有 $4$ 个整数，}
\step{①当 $4$ 个整数为 $0$，$1$，$2$，$3$ 时，则 $\begin{cases}-1\leqslant-\dfrac a2<0\\ 3<2a\leqslant4\end{cases}$，所以 $\dfrac32<a\leqslant2$，}
\step{②当 $4$ 个整数为 $-1$，$0$，$1$，$2$ 时，则 $\begin{cases}-2\leqslant-\dfrac a2<-1\\ 2<2a\leqslant3\end{cases}$，无解，}
\step{综上，$a$ 的取值范围是 $\left(\dfrac32,2\right]$.}
\solend
\fi

\item 集合 $M$、$N$、$S$ 都是非空集合，现规定如下运算：

$M\odot N\odot S=\{x\mid x\in(M\cap N)\cup(N\cap S)\cup(S\cap M)$ 且 $x\notin M\cap N\cap S\}$．假设集合 $A=\{x\mid a<x<b\}$，$B=\{x\mid c<x<d\}$，$C=\{x\mid e<x<f\}$，其中实数 $a$、$b$、$c$、$d$、$e$、$f$ 满足：$a<c<e<0<b<d<f$．计算 $A\odot B\odot C=\blankline[4.4em]$.

\ifanswer
\solbegin
\step{因为集合 $A=\{x\mid a<x<b\}$，$B=\{x\mid c<x<d\}$，$C=\{x\mid e<x<f\}$，}
\step{又因为 $a<c<e<0<b<d<f$，}
\step{所以 $A\cap B=\{x\mid c<x<b\}$，$B\cap C=\{x\mid e<x<d\}$，$C\cap A=\{x\mid e<x<b\}$，}
\step{$A\cap B\cap C=\{x\mid e<x<b\}$，故 $A\odot B\odot C=\{x\mid c<x\leqslant e$ 或 $b\leqslant x<d\}$.}
\solend
\fi

\end{enumerate}

\bigsec{二、选择题}

\begin{enumerate}
\setcounter{enumi}{12}

\item “$x\neq1$”是“$x^2-1\neq0$”的\mbox{（\quad）}条件.

\keepwithprev
A. 充分非必要\qquad B. 必要非充分

C. 充要\qquad D. 既非充分也非必要

\ifanswer
\solbegin
\step{$x^2-1\neq0\Leftrightarrow x\neq1$ 且 $x\neq-1\Rightarrow x\neq1$.}
\step{反之不成立，如 $x=-1$ 时，$x\neq1$，但 $x^2-1=0$.}
\step{因此“$x\neq1$”是“$x^2-1\neq0$”的必要非充分条件，故选 $B$.}
\solend
\fi

\item 若 $\dfrac1a<\dfrac1b<0$，则下列结论中不正确的是\mbox{（\quad）}

\keepwithprev
A. $a+b<0$\qquad B. $a^2<b^2$\qquad C. $ab<b^2$\qquad D. $ab<a^2$

\ifanswer
\solbegin
\step{由 $\dfrac1a<\dfrac1b<0$ 得 $b<a<0$，因此 $a+b<0$，且 $a^2<ab<b^2$，}
\step{所以 A、B、C 正确，D 不正确．故选 $D$.}
\solend
\fi

\item 非空集合 $X=\{x\mid a+1\leqslant x\leqslant3a-5\}$，$Y=\{x\mid1\leqslant x\leqslant16\}$，使得 $X\subseteq(X\cap Y)$ 成立的所有 $a$ 的集合是\mbox{（\quad）}

\keepwithprev
A. $\{a\mid3\leqslant a\leqslant7\}$\qquad B. $\{a\mid0\leqslant a\leqslant7\}$\qquad C. $\{a\mid3<a\leqslant7\}$\qquad D. $\{a\mid a\leqslant7\}$

\ifanswer
\solbegin
\step{因为 $X\subseteq(X\cap Y)$ 成立，所以 $X\subseteq Y$ 成立.}
\step{$X\neq\varnothing$，要使 $X\subseteq Y$，则 $\begin{cases}a\geqslant3\\ a+1\geqslant1\\ 3a-5\leqslant16\end{cases}$，即 $\begin{cases}a\geqslant3\\ a\geqslant0\\ a\leqslant7\end{cases}$，所以 $3\leqslant a\leqslant7$，故选 $A$.}
\solend
\fi

\item 已知 $M$ 是满足下列条件的集合：①$0\in M$，$1\in M$；②若 $x$、$y\in M$，则 $x-y\in M$；③若 $x\in M$ 且 $x\neq0$，则 $\dfrac1x\in M$．则下列说法正确的个数为\mbox{（\quad）}

（1）$\dfrac12\in M$；（2）$x+y\in M$；（3）$x^2\in M$．

\keepwithprev
A. $0$\qquad B. $1$\qquad C. $2$\qquad D. $3$

\ifanswer
\solbegin
\step{（1）$\dfrac12\in M$ 正确，理由如下：}
\step{由①得 $0\in M$，$1\in M$，由②得 $0-1=-1\in M$，$1-(-1)=2\in M$，}
\step{由③得 $\dfrac12\in M$；}
\step{（2）由①得 $0\in M$，由题意 $y\in M$，}
\step{所以由②得 $0-y=-y\in M$，又 $x\in M$，所以 $x-(-y)=x+y\in M$；}
\step{（3）$x\in M$，由②得 $x-1\in M$，由③得 $\dfrac1x\in M$，$\dfrac1{x-1}\in M$，}
\step{所以 $\dfrac1x-\dfrac1{x-1}\in M$，即 $\dfrac1{x(x-1)}\in M$，所以 $x(x-1)\in M$，}
\step{由（2）结论得 $x(x-1)+x\in M$，即 $x^2\in M$，}
\step{所以说法正确的个数为 $3$．故选 $D$.}
\solend
\fi

\end{enumerate}

\bigsec{三、解答题}

\begin{enumerate}
\setcounter{enumi}{16}

\item 已知 $a>b$，用比较法证明：$a^3>b^3$.

\ansspace*[8]

\ifanswer
\solbegin
\step{$a>b$，$a^3-b^3=(a-b)(a^2+ab+b^2)=(a-b)\left[\left(a+\dfrac b2\right)^2+\dfrac34b^2\right]>0$，}
\step{故 $a^3>b^3$.}
\solend
\fi

\item 求下列关于 $x$ 的不等式的解集：

（1）$\dfrac{x^2(x+3)(x^2-x+1)}{x^2-5x+4}>0$；

（2）$\dfrac{ax-1}{x-2}\geqslant0\left(a<\dfrac12\right)$.

\ansspace*[12]

\ifanswer
\solbegin
\stepm{（1）}{分母 $x^2-5x+4=(x-1)(x-4)$，要求 $x\neq1,4$；}
\step{又 $x=0$ 时原不等式不成立，且 $x^2-x+1=\left(x-\dfrac12\right)^2+\dfrac34>0$，}
\step{因此原不等式等价于 $\dfrac{x+3}{(x-1)(x-4)}>0,x\neq0$.}
\step{在数轴上标出 $-3,1,4$，分段判号得 $-3<x<1$ 或 $x>4$，}
\step{再排除 $x=0$，解集为 $(-3,0)\cup(0,1)\cup(4,+\infty)$.}
\stepm{（2）}{原不等式要求 $x\neq2$，按 $a$ 的符号分类讨论.}
\step{①当 $a<0$ 时，$\dfrac1a<2$，原不等式等价于 $\dfrac{x-\dfrac1a}{x-2}\leqslant0\Leftrightarrow\dfrac1a\leqslant x<2$，}
\step{解集为 $\left[\dfrac1a,2\right)$.}
\step{②当 $a=0$ 时，原不等式为 $\dfrac{-1}{x-2}\geqslant0\Leftrightarrow x<2$，解集为 $(-\infty,2)$.}
\step{③当 $0<a<\dfrac12$ 时，$\dfrac1a>2$，}
\step{原不等式等价于 $\dfrac{x-\dfrac1a}{x-2}\geqslant0\Leftrightarrow x<2$ 或 $x\geqslant\dfrac1a$，解集为 $(-\infty,2)\cup\left[\dfrac1a,+\infty\right)$.}
\solend
\fi

\item 集合 $A=\{x\mid-2\leqslant x\leqslant5\}$，$B=\{x\mid m+1\leqslant x\leqslant2m-1\}$.

（1）若 $B\subseteq A$，求实数 $m$ 的取值范围；

（2）若 $A\cap B\neq\varnothing$，求实数 $m$ 的取值范围.

\ansspace*[10]

\ifanswer
\solbegin
\stepm{（1）}{当 $B=\varnothing$ 时，$m+1>2m-1$，解得 $m<2$，}
\step{当 $B\neq\varnothing$ 时，$\begin{cases}m+1\leqslant2m-1\\ m+1\geqslant-2\\ 2m-1\leqslant5\end{cases}$，解得 $2\leqslant m\leqslant3$，}
\step{综上，实数 $m$ 的取值范围是 $(-\infty,3]$；}
\stepm{（2）}{当 $B=\varnothing$ 时，$m+1>2m-1$，解得 $m<2$，}
\step{当 $B\neq\varnothing$ 时，$\begin{cases}m+1\leqslant2m-1\\ 2m-1<-2\end{cases}$ 或 $\begin{cases}m+1\leqslant2m-1\\ m+1>5\end{cases}$，解得 $m\in\varnothing$ 或 $m>4$.}
\step{所以当 $A\cap B\neq\varnothing$ 时，实数 $m$ 的取值范围是 $[2,4]$.}
\solend
\fi

\item 在解决实际问题时，往往会有不同的思路和方法，这些方法有些正确，有些错误；有些简洁，有些复杂.

问题①：设 $a\in R$，集合 $M=\{x\mid(x-1)(x-a)=0\}$，若 $x\in M$ 是 $x\in\{1,2,3\}$ 的充分条件，求：$a$ 的取值集合.

问题②：设 $x\in Z$，$a=x+3$，$b=18-x$，若 $c\in Z$，求证：$c-a$ 和 $c-b$ 至少有一个数是奇数

（1）小明在解决问题①，他认为原问题等价于 $\{1,a\}\subseteq\{1,2,3\}$，解得 $a$ 的取值集合为 $\{2,3\}$，张老师判断小明解题错误，请计算正确的 $a$ 的取值集合；

（2）小红认为既然 $c\in Z$，只需根据 $c$ 是奇数还是偶数，分类讨论即可；小华则认为可以使用反证法解决问题，请你选一种你认为更好的方法并证明.

\ansspace*[12]

\ifanswer
\solbegin
\stepm{（1）}{当 $a=1$ 时，集合 $M=\{x\mid(x-1)^2=0\}=\{1\}$，}
\step{此时 $x\in M\Rightarrow x\in\{1,2,3\}$，即 $x\in M$ 是 $x\in\{1,2,3\}$ 的充分条件.}
\step{当 $a\neq1$ 时，集合 $M=\{x\mid(x-1)(x-a)=0\}=\{1,a\}$，}
\step{由 $x\in M$ 是 $x\in\{1,2,3\}$ 的充分条件，得 $\{1,a\}\subseteq\{1,2,3\}$，}
\step{所以 $a=2$ 或 $a=3$.}
\step{综上，实数 $a$ 的取值集合为 $\{1,2,3\}$；}
\stepm{（2）}{利用反证法证明．证明如下：}
\step{假设 $c-a$、$c-b$ 均为偶数，则 $(c-b)-(c-a)=a-b$ 为偶数.}
\step{因为 $a-b=x+3-18+x=2x-15$ 为奇数，}
\step{与 $b-a$ 为偶数相矛盾，所以原假设不成立.}
\step{故 $c-a$、$c-b$ 至少有一个数是奇数.}
\solend
\fi

\item 已知实数 $x_1$、$x_2$、$x_3$、$x_4$、$x_5$，满足 $x_1+x_2+x_3+x_4+x_5=5$.

（1）求证：$x_1$、$x_2$、$x_2$、$x_4$、$x_5$ 中至少有一个实数不小于 $1$；% 原卷把 x_3 误印作 x_2，照录

（2）设 $x_1$、$x_2$、$x_3$、$x_4$、$x_5$ 这五个实数两两不等，集合 $A=\{x_1,x_2,x_3,x_4,x_5\}$，若 $B\neq\varnothing$ 且 $B\subseteq A$，记 $G(B)$ 是 $B$ 中所有元素之和，对所有的 $B$，求 $G(B)$ 的平均值.

\ansspace*[13]

\ifanswer
\solbegin
\stepm{（1）}{假设 $x_1$、$x_2$、$x_2$、$x_4$、$x_5$ 中的每一个数都小于 $1$，}% 原卷 x_3 误印作 x_2，照录
\step{得 $x_1+x_2+x_2+x_4+x_5<5$，这与 $x_1+x_2+x_3+x_4+x_5=5$ 矛盾，}% 前半句原卷 x_3 误印作 x_2，照录
\step{故 $x_1$、$x_2$、$x_2$、$x_4$、$x_5$ 中至少有一个实数不小于 $1$；}% 原卷 x_3 误印作 x_2，照录
\stepm{（2）}{集合 $A=\{x_1,x_2,x_3,x_4,x_5\}$ 的非空子集个数为 $32-1=31$，}
\step{由于 $G(B)$ 是 $B$ 中所有元素之和，}
\step{得 $G(B)=16(x_1+x_2+x_3+x_4+x_5)=16\times5=80$，}
\step{则 $G(B)$ 的平均值为 $\dfrac{80}{31}$.}
\solend
\fi

\end{enumerate}

\end{document}
"
%
% 原始材料：微信公众号“魔都数学加油站”文章，连载“27学年高一98”；
% 原文标题“27学年高一98——2026-2027年上海市市西中学高一上10月月考”；
% 原文链接：https://mp.weixin.qq.com/s/duK4wz-HT3W9tjOeYQaVfg。
% 正文为 8 张 PNG 页（`原图/高一98_市西中学_10月月考/01..08.png`），
% 尺寸混排（第 1—4、6、7 页为 4762×6735，第 5、8 页为 2560×3620），按页序读取。
%
% 篇幅与题号：原卷自第 3 题起编号（第 1、2 题不在原文内）。填空题 3—12（10 题）、
% 选择题 13—16（4 题）、解答题 17—21（5 题），共 19 题。
%
% 讲评版：全部 19 题均印【解析】，其中第 13—16 题的答案在解析末句
% “故选 $B$／$D$／$A$／$D$”（题干括号内未印答案）。全卷无插图。
%
% 副标题：原卷未印副标题。本稿按内容标作“集合与常用逻辑用语、不等式”。
%
% 原卷笔误/存疑（均照录未改，正文对应处另有 % 注）：
%   ① 第 21 题 (1) 求证行、解析“假设”行与“故”行三处把 $x_3$ 误印作 $x_2$
%      （“$x_1$、$x_2$、$x_2$、$x_4$、$x_5$”）；解析同一行后半句
%      “这与 $x_1+x_2+x_3+x_4+x_5=5$ 矛盾”却用 $x_3$，显系笔误，照录；
%   ② 第 6 题解析的判别式行印作“$\Delta=(-4)^2-4a\geqslant0$”，与
%      $f(1)>0$ 并列，照录。
%
% 与原卷的排版差异：
%   ① 原卷每题下直接印【解析】，本稿按全稿惯例将答案与解析置于答案版；
%   ② 原卷未印副标题，本稿补出；填空横线统一排 \blankline[4.4em]；
%   ③ 原卷数集记号排斜体 $R$、$\Z$，本稿按全稿惯例改用 $\R$、$\Z$；
%   ④ 解答题按惯例加 \ansspace；选择题括号排 \mbox{（\quad）}；
%   ⑤ 第 12 题定义的运算 $\odot$ 原卷排实心小圆，本稿排 $\odot$。
%
% 行数核对：逐题数 \step/\stepm 条数，合计 82 条（Q3：1；Q4：2；Q5：1；Q6：4；Q7：3；
% Q8：2；Q9：5；Q10：2；Q11：5；Q12：4；Q13：3；Q14：2；Q15：2；Q16：9；Q17：2；
% Q18：11；Q19：6；Q20：11；Q21：7）。全卷 19 题均印【解析】。
% 原卷解析的换行按语意合并，未逐行强拆。
%
\documentclass[a4paper,11pt]{article}
\usepackage{common}
\begin{document}

\examtitlest[7]{2026--2027 年市西中学高一上 10 月月考}{集合与常用逻辑用语、不等式}

\bigsec{一、填空题}

\begin{enumerate}
\setcounter{enumi}{2}

\item 集合 $A=\{x\mid y=\sqrt{1-x^2}\},B=\{y\mid y=x^2+1,x\in A\}$，则 $A\cap\overline B=\blankline[4.4em]$.

\ifanswer
\solbegin
\step{$1-x^2\geqslant0\Rightarrow A=[-1,1]$，从而 $B=[1,2]$，所以 $A\cap\overline B=[-1,1)$.}
\solend
\fi

\item 已知 $b,c\in R$，关于 $x$ 的不等式 $x^2-bx+c<0$ 的解集为 $(-3,2)$，则 $b+c=\blankline[4.4em]$.

\ifanswer
\solbegin
\step{由题意得 $x=-3$ 与 $x=2$ 是 $x^2-bx+c=0$ 的两根，}
\step{所以 $\begin{cases}-3+2=b\\ -3\times2=c\end{cases}$，解得 $\begin{cases}b=-1\\ c=-6\end{cases}$，所以 $b+c=-7$.}
\solend
\fi

\item 若 $\alpha:\lvert x-m\rvert<1$ 是 $\beta:x^2-4x<0$ 的充分条件，则实数 $m$ 的取值范围是\blankline[4.4em].

\ifanswer
\solbegin
\step{$(m-1,m+1)\subseteq(0,4)\Leftrightarrow m-1\geqslant0$ 且 $m+1\leqslant4\Leftrightarrow m\in[1,3]$.}
\solend
\fi

\item 已知方程 $x^2-4x+a=0$ 的两根都大于 $1$，则 $a$ 的取值范围为\blankline[4.4em].

\ifanswer
\solbegin
\step{设函数 $f(x)=x^2-4x+a$，对称轴为 $x=2$，}
\step{因为方程 $x^2-4x+a=0$ 的两根都大于 $1$，}
\step{所以 $\begin{cases}\Delta=(-4)^2-4a\geqslant0\\ f(1)>0\end{cases}$，即 $\begin{cases}16-4a\geqslant0\\ 1-4+a>0\end{cases}$，解得 $3<a\leqslant4$，}
\step{即 $a$ 的取值范围为 $(3,4]$.}
\solend
\fi

\item 若不等式 $(a-2)x^2+2(a-2)x-4<0$ 对 $x\in R$ 恒成立，则实数 $a$ 的取值范围为\blankline[4.4em].

\ifanswer
\solbegin
\step{①当 $a-2=0$，即 $a=2$ 时，不等式为 $-4<0$ 恒成立，符合题意；}
\step{②当 $a\neq2$ 时，则 $\begin{cases}a-2<0\\ \Delta=4(a-2)^2+16(a-2)<0\end{cases}$，解得 $-2<a<2$.}
\step{综上所述，实数 $a$ 的取值范围是 $(-2,2]$.}
\solend
\fi

\item 已知 $-1\leqslant a+b\leqslant4,2\leqslant a-b\leqslant3$，则 $3a-2b$ 的取值范围为\blankline[4.4em].

\ifanswer
\solbegin
\step{设 $3a-2b=m(a+b)+n(a-b)$，则 $m+n=3,m-n=-2$，}
\step{所以 $m=\dfrac12,n=\dfrac52$，所以 $3a-2b=\dfrac12(a+b)+\dfrac52(a-b)\in\left[\dfrac92,\dfrac{19}2\right]$.}
\solend
\fi

\item 设 $x_1$、$x_2$ 是关于 $x$ 的方程 $2x^2-4mx+2m^2+3m-2=0$ 的两个实数根，则 $\lvert x_1-x_2\rvert$ 的最小值为\blankline[4.4em].

\ifanswer
\solbegin
\step{$x_1$、$x_2$ 是方程 $2x^2-4mx+2m^2+3m-2=0$ 的两个实数根，}
\step{所以 $\Delta=(4m)^2-8(2m^2+3m-2)\geqslant0$，解得 $m\leqslant\dfrac23$，}
\step{由韦达定理得 $x_1+x_2=2m,x_1\cdot x_2=\dfrac{2m^2+3m-2}2$，}
\step{则 $\lvert x_1-x_2\rvert=\sqrt{(x_1+x_2)^2-4x_1\cdot x_2}=\sqrt{4-6m}$，}
\step{当 $m=\dfrac23$ 时，$\lvert x_1-x_2\rvert$ 的最小值为 $0$.}
\solend
\fi

\item 设 $x\in R$，$[x]$ 表示不大于 $x$ 的最大整数，如 $[\pi]=3$，$[-1.2]=-2$，$\left[\dfrac12\right]=0$，则使 $\lvert\lvert x^2-1\rvert\rvert=3$ 成立的 $x$ 的取值范围是\blankline[4.4em].

\ifanswer
\solbegin
\step{若 $\lvert\lvert x^2-1\rvert\rvert=3$，则 $3\leqslant\lvert x^2-1\rvert<4$，即 $4\leqslant x^2<5$，}
\step{解得 $2\leqslant x<\sqrt5$ 或 $-\sqrt5<x\leqslant-2$，故 $x$ 的取值范围是 $(-\sqrt5,-2]\cup[2,\sqrt5)$.}
\solend
\fi

\item 设 $a$ 为实数，若关于 $x$ 的一元一次不等式组 $\begin{cases}2x+a>0\\ 3x-6a<0\end{cases}$ 的解集中有且仅有 $4$ 个整数，则 $a$ 的取值范围是\blankline[4.4em].

\ifanswer
\solbegin
\step{一元一次不等式组 $\begin{cases}2x+a>0\\ 3x-6a<0\end{cases}\Leftrightarrow-\dfrac a2<x<2a$，且 $a>0$，}
\step{关于 $x$ 的一元一次不等式组 $\begin{cases}2x+a>0\\ 3x-6a<0\end{cases}$ 的解集中有且仅有 $4$ 个整数，}
\step{①当 $4$ 个整数为 $0$，$1$，$2$，$3$ 时，则 $\begin{cases}-1\leqslant-\dfrac a2<0\\ 3<2a\leqslant4\end{cases}$，所以 $\dfrac32<a\leqslant2$，}
\step{②当 $4$ 个整数为 $-1$，$0$，$1$，$2$ 时，则 $\begin{cases}-2\leqslant-\dfrac a2<-1\\ 2<2a\leqslant3\end{cases}$，无解，}
\step{综上，$a$ 的取值范围是 $\left(\dfrac32,2\right]$.}
\solend
\fi

\item 集合 $M$、$N$、$S$ 都是非空集合，现规定如下运算：

$M\odot N\odot S=\{x\mid x\in(M\cap N)\cup(N\cap S)\cup(S\cap M)$ 且 $x\notin M\cap N\cap S\}$．假设集合 $A=\{x\mid a<x<b\}$，$B=\{x\mid c<x<d\}$，$C=\{x\mid e<x<f\}$，其中实数 $a$、$b$、$c$、$d$、$e$、$f$ 满足：$a<c<e<0<b<d<f$．计算 $A\odot B\odot C=\blankline[4.4em]$.

\ifanswer
\solbegin
\step{因为集合 $A=\{x\mid a<x<b\}$，$B=\{x\mid c<x<d\}$，$C=\{x\mid e<x<f\}$，}
\step{又因为 $a<c<e<0<b<d<f$，}
\step{所以 $A\cap B=\{x\mid c<x<b\}$，$B\cap C=\{x\mid e<x<d\}$，$C\cap A=\{x\mid e<x<b\}$，}
\step{$A\cap B\cap C=\{x\mid e<x<b\}$，故 $A\odot B\odot C=\{x\mid c<x\leqslant e$ 或 $b\leqslant x<d\}$.}
\solend
\fi

\end{enumerate}

\bigsec{二、选择题}

\begin{enumerate}
\setcounter{enumi}{12}

\item “$x\neq1$”是“$x^2-1\neq0$”的\mbox{（\quad）}条件.

\keepwithprev
A. 充分非必要\qquad B. 必要非充分

C. 充要\qquad D. 既非充分也非必要

\ifanswer
\solbegin
\step{$x^2-1\neq0\Leftrightarrow x\neq1$ 且 $x\neq-1\Rightarrow x\neq1$.}
\step{反之不成立，如 $x=-1$ 时，$x\neq1$，但 $x^2-1=0$.}
\step{因此“$x\neq1$”是“$x^2-1\neq0$”的必要非充分条件，故选 $B$.}
\solend
\fi

\item 若 $\dfrac1a<\dfrac1b<0$，则下列结论中不正确的是\mbox{（\quad）}

\keepwithprev
A. $a+b<0$\qquad B. $a^2<b^2$\qquad C. $ab<b^2$\qquad D. $ab<a^2$

\ifanswer
\solbegin
\step{由 $\dfrac1a<\dfrac1b<0$ 得 $b<a<0$，因此 $a+b<0$，且 $a^2<ab<b^2$，}
\step{所以 A、B、C 正确，D 不正确．故选 $D$.}
\solend
\fi

\item 非空集合 $X=\{x\mid a+1\leqslant x\leqslant3a-5\}$，$Y=\{x\mid1\leqslant x\leqslant16\}$，使得 $X\subseteq(X\cap Y)$ 成立的所有 $a$ 的集合是\mbox{（\quad）}

\keepwithprev
A. $\{a\mid3\leqslant a\leqslant7\}$\qquad B. $\{a\mid0\leqslant a\leqslant7\}$\qquad C. $\{a\mid3<a\leqslant7\}$\qquad D. $\{a\mid a\leqslant7\}$

\ifanswer
\solbegin
\step{因为 $X\subseteq(X\cap Y)$ 成立，所以 $X\subseteq Y$ 成立.}
\step{$X\neq\varnothing$，要使 $X\subseteq Y$，则 $\begin{cases}a\geqslant3\\ a+1\geqslant1\\ 3a-5\leqslant16\end{cases}$，即 $\begin{cases}a\geqslant3\\ a\geqslant0\\ a\leqslant7\end{cases}$，所以 $3\leqslant a\leqslant7$，故选 $A$.}
\solend
\fi

\item 已知 $M$ 是满足下列条件的集合：①$0\in M$，$1\in M$；②若 $x$、$y\in M$，则 $x-y\in M$；③若 $x\in M$ 且 $x\neq0$，则 $\dfrac1x\in M$．则下列说法正确的个数为\mbox{（\quad）}

（1）$\dfrac12\in M$；（2）$x+y\in M$；（3）$x^2\in M$．

\keepwithprev
A. $0$\qquad B. $1$\qquad C. $2$\qquad D. $3$

\ifanswer
\solbegin
\step{（1）$\dfrac12\in M$ 正确，理由如下：}
\step{由①得 $0\in M$，$1\in M$，由②得 $0-1=-1\in M$，$1-(-1)=2\in M$，}
\step{由③得 $\dfrac12\in M$；}
\step{（2）由①得 $0\in M$，由题意 $y\in M$，}
\step{所以由②得 $0-y=-y\in M$，又 $x\in M$，所以 $x-(-y)=x+y\in M$；}
\step{（3）$x\in M$，由②得 $x-1\in M$，由③得 $\dfrac1x\in M$，$\dfrac1{x-1}\in M$，}
\step{所以 $\dfrac1x-\dfrac1{x-1}\in M$，即 $\dfrac1{x(x-1)}\in M$，所以 $x(x-1)\in M$，}
\step{由（2）结论得 $x(x-1)+x\in M$，即 $x^2\in M$，}
\step{所以说法正确的个数为 $3$．故选 $D$.}
\solend
\fi

\end{enumerate}

\bigsec{三、解答题}

\begin{enumerate}
\setcounter{enumi}{16}

\item 已知 $a>b$，用比较法证明：$a^3>b^3$.

\ansspace*[8]

\ifanswer
\solbegin
\step{$a>b$，$a^3-b^3=(a-b)(a^2+ab+b^2)=(a-b)\left[\left(a+\dfrac b2\right)^2+\dfrac34b^2\right]>0$，}
\step{故 $a^3>b^3$.}
\solend
\fi

\item 求下列关于 $x$ 的不等式的解集：

（1）$\dfrac{x^2(x+3)(x^2-x+1)}{x^2-5x+4}>0$；

（2）$\dfrac{ax-1}{x-2}\geqslant0\left(a<\dfrac12\right)$.

\ansspace*[12]

\ifanswer
\solbegin
\stepm{（1）}{分母 $x^2-5x+4=(x-1)(x-4)$，要求 $x\neq1,4$；}
\step{又 $x=0$ 时原不等式不成立，且 $x^2-x+1=\left(x-\dfrac12\right)^2+\dfrac34>0$，}
\step{因此原不等式等价于 $\dfrac{x+3}{(x-1)(x-4)}>0,x\neq0$.}
\step{在数轴上标出 $-3,1,4$，分段判号得 $-3<x<1$ 或 $x>4$，}
\step{再排除 $x=0$，解集为 $(-3,0)\cup(0,1)\cup(4,+\infty)$.}
\stepm{（2）}{原不等式要求 $x\neq2$，按 $a$ 的符号分类讨论.}
\step{①当 $a<0$ 时，$\dfrac1a<2$，原不等式等价于 $\dfrac{x-\dfrac1a}{x-2}\leqslant0\Leftrightarrow\dfrac1a\leqslant x<2$，}
\step{解集为 $\left[\dfrac1a,2\right)$.}
\step{②当 $a=0$ 时，原不等式为 $\dfrac{-1}{x-2}\geqslant0\Leftrightarrow x<2$，解集为 $(-\infty,2)$.}
\step{③当 $0<a<\dfrac12$ 时，$\dfrac1a>2$，}
\step{原不等式等价于 $\dfrac{x-\dfrac1a}{x-2}\geqslant0\Leftrightarrow x<2$ 或 $x\geqslant\dfrac1a$，解集为 $(-\infty,2)\cup\left[\dfrac1a,+\infty\right)$.}
\solend
\fi

\item 集合 $A=\{x\mid-2\leqslant x\leqslant5\}$，$B=\{x\mid m+1\leqslant x\leqslant2m-1\}$.

（1）若 $B\subseteq A$，求实数 $m$ 的取值范围；

（2）若 $A\cap B\neq\varnothing$，求实数 $m$ 的取值范围.

\ansspace*[10]

\ifanswer
\solbegin
\stepm{（1）}{当 $B=\varnothing$ 时，$m+1>2m-1$，解得 $m<2$，}
\step{当 $B\neq\varnothing$ 时，$\begin{cases}m+1\leqslant2m-1\\ m+1\geqslant-2\\ 2m-1\leqslant5\end{cases}$，解得 $2\leqslant m\leqslant3$，}
\step{综上，实数 $m$ 的取值范围是 $(-\infty,3]$；}
\stepm{（2）}{当 $B=\varnothing$ 时，$m+1>2m-1$，解得 $m<2$，}
\step{当 $B\neq\varnothing$ 时，$\begin{cases}m+1\leqslant2m-1\\ 2m-1<-2\end{cases}$ 或 $\begin{cases}m+1\leqslant2m-1\\ m+1>5\end{cases}$，解得 $m\in\varnothing$ 或 $m>4$.}
\step{所以当 $A\cap B\neq\varnothing$ 时，实数 $m$ 的取值范围是 $[2,4]$.}
\solend
\fi

\item 在解决实际问题时，往往会有不同的思路和方法，这些方法有些正确，有些错误；有些简洁，有些复杂.

问题①：设 $a\in R$，集合 $M=\{x\mid(x-1)(x-a)=0\}$，若 $x\in M$ 是 $x\in\{1,2,3\}$ 的充分条件，求：$a$ 的取值集合.

问题②：设 $x\in Z$，$a=x+3$，$b=18-x$，若 $c\in Z$，求证：$c-a$ 和 $c-b$ 至少有一个数是奇数

（1）小明在解决问题①，他认为原问题等价于 $\{1,a\}\subseteq\{1,2,3\}$，解得 $a$ 的取值集合为 $\{2,3\}$，张老师判断小明解题错误，请计算正确的 $a$ 的取值集合；

（2）小红认为既然 $c\in Z$，只需根据 $c$ 是奇数还是偶数，分类讨论即可；小华则认为可以使用反证法解决问题，请你选一种你认为更好的方法并证明.

\ansspace*[12]

\ifanswer
\solbegin
\stepm{（1）}{当 $a=1$ 时，集合 $M=\{x\mid(x-1)^2=0\}=\{1\}$，}
\step{此时 $x\in M\Rightarrow x\in\{1,2,3\}$，即 $x\in M$ 是 $x\in\{1,2,3\}$ 的充分条件.}
\step{当 $a\neq1$ 时，集合 $M=\{x\mid(x-1)(x-a)=0\}=\{1,a\}$，}
\step{由 $x\in M$ 是 $x\in\{1,2,3\}$ 的充分条件，得 $\{1,a\}\subseteq\{1,2,3\}$，}
\step{所以 $a=2$ 或 $a=3$.}
\step{综上，实数 $a$ 的取值集合为 $\{1,2,3\}$；}
\stepm{（2）}{利用反证法证明．证明如下：}
\step{假设 $c-a$、$c-b$ 均为偶数，则 $(c-b)-(c-a)=a-b$ 为偶数.}
\step{因为 $a-b=x+3-18+x=2x-15$ 为奇数，}
\step{与 $b-a$ 为偶数相矛盾，所以原假设不成立.}
\step{故 $c-a$、$c-b$ 至少有一个数是奇数.}
\solend
\fi

\item 已知实数 $x_1$、$x_2$、$x_3$、$x_4$、$x_5$，满足 $x_1+x_2+x_3+x_4+x_5=5$.

（1）求证：$x_1$、$x_2$、$x_2$、$x_4$、$x_5$ 中至少有一个实数不小于 $1$；% 原卷把 x_3 误印作 x_2，照录

（2）设 $x_1$、$x_2$、$x_3$、$x_4$、$x_5$ 这五个实数两两不等，集合 $A=\{x_1,x_2,x_3,x_4,x_5\}$，若 $B\neq\varnothing$ 且 $B\subseteq A$，记 $G(B)$ 是 $B$ 中所有元素之和，对所有的 $B$，求 $G(B)$ 的平均值.

\ansspace*[13]

\ifanswer
\solbegin
\stepm{（1）}{假设 $x_1$、$x_2$、$x_2$、$x_4$、$x_5$ 中的每一个数都小于 $1$，}% 原卷 x_3 误印作 x_2，照录
\step{得 $x_1+x_2+x_2+x_4+x_5<5$，这与 $x_1+x_2+x_3+x_4+x_5=5$ 矛盾，}% 前半句原卷 x_3 误印作 x_2，照录
\step{故 $x_1$、$x_2$、$x_2$、$x_4$、$x_5$ 中至少有一个实数不小于 $1$；}% 原卷 x_3 误印作 x_2，照录
\stepm{（2）}{集合 $A=\{x_1,x_2,x_3,x_4,x_5\}$ 的非空子集个数为 $32-1=31$，}
\step{由于 $G(B)$ 是 $B$ 中所有元素之和，}
\step{得 $G(B)=16(x_1+x_2+x_3+x_4+x_5)=16\times5=80$，}
\step{则 $G(B)$ 的平均值为 $\dfrac{80}{31}$.}
\solend
\fi

\end{enumerate}

\end{document}
"
%
% 原始材料：微信公众号“魔都数学加油站”文章，连载“27学年高一98”；
% 原文标题“27学年高一98——2026-2027年上海市市西中学高一上10月月考”；
% 原文链接：https://mp.weixin.qq.com/s/duK4wz-HT3W9tjOeYQaVfg。
% 正文为 8 张 PNG 页（`原图/高一98_市西中学_10月月考/01..08.png`），
% 尺寸混排（第 1—4、6、7 页为 4762×6735，第 5、8 页为 2560×3620），按页序读取。
%
% 篇幅与题号：原卷自第 3 题起编号（第 1、2 题不在原文内）。填空题 3—12（10 题）、
% 选择题 13—16（4 题）、解答题 17—21（5 题），共 19 题。
%
% 讲评版：全部 19 题均印【解析】，其中第 13—16 题的答案在解析末句
% “故选 $B$／$D$／$A$／$D$”（题干括号内未印答案）。全卷无插图。
%
% 副标题：原卷未印副标题。本稿按内容标作“集合与常用逻辑用语、不等式”。
%
% 原卷笔误/存疑（均照录未改，正文对应处另有 % 注）：
%   ① 第 21 题 (1) 求证行、解析“假设”行与“故”行三处把 $x_3$ 误印作 $x_2$
%      （“$x_1$、$x_2$、$x_2$、$x_4$、$x_5$”）；解析同一行后半句
%      “这与 $x_1+x_2+x_3+x_4+x_5=5$ 矛盾”却用 $x_3$，显系笔误，照录；
%   ② 第 6 题解析的判别式行印作“$\Delta=(-4)^2-4a\geqslant0$”，与
%      $f(1)>0$ 并列，照录。
%
% 与原卷的排版差异：
%   ① 原卷每题下直接印【解析】，本稿按全稿惯例将答案与解析置于答案版；
%   ② 原卷未印副标题，本稿补出；填空横线统一排 \blankline[4.4em]；
%   ③ 原卷数集记号排斜体 $R$、$\Z$，本稿按全稿惯例改用 $\R$、$\Z$；
%   ④ 解答题按惯例加 \ansspace；选择题括号排 \mbox{（\quad）}；
%   ⑤ 第 12 题定义的运算 $\odot$ 原卷排实心小圆，本稿排 $\odot$。
%
% 行数核对：逐题数 \step/\stepm 条数，合计 82 条（Q3：1；Q4：2；Q5：1；Q6：4；Q7：3；
% Q8：2；Q9：5；Q10：2；Q11：5；Q12：4；Q13：3；Q14：2；Q15：2；Q16：9；Q17：2；
% Q18：11；Q19：6；Q20：11；Q21：7）。全卷 19 题均印【解析】。
% 原卷解析的换行按语意合并，未逐行强拆。
%
\documentclass[a4paper,11pt]{article}
\usepackage{common}
\begin{document}

\examtitlest[7]{2026--2027 年市西中学高一上 10 月月考}{集合与常用逻辑用语、不等式}

\bigsec{一、填空题}

\begin{enumerate}
\setcounter{enumi}{2}

\item 集合 $A=\{x\mid y=\sqrt{1-x^2}\},B=\{y\mid y=x^2+1,x\in A\}$，则 $A\cap\overline B=\blankline[4.4em]$.

\ifanswer
\solbegin
\step{$1-x^2\geqslant0\Rightarrow A=[-1,1]$，从而 $B=[1,2]$，所以 $A\cap\overline B=[-1,1)$.}
\solend
\fi

\item 已知 $b,c\in R$，关于 $x$ 的不等式 $x^2-bx+c<0$ 的解集为 $(-3,2)$，则 $b+c=\blankline[4.4em]$.

\ifanswer
\solbegin
\step{由题意得 $x=-3$ 与 $x=2$ 是 $x^2-bx+c=0$ 的两根，}
\step{所以 $\begin{cases}-3+2=b\\ -3\times2=c\end{cases}$，解得 $\begin{cases}b=-1\\ c=-6\end{cases}$，所以 $b+c=-7$.}
\solend
\fi

\item 若 $\alpha:\lvert x-m\rvert<1$ 是 $\beta:x^2-4x<0$ 的充分条件，则实数 $m$ 的取值范围是\blankline[4.4em].

\ifanswer
\solbegin
\step{$(m-1,m+1)\subseteq(0,4)\Leftrightarrow m-1\geqslant0$ 且 $m+1\leqslant4\Leftrightarrow m\in[1,3]$.}
\solend
\fi

\item 已知方程 $x^2-4x+a=0$ 的两根都大于 $1$，则 $a$ 的取值范围为\blankline[4.4em].

\ifanswer
\solbegin
\step{设函数 $f(x)=x^2-4x+a$，对称轴为 $x=2$，}
\step{因为方程 $x^2-4x+a=0$ 的两根都大于 $1$，}
\step{所以 $\begin{cases}\Delta=(-4)^2-4a\geqslant0\\ f(1)>0\end{cases}$，即 $\begin{cases}16-4a\geqslant0\\ 1-4+a>0\end{cases}$，解得 $3<a\leqslant4$，}
\step{即 $a$ 的取值范围为 $(3,4]$.}
\solend
\fi

\item 若不等式 $(a-2)x^2+2(a-2)x-4<0$ 对 $x\in R$ 恒成立，则实数 $a$ 的取值范围为\blankline[4.4em].

\ifanswer
\solbegin
\step{①当 $a-2=0$，即 $a=2$ 时，不等式为 $-4<0$ 恒成立，符合题意；}
\step{②当 $a\neq2$ 时，则 $\begin{cases}a-2<0\\ \Delta=4(a-2)^2+16(a-2)<0\end{cases}$，解得 $-2<a<2$.}
\step{综上所述，实数 $a$ 的取值范围是 $(-2,2]$.}
\solend
\fi

\item 已知 $-1\leqslant a+b\leqslant4,2\leqslant a-b\leqslant3$，则 $3a-2b$ 的取值范围为\blankline[4.4em].

\ifanswer
\solbegin
\step{设 $3a-2b=m(a+b)+n(a-b)$，则 $m+n=3,m-n=-2$，}
\step{所以 $m=\dfrac12,n=\dfrac52$，所以 $3a-2b=\dfrac12(a+b)+\dfrac52(a-b)\in\left[\dfrac92,\dfrac{19}2\right]$.}
\solend
\fi

\item 设 $x_1$、$x_2$ 是关于 $x$ 的方程 $2x^2-4mx+2m^2+3m-2=0$ 的两个实数根，则 $\lvert x_1-x_2\rvert$ 的最小值为\blankline[4.4em].

\ifanswer
\solbegin
\step{$x_1$、$x_2$ 是方程 $2x^2-4mx+2m^2+3m-2=0$ 的两个实数根，}
\step{所以 $\Delta=(4m)^2-8(2m^2+3m-2)\geqslant0$，解得 $m\leqslant\dfrac23$，}
\step{由韦达定理得 $x_1+x_2=2m,x_1\cdot x_2=\dfrac{2m^2+3m-2}2$，}
\step{则 $\lvert x_1-x_2\rvert=\sqrt{(x_1+x_2)^2-4x_1\cdot x_2}=\sqrt{4-6m}$，}
\step{当 $m=\dfrac23$ 时，$\lvert x_1-x_2\rvert$ 的最小值为 $0$.}
\solend
\fi

\item 设 $x\in R$，$[x]$ 表示不大于 $x$ 的最大整数，如 $[\pi]=3$，$[-1.2]=-2$，$\left[\dfrac12\right]=0$，则使 $\lvert\lvert x^2-1\rvert\rvert=3$ 成立的 $x$ 的取值范围是\blankline[4.4em].

\ifanswer
\solbegin
\step{若 $\lvert\lvert x^2-1\rvert\rvert=3$，则 $3\leqslant\lvert x^2-1\rvert<4$，即 $4\leqslant x^2<5$，}
\step{解得 $2\leqslant x<\sqrt5$ 或 $-\sqrt5<x\leqslant-2$，故 $x$ 的取值范围是 $(-\sqrt5,-2]\cup[2,\sqrt5)$.}
\solend
\fi

\item 设 $a$ 为实数，若关于 $x$ 的一元一次不等式组 $\begin{cases}2x+a>0\\ 3x-6a<0\end{cases}$ 的解集中有且仅有 $4$ 个整数，则 $a$ 的取值范围是\blankline[4.4em].

\ifanswer
\solbegin
\step{一元一次不等式组 $\begin{cases}2x+a>0\\ 3x-6a<0\end{cases}\Leftrightarrow-\dfrac a2<x<2a$，且 $a>0$，}
\step{关于 $x$ 的一元一次不等式组 $\begin{cases}2x+a>0\\ 3x-6a<0\end{cases}$ 的解集中有且仅有 $4$ 个整数，}
\step{①当 $4$ 个整数为 $0$，$1$，$2$，$3$ 时，则 $\begin{cases}-1\leqslant-\dfrac a2<0\\ 3<2a\leqslant4\end{cases}$，所以 $\dfrac32<a\leqslant2$，}
\step{②当 $4$ 个整数为 $-1$，$0$，$1$，$2$ 时，则 $\begin{cases}-2\leqslant-\dfrac a2<-1\\ 2<2a\leqslant3\end{cases}$，无解，}
\step{综上，$a$ 的取值范围是 $\left(\dfrac32,2\right]$.}
\solend
\fi

\item 集合 $M$、$N$、$S$ 都是非空集合，现规定如下运算：

$M\odot N\odot S=\{x\mid x\in(M\cap N)\cup(N\cap S)\cup(S\cap M)$ 且 $x\notin M\cap N\cap S\}$．假设集合 $A=\{x\mid a<x<b\}$，$B=\{x\mid c<x<d\}$，$C=\{x\mid e<x<f\}$，其中实数 $a$、$b$、$c$、$d$、$e$、$f$ 满足：$a<c<e<0<b<d<f$．计算 $A\odot B\odot C=\blankline[4.4em]$.

\ifanswer
\solbegin
\step{因为集合 $A=\{x\mid a<x<b\}$，$B=\{x\mid c<x<d\}$，$C=\{x\mid e<x<f\}$，}
\step{又因为 $a<c<e<0<b<d<f$，}
\step{所以 $A\cap B=\{x\mid c<x<b\}$，$B\cap C=\{x\mid e<x<d\}$，$C\cap A=\{x\mid e<x<b\}$，}
\step{$A\cap B\cap C=\{x\mid e<x<b\}$，故 $A\odot B\odot C=\{x\mid c<x\leqslant e$ 或 $b\leqslant x<d\}$.}
\solend
\fi

\end{enumerate}

\bigsec{二、选择题}

\begin{enumerate}
\setcounter{enumi}{12}

\item “$x\neq1$”是“$x^2-1\neq0$”的\mbox{（\quad）}条件.

\keepwithprev
A. 充分非必要\qquad B. 必要非充分

C. 充要\qquad D. 既非充分也非必要

\ifanswer
\solbegin
\step{$x^2-1\neq0\Leftrightarrow x\neq1$ 且 $x\neq-1\Rightarrow x\neq1$.}
\step{反之不成立，如 $x=-1$ 时，$x\neq1$，但 $x^2-1=0$.}
\step{因此“$x\neq1$”是“$x^2-1\neq0$”的必要非充分条件，故选 $B$.}
\solend
\fi

\item 若 $\dfrac1a<\dfrac1b<0$，则下列结论中不正确的是\mbox{（\quad）}

\keepwithprev
A. $a+b<0$\qquad B. $a^2<b^2$\qquad C. $ab<b^2$\qquad D. $ab<a^2$

\ifanswer
\solbegin
\step{由 $\dfrac1a<\dfrac1b<0$ 得 $b<a<0$，因此 $a+b<0$，且 $a^2<ab<b^2$，}
\step{所以 A、B、C 正确，D 不正确．故选 $D$.}
\solend
\fi

\item 非空集合 $X=\{x\mid a+1\leqslant x\leqslant3a-5\}$，$Y=\{x\mid1\leqslant x\leqslant16\}$，使得 $X\subseteq(X\cap Y)$ 成立的所有 $a$ 的集合是\mbox{（\quad）}

\keepwithprev
A. $\{a\mid3\leqslant a\leqslant7\}$\qquad B. $\{a\mid0\leqslant a\leqslant7\}$\qquad C. $\{a\mid3<a\leqslant7\}$\qquad D. $\{a\mid a\leqslant7\}$

\ifanswer
\solbegin
\step{因为 $X\subseteq(X\cap Y)$ 成立，所以 $X\subseteq Y$ 成立.}
\step{$X\neq\varnothing$，要使 $X\subseteq Y$，则 $\begin{cases}a\geqslant3\\ a+1\geqslant1\\ 3a-5\leqslant16\end{cases}$，即 $\begin{cases}a\geqslant3\\ a\geqslant0\\ a\leqslant7\end{cases}$，所以 $3\leqslant a\leqslant7$，故选 $A$.}
\solend
\fi

\item 已知 $M$ 是满足下列条件的集合：①$0\in M$，$1\in M$；②若 $x$、$y\in M$，则 $x-y\in M$；③若 $x\in M$ 且 $x\neq0$，则 $\dfrac1x\in M$．则下列说法正确的个数为\mbox{（\quad）}

（1）$\dfrac12\in M$；（2）$x+y\in M$；（3）$x^2\in M$．

\keepwithprev
A. $0$\qquad B. $1$\qquad C. $2$\qquad D. $3$

\ifanswer
\solbegin
\step{（1）$\dfrac12\in M$ 正确，理由如下：}
\step{由①得 $0\in M$，$1\in M$，由②得 $0-1=-1\in M$，$1-(-1)=2\in M$，}
\step{由③得 $\dfrac12\in M$；}
\step{（2）由①得 $0\in M$，由题意 $y\in M$，}
\step{所以由②得 $0-y=-y\in M$，又 $x\in M$，所以 $x-(-y)=x+y\in M$；}
\step{（3）$x\in M$，由②得 $x-1\in M$，由③得 $\dfrac1x\in M$，$\dfrac1{x-1}\in M$，}
\step{所以 $\dfrac1x-\dfrac1{x-1}\in M$，即 $\dfrac1{x(x-1)}\in M$，所以 $x(x-1)\in M$，}
\step{由（2）结论得 $x(x-1)+x\in M$，即 $x^2\in M$，}
\step{所以说法正确的个数为 $3$．故选 $D$.}
\solend
\fi

\end{enumerate}

\bigsec{三、解答题}

\begin{enumerate}
\setcounter{enumi}{16}

\item 已知 $a>b$，用比较法证明：$a^3>b^3$.

\ansspace*[8]

\ifanswer
\solbegin
\step{$a>b$，$a^3-b^3=(a-b)(a^2+ab+b^2)=(a-b)\left[\left(a+\dfrac b2\right)^2+\dfrac34b^2\right]>0$，}
\step{故 $a^3>b^3$.}
\solend
\fi

\item 求下列关于 $x$ 的不等式的解集：

（1）$\dfrac{x^2(x+3)(x^2-x+1)}{x^2-5x+4}>0$；

（2）$\dfrac{ax-1}{x-2}\geqslant0\left(a<\dfrac12\right)$.

\ansspace*[12]

\ifanswer
\solbegin
\stepm{（1）}{分母 $x^2-5x+4=(x-1)(x-4)$，要求 $x\neq1,4$；}
\step{又 $x=0$ 时原不等式不成立，且 $x^2-x+1=\left(x-\dfrac12\right)^2+\dfrac34>0$，}
\step{因此原不等式等价于 $\dfrac{x+3}{(x-1)(x-4)}>0,x\neq0$.}
\step{在数轴上标出 $-3,1,4$，分段判号得 $-3<x<1$ 或 $x>4$，}
\step{再排除 $x=0$，解集为 $(-3,0)\cup(0,1)\cup(4,+\infty)$.}
\stepm{（2）}{原不等式要求 $x\neq2$，按 $a$ 的符号分类讨论.}
\step{①当 $a<0$ 时，$\dfrac1a<2$，原不等式等价于 $\dfrac{x-\dfrac1a}{x-2}\leqslant0\Leftrightarrow\dfrac1a\leqslant x<2$，}
\step{解集为 $\left[\dfrac1a,2\right)$.}
\step{②当 $a=0$ 时，原不等式为 $\dfrac{-1}{x-2}\geqslant0\Leftrightarrow x<2$，解集为 $(-\infty,2)$.}
\step{③当 $0<a<\dfrac12$ 时，$\dfrac1a>2$，}
\step{原不等式等价于 $\dfrac{x-\dfrac1a}{x-2}\geqslant0\Leftrightarrow x<2$ 或 $x\geqslant\dfrac1a$，解集为 $(-\infty,2)\cup\left[\dfrac1a,+\infty\right)$.}
\solend
\fi

\item 集合 $A=\{x\mid-2\leqslant x\leqslant5\}$，$B=\{x\mid m+1\leqslant x\leqslant2m-1\}$.

（1）若 $B\subseteq A$，求实数 $m$ 的取值范围；

（2）若 $A\cap B\neq\varnothing$，求实数 $m$ 的取值范围.

\ansspace*[10]

\ifanswer
\solbegin
\stepm{（1）}{当 $B=\varnothing$ 时，$m+1>2m-1$，解得 $m<2$，}
\step{当 $B\neq\varnothing$ 时，$\begin{cases}m+1\leqslant2m-1\\ m+1\geqslant-2\\ 2m-1\leqslant5\end{cases}$，解得 $2\leqslant m\leqslant3$，}
\step{综上，实数 $m$ 的取值范围是 $(-\infty,3]$；}
\stepm{（2）}{当 $B=\varnothing$ 时，$m+1>2m-1$，解得 $m<2$，}
\step{当 $B\neq\varnothing$ 时，$\begin{cases}m+1\leqslant2m-1\\ 2m-1<-2\end{cases}$ 或 $\begin{cases}m+1\leqslant2m-1\\ m+1>5\end{cases}$，解得 $m\in\varnothing$ 或 $m>4$.}
\step{所以当 $A\cap B\neq\varnothing$ 时，实数 $m$ 的取值范围是 $[2,4]$.}
\solend
\fi

\item 在解决实际问题时，往往会有不同的思路和方法，这些方法有些正确，有些错误；有些简洁，有些复杂.

问题①：设 $a\in R$，集合 $M=\{x\mid(x-1)(x-a)=0\}$，若 $x\in M$ 是 $x\in\{1,2,3\}$ 的充分条件，求：$a$ 的取值集合.

问题②：设 $x\in Z$，$a=x+3$，$b=18-x$，若 $c\in Z$，求证：$c-a$ 和 $c-b$ 至少有一个数是奇数

（1）小明在解决问题①，他认为原问题等价于 $\{1,a\}\subseteq\{1,2,3\}$，解得 $a$ 的取值集合为 $\{2,3\}$，张老师判断小明解题错误，请计算正确的 $a$ 的取值集合；

（2）小红认为既然 $c\in Z$，只需根据 $c$ 是奇数还是偶数，分类讨论即可；小华则认为可以使用反证法解决问题，请你选一种你认为更好的方法并证明.

\ansspace*[12]

\ifanswer
\solbegin
\stepm{（1）}{当 $a=1$ 时，集合 $M=\{x\mid(x-1)^2=0\}=\{1\}$，}
\step{此时 $x\in M\Rightarrow x\in\{1,2,3\}$，即 $x\in M$ 是 $x\in\{1,2,3\}$ 的充分条件.}
\step{当 $a\neq1$ 时，集合 $M=\{x\mid(x-1)(x-a)=0\}=\{1,a\}$，}
\step{由 $x\in M$ 是 $x\in\{1,2,3\}$ 的充分条件，得 $\{1,a\}\subseteq\{1,2,3\}$，}
\step{所以 $a=2$ 或 $a=3$.}
\step{综上，实数 $a$ 的取值集合为 $\{1,2,3\}$；}
\stepm{（2）}{利用反证法证明．证明如下：}
\step{假设 $c-a$、$c-b$ 均为偶数，则 $(c-b)-(c-a)=a-b$ 为偶数.}
\step{因为 $a-b=x+3-18+x=2x-15$ 为奇数，}
\step{与 $b-a$ 为偶数相矛盾，所以原假设不成立.}
\step{故 $c-a$、$c-b$ 至少有一个数是奇数.}
\solend
\fi

\item 已知实数 $x_1$、$x_2$、$x_3$、$x_4$、$x_5$，满足 $x_1+x_2+x_3+x_4+x_5=5$.

（1）求证：$x_1$、$x_2$、$x_2$、$x_4$、$x_5$ 中至少有一个实数不小于 $1$；% 原卷把 x_3 误印作 x_2，照录

（2）设 $x_1$、$x_2$、$x_3$、$x_4$、$x_5$ 这五个实数两两不等，集合 $A=\{x_1,x_2,x_3,x_4,x_5\}$，若 $B\neq\varnothing$ 且 $B\subseteq A$，记 $G(B)$ 是 $B$ 中所有元素之和，对所有的 $B$，求 $G(B)$ 的平均值.

\ansspace*[13]

\ifanswer
\solbegin
\stepm{（1）}{假设 $x_1$、$x_2$、$x_2$、$x_4$、$x_5$ 中的每一个数都小于 $1$，}% 原卷 x_3 误印作 x_2，照录
\step{得 $x_1+x_2+x_2+x_4+x_5<5$，这与 $x_1+x_2+x_3+x_4+x_5=5$ 矛盾，}% 前半句原卷 x_3 误印作 x_2，照录
\step{故 $x_1$、$x_2$、$x_2$、$x_4$、$x_5$ 中至少有一个实数不小于 $1$；}% 原卷 x_3 误印作 x_2，照录
\stepm{（2）}{集合 $A=\{x_1,x_2,x_3,x_4,x_5\}$ 的非空子集个数为 $32-1=31$，}
\step{由于 $G(B)$ 是 $B$ 中所有元素之和，}
\step{得 $G(B)=16(x_1+x_2+x_3+x_4+x_5)=16\times5=80$，}
\step{则 $G(B)$ 的平均值为 $\dfrac{80}{31}$.}
\solend
\fi

\end{enumerate}

\end{document}
"
% 紧凑版：xelatex "\def\iscompact{1}% 高一98_市西中学_10月月考.tex
%   —— 高一上数学“2026-2027 年市西中学高一上 10 月月考”（讲评版）
% 试卷版：xelatex 高一98_市西中学_10月月考.tex
% 答案版：xelatex "\def\isanswer{1}% 高一98_市西中学_10月月考.tex
%   —— 高一上数学“2026-2027 年市西中学高一上 10 月月考”（讲评版）
% 试卷版：xelatex 高一98_市西中学_10月月考.tex
% 答案版：xelatex "\def\isanswer{1}% 高一98_市西中学_10月月考.tex
%   —— 高一上数学“2026-2027 年市西中学高一上 10 月月考”（讲评版）
% 试卷版：xelatex 高一98_市西中学_10月月考.tex
% 答案版：xelatex "\def\isanswer{1}\input{高一98_市西中学_10月月考.tex}"
% 紧凑版：xelatex "\def\iscompact{1}\input{高一98_市西中学_10月月考.tex}"
%
% 原始材料：微信公众号“魔都数学加油站”文章，连载“27学年高一98”；
% 原文标题“27学年高一98——2026-2027年上海市市西中学高一上10月月考”；
% 原文链接：https://mp.weixin.qq.com/s/duK4wz-HT3W9tjOeYQaVfg。
% 正文为 8 张 PNG 页（`原图/高一98_市西中学_10月月考/01..08.png`），
% 尺寸混排（第 1—4、6、7 页为 4762×6735，第 5、8 页为 2560×3620），按页序读取。
%
% 篇幅与题号：原卷自第 3 题起编号（第 1、2 题不在原文内）。填空题 3—12（10 题）、
% 选择题 13—16（4 题）、解答题 17—21（5 题），共 19 题。
%
% 讲评版：全部 19 题均印【解析】，其中第 13—16 题的答案在解析末句
% “故选 $B$／$D$／$A$／$D$”（题干括号内未印答案）。全卷无插图。
%
% 副标题：原卷未印副标题。本稿按内容标作“集合与常用逻辑用语、不等式”。
%
% 原卷笔误/存疑（均照录未改，正文对应处另有 % 注）：
%   ① 第 21 题 (1) 求证行、解析“假设”行与“故”行三处把 $x_3$ 误印作 $x_2$
%      （“$x_1$、$x_2$、$x_2$、$x_4$、$x_5$”）；解析同一行后半句
%      “这与 $x_1+x_2+x_3+x_4+x_5=5$ 矛盾”却用 $x_3$，显系笔误，照录；
%   ② 第 6 题解析的判别式行印作“$\Delta=(-4)^2-4a\geqslant0$”，与
%      $f(1)>0$ 并列，照录。
%
% 与原卷的排版差异：
%   ① 原卷每题下直接印【解析】，本稿按全稿惯例将答案与解析置于答案版；
%   ② 原卷未印副标题，本稿补出；填空横线统一排 \blankline[4.4em]；
%   ③ 原卷数集记号排斜体 $R$、$\Z$，本稿按全稿惯例改用 $\R$、$\Z$；
%   ④ 解答题按惯例加 \ansspace；选择题括号排 \mbox{（\quad）}；
%   ⑤ 第 12 题定义的运算 $\odot$ 原卷排实心小圆，本稿排 $\odot$。
%
% 行数核对：逐题数 \step/\stepm 条数，合计 82 条（Q3：1；Q4：2；Q5：1；Q6：4；Q7：3；
% Q8：2；Q9：5；Q10：2；Q11：5；Q12：4；Q13：3；Q14：2；Q15：2；Q16：9；Q17：2；
% Q18：11；Q19：6；Q20：11；Q21：7）。全卷 19 题均印【解析】。
% 原卷解析的换行按语意合并，未逐行强拆。
%
\documentclass[a4paper,11pt]{article}
\usepackage{common}
\begin{document}

\examtitlest[7]{2026--2027 年市西中学高一上 10 月月考}{集合与常用逻辑用语、不等式}

\bigsec{一、填空题}

\begin{enumerate}
\setcounter{enumi}{2}

\item 集合 $A=\{x\mid y=\sqrt{1-x^2}\},B=\{y\mid y=x^2+1,x\in A\}$，则 $A\cap\overline B=\blankline[4.4em]$.

\ifanswer
\solbegin
\step{$1-x^2\geqslant0\Rightarrow A=[-1,1]$，从而 $B=[1,2]$，所以 $A\cap\overline B=[-1,1)$.}
\solend
\fi

\item 已知 $b,c\in R$，关于 $x$ 的不等式 $x^2-bx+c<0$ 的解集为 $(-3,2)$，则 $b+c=\blankline[4.4em]$.

\ifanswer
\solbegin
\step{由题意得 $x=-3$ 与 $x=2$ 是 $x^2-bx+c=0$ 的两根，}
\step{所以 $\begin{cases}-3+2=b\\ -3\times2=c\end{cases}$，解得 $\begin{cases}b=-1\\ c=-6\end{cases}$，所以 $b+c=-7$.}
\solend
\fi

\item 若 $\alpha:\lvert x-m\rvert<1$ 是 $\beta:x^2-4x<0$ 的充分条件，则实数 $m$ 的取值范围是\blankline[4.4em].

\ifanswer
\solbegin
\step{$(m-1,m+1)\subseteq(0,4)\Leftrightarrow m-1\geqslant0$ 且 $m+1\leqslant4\Leftrightarrow m\in[1,3]$.}
\solend
\fi

\item 已知方程 $x^2-4x+a=0$ 的两根都大于 $1$，则 $a$ 的取值范围为\blankline[4.4em].

\ifanswer
\solbegin
\step{设函数 $f(x)=x^2-4x+a$，对称轴为 $x=2$，}
\step{因为方程 $x^2-4x+a=0$ 的两根都大于 $1$，}
\step{所以 $\begin{cases}\Delta=(-4)^2-4a\geqslant0\\ f(1)>0\end{cases}$，即 $\begin{cases}16-4a\geqslant0\\ 1-4+a>0\end{cases}$，解得 $3<a\leqslant4$，}
\step{即 $a$ 的取值范围为 $(3,4]$.}
\solend
\fi

\item 若不等式 $(a-2)x^2+2(a-2)x-4<0$ 对 $x\in R$ 恒成立，则实数 $a$ 的取值范围为\blankline[4.4em].

\ifanswer
\solbegin
\step{①当 $a-2=0$，即 $a=2$ 时，不等式为 $-4<0$ 恒成立，符合题意；}
\step{②当 $a\neq2$ 时，则 $\begin{cases}a-2<0\\ \Delta=4(a-2)^2+16(a-2)<0\end{cases}$，解得 $-2<a<2$.}
\step{综上所述，实数 $a$ 的取值范围是 $(-2,2]$.}
\solend
\fi

\item 已知 $-1\leqslant a+b\leqslant4,2\leqslant a-b\leqslant3$，则 $3a-2b$ 的取值范围为\blankline[4.4em].

\ifanswer
\solbegin
\step{设 $3a-2b=m(a+b)+n(a-b)$，则 $m+n=3,m-n=-2$，}
\step{所以 $m=\dfrac12,n=\dfrac52$，所以 $3a-2b=\dfrac12(a+b)+\dfrac52(a-b)\in\left[\dfrac92,\dfrac{19}2\right]$.}
\solend
\fi

\item 设 $x_1$、$x_2$ 是关于 $x$ 的方程 $2x^2-4mx+2m^2+3m-2=0$ 的两个实数根，则 $\lvert x_1-x_2\rvert$ 的最小值为\blankline[4.4em].

\ifanswer
\solbegin
\step{$x_1$、$x_2$ 是方程 $2x^2-4mx+2m^2+3m-2=0$ 的两个实数根，}
\step{所以 $\Delta=(4m)^2-8(2m^2+3m-2)\geqslant0$，解得 $m\leqslant\dfrac23$，}
\step{由韦达定理得 $x_1+x_2=2m,x_1\cdot x_2=\dfrac{2m^2+3m-2}2$，}
\step{则 $\lvert x_1-x_2\rvert=\sqrt{(x_1+x_2)^2-4x_1\cdot x_2}=\sqrt{4-6m}$，}
\step{当 $m=\dfrac23$ 时，$\lvert x_1-x_2\rvert$ 的最小值为 $0$.}
\solend
\fi

\item 设 $x\in R$，$[x]$ 表示不大于 $x$ 的最大整数，如 $[\pi]=3$，$[-1.2]=-2$，$\left[\dfrac12\right]=0$，则使 $\lvert\lvert x^2-1\rvert\rvert=3$ 成立的 $x$ 的取值范围是\blankline[4.4em].

\ifanswer
\solbegin
\step{若 $\lvert\lvert x^2-1\rvert\rvert=3$，则 $3\leqslant\lvert x^2-1\rvert<4$，即 $4\leqslant x^2<5$，}
\step{解得 $2\leqslant x<\sqrt5$ 或 $-\sqrt5<x\leqslant-2$，故 $x$ 的取值范围是 $(-\sqrt5,-2]\cup[2,\sqrt5)$.}
\solend
\fi

\item 设 $a$ 为实数，若关于 $x$ 的一元一次不等式组 $\begin{cases}2x+a>0\\ 3x-6a<0\end{cases}$ 的解集中有且仅有 $4$ 个整数，则 $a$ 的取值范围是\blankline[4.4em].

\ifanswer
\solbegin
\step{一元一次不等式组 $\begin{cases}2x+a>0\\ 3x-6a<0\end{cases}\Leftrightarrow-\dfrac a2<x<2a$，且 $a>0$，}
\step{关于 $x$ 的一元一次不等式组 $\begin{cases}2x+a>0\\ 3x-6a<0\end{cases}$ 的解集中有且仅有 $4$ 个整数，}
\step{①当 $4$ 个整数为 $0$，$1$，$2$，$3$ 时，则 $\begin{cases}-1\leqslant-\dfrac a2<0\\ 3<2a\leqslant4\end{cases}$，所以 $\dfrac32<a\leqslant2$，}
\step{②当 $4$ 个整数为 $-1$，$0$，$1$，$2$ 时，则 $\begin{cases}-2\leqslant-\dfrac a2<-1\\ 2<2a\leqslant3\end{cases}$，无解，}
\step{综上，$a$ 的取值范围是 $\left(\dfrac32,2\right]$.}
\solend
\fi

\item 集合 $M$、$N$、$S$ 都是非空集合，现规定如下运算：

$M\odot N\odot S=\{x\mid x\in(M\cap N)\cup(N\cap S)\cup(S\cap M)$ 且 $x\notin M\cap N\cap S\}$．假设集合 $A=\{x\mid a<x<b\}$，$B=\{x\mid c<x<d\}$，$C=\{x\mid e<x<f\}$，其中实数 $a$、$b$、$c$、$d$、$e$、$f$ 满足：$a<c<e<0<b<d<f$．计算 $A\odot B\odot C=\blankline[4.4em]$.

\ifanswer
\solbegin
\step{因为集合 $A=\{x\mid a<x<b\}$，$B=\{x\mid c<x<d\}$，$C=\{x\mid e<x<f\}$，}
\step{又因为 $a<c<e<0<b<d<f$，}
\step{所以 $A\cap B=\{x\mid c<x<b\}$，$B\cap C=\{x\mid e<x<d\}$，$C\cap A=\{x\mid e<x<b\}$，}
\step{$A\cap B\cap C=\{x\mid e<x<b\}$，故 $A\odot B\odot C=\{x\mid c<x\leqslant e$ 或 $b\leqslant x<d\}$.}
\solend
\fi

\end{enumerate}

\bigsec{二、选择题}

\begin{enumerate}
\setcounter{enumi}{12}

\item “$x\neq1$”是“$x^2-1\neq0$”的\mbox{（\quad）}条件.

\keepwithprev
A. 充分非必要\qquad B. 必要非充分

C. 充要\qquad D. 既非充分也非必要

\ifanswer
\solbegin
\step{$x^2-1\neq0\Leftrightarrow x\neq1$ 且 $x\neq-1\Rightarrow x\neq1$.}
\step{反之不成立，如 $x=-1$ 时，$x\neq1$，但 $x^2-1=0$.}
\step{因此“$x\neq1$”是“$x^2-1\neq0$”的必要非充分条件，故选 $B$.}
\solend
\fi

\item 若 $\dfrac1a<\dfrac1b<0$，则下列结论中不正确的是\mbox{（\quad）}

\keepwithprev
A. $a+b<0$\qquad B. $a^2<b^2$\qquad C. $ab<b^2$\qquad D. $ab<a^2$

\ifanswer
\solbegin
\step{由 $\dfrac1a<\dfrac1b<0$ 得 $b<a<0$，因此 $a+b<0$，且 $a^2<ab<b^2$，}
\step{所以 A、B、C 正确，D 不正确．故选 $D$.}
\solend
\fi

\item 非空集合 $X=\{x\mid a+1\leqslant x\leqslant3a-5\}$，$Y=\{x\mid1\leqslant x\leqslant16\}$，使得 $X\subseteq(X\cap Y)$ 成立的所有 $a$ 的集合是\mbox{（\quad）}

\keepwithprev
A. $\{a\mid3\leqslant a\leqslant7\}$\qquad B. $\{a\mid0\leqslant a\leqslant7\}$\qquad C. $\{a\mid3<a\leqslant7\}$\qquad D. $\{a\mid a\leqslant7\}$

\ifanswer
\solbegin
\step{因为 $X\subseteq(X\cap Y)$ 成立，所以 $X\subseteq Y$ 成立.}
\step{$X\neq\varnothing$，要使 $X\subseteq Y$，则 $\begin{cases}a\geqslant3\\ a+1\geqslant1\\ 3a-5\leqslant16\end{cases}$，即 $\begin{cases}a\geqslant3\\ a\geqslant0\\ a\leqslant7\end{cases}$，所以 $3\leqslant a\leqslant7$，故选 $A$.}
\solend
\fi

\item 已知 $M$ 是满足下列条件的集合：①$0\in M$，$1\in M$；②若 $x$、$y\in M$，则 $x-y\in M$；③若 $x\in M$ 且 $x\neq0$，则 $\dfrac1x\in M$．则下列说法正确的个数为\mbox{（\quad）}

（1）$\dfrac12\in M$；（2）$x+y\in M$；（3）$x^2\in M$．

\keepwithprev
A. $0$\qquad B. $1$\qquad C. $2$\qquad D. $3$

\ifanswer
\solbegin
\step{（1）$\dfrac12\in M$ 正确，理由如下：}
\step{由①得 $0\in M$，$1\in M$，由②得 $0-1=-1\in M$，$1-(-1)=2\in M$，}
\step{由③得 $\dfrac12\in M$；}
\step{（2）由①得 $0\in M$，由题意 $y\in M$，}
\step{所以由②得 $0-y=-y\in M$，又 $x\in M$，所以 $x-(-y)=x+y\in M$；}
\step{（3）$x\in M$，由②得 $x-1\in M$，由③得 $\dfrac1x\in M$，$\dfrac1{x-1}\in M$，}
\step{所以 $\dfrac1x-\dfrac1{x-1}\in M$，即 $\dfrac1{x(x-1)}\in M$，所以 $x(x-1)\in M$，}
\step{由（2）结论得 $x(x-1)+x\in M$，即 $x^2\in M$，}
\step{所以说法正确的个数为 $3$．故选 $D$.}
\solend
\fi

\end{enumerate}

\bigsec{三、解答题}

\begin{enumerate}
\setcounter{enumi}{16}

\item 已知 $a>b$，用比较法证明：$a^3>b^3$.

\ansspace*[8]

\ifanswer
\solbegin
\step{$a>b$，$a^3-b^3=(a-b)(a^2+ab+b^2)=(a-b)\left[\left(a+\dfrac b2\right)^2+\dfrac34b^2\right]>0$，}
\step{故 $a^3>b^3$.}
\solend
\fi

\item 求下列关于 $x$ 的不等式的解集：

（1）$\dfrac{x^2(x+3)(x^2-x+1)}{x^2-5x+4}>0$；

（2）$\dfrac{ax-1}{x-2}\geqslant0\left(a<\dfrac12\right)$.

\ansspace*[12]

\ifanswer
\solbegin
\stepm{（1）}{分母 $x^2-5x+4=(x-1)(x-4)$，要求 $x\neq1,4$；}
\step{又 $x=0$ 时原不等式不成立，且 $x^2-x+1=\left(x-\dfrac12\right)^2+\dfrac34>0$，}
\step{因此原不等式等价于 $\dfrac{x+3}{(x-1)(x-4)}>0,x\neq0$.}
\step{在数轴上标出 $-3,1,4$，分段判号得 $-3<x<1$ 或 $x>4$，}
\step{再排除 $x=0$，解集为 $(-3,0)\cup(0,1)\cup(4,+\infty)$.}
\stepm{（2）}{原不等式要求 $x\neq2$，按 $a$ 的符号分类讨论.}
\step{①当 $a<0$ 时，$\dfrac1a<2$，原不等式等价于 $\dfrac{x-\dfrac1a}{x-2}\leqslant0\Leftrightarrow\dfrac1a\leqslant x<2$，}
\step{解集为 $\left[\dfrac1a,2\right)$.}
\step{②当 $a=0$ 时，原不等式为 $\dfrac{-1}{x-2}\geqslant0\Leftrightarrow x<2$，解集为 $(-\infty,2)$.}
\step{③当 $0<a<\dfrac12$ 时，$\dfrac1a>2$，}
\step{原不等式等价于 $\dfrac{x-\dfrac1a}{x-2}\geqslant0\Leftrightarrow x<2$ 或 $x\geqslant\dfrac1a$，解集为 $(-\infty,2)\cup\left[\dfrac1a,+\infty\right)$.}
\solend
\fi

\item 集合 $A=\{x\mid-2\leqslant x\leqslant5\}$，$B=\{x\mid m+1\leqslant x\leqslant2m-1\}$.

（1）若 $B\subseteq A$，求实数 $m$ 的取值范围；

（2）若 $A\cap B\neq\varnothing$，求实数 $m$ 的取值范围.

\ansspace*[10]

\ifanswer
\solbegin
\stepm{（1）}{当 $B=\varnothing$ 时，$m+1>2m-1$，解得 $m<2$，}
\step{当 $B\neq\varnothing$ 时，$\begin{cases}m+1\leqslant2m-1\\ m+1\geqslant-2\\ 2m-1\leqslant5\end{cases}$，解得 $2\leqslant m\leqslant3$，}
\step{综上，实数 $m$ 的取值范围是 $(-\infty,3]$；}
\stepm{（2）}{当 $B=\varnothing$ 时，$m+1>2m-1$，解得 $m<2$，}
\step{当 $B\neq\varnothing$ 时，$\begin{cases}m+1\leqslant2m-1\\ 2m-1<-2\end{cases}$ 或 $\begin{cases}m+1\leqslant2m-1\\ m+1>5\end{cases}$，解得 $m\in\varnothing$ 或 $m>4$.}
\step{所以当 $A\cap B\neq\varnothing$ 时，实数 $m$ 的取值范围是 $[2,4]$.}
\solend
\fi

\item 在解决实际问题时，往往会有不同的思路和方法，这些方法有些正确，有些错误；有些简洁，有些复杂.

问题①：设 $a\in R$，集合 $M=\{x\mid(x-1)(x-a)=0\}$，若 $x\in M$ 是 $x\in\{1,2,3\}$ 的充分条件，求：$a$ 的取值集合.

问题②：设 $x\in Z$，$a=x+3$，$b=18-x$，若 $c\in Z$，求证：$c-a$ 和 $c-b$ 至少有一个数是奇数

（1）小明在解决问题①，他认为原问题等价于 $\{1,a\}\subseteq\{1,2,3\}$，解得 $a$ 的取值集合为 $\{2,3\}$，张老师判断小明解题错误，请计算正确的 $a$ 的取值集合；

（2）小红认为既然 $c\in Z$，只需根据 $c$ 是奇数还是偶数，分类讨论即可；小华则认为可以使用反证法解决问题，请你选一种你认为更好的方法并证明.

\ansspace*[12]

\ifanswer
\solbegin
\stepm{（1）}{当 $a=1$ 时，集合 $M=\{x\mid(x-1)^2=0\}=\{1\}$，}
\step{此时 $x\in M\Rightarrow x\in\{1,2,3\}$，即 $x\in M$ 是 $x\in\{1,2,3\}$ 的充分条件.}
\step{当 $a\neq1$ 时，集合 $M=\{x\mid(x-1)(x-a)=0\}=\{1,a\}$，}
\step{由 $x\in M$ 是 $x\in\{1,2,3\}$ 的充分条件，得 $\{1,a\}\subseteq\{1,2,3\}$，}
\step{所以 $a=2$ 或 $a=3$.}
\step{综上，实数 $a$ 的取值集合为 $\{1,2,3\}$；}
\stepm{（2）}{利用反证法证明．证明如下：}
\step{假设 $c-a$、$c-b$ 均为偶数，则 $(c-b)-(c-a)=a-b$ 为偶数.}
\step{因为 $a-b=x+3-18+x=2x-15$ 为奇数，}
\step{与 $b-a$ 为偶数相矛盾，所以原假设不成立.}
\step{故 $c-a$、$c-b$ 至少有一个数是奇数.}
\solend
\fi

\item 已知实数 $x_1$、$x_2$、$x_3$、$x_4$、$x_5$，满足 $x_1+x_2+x_3+x_4+x_5=5$.

（1）求证：$x_1$、$x_2$、$x_2$、$x_4$、$x_5$ 中至少有一个实数不小于 $1$；% 原卷把 x_3 误印作 x_2，照录

（2）设 $x_1$、$x_2$、$x_3$、$x_4$、$x_5$ 这五个实数两两不等，集合 $A=\{x_1,x_2,x_3,x_4,x_5\}$，若 $B\neq\varnothing$ 且 $B\subseteq A$，记 $G(B)$ 是 $B$ 中所有元素之和，对所有的 $B$，求 $G(B)$ 的平均值.

\ansspace*[13]

\ifanswer
\solbegin
\stepm{（1）}{假设 $x_1$、$x_2$、$x_2$、$x_4$、$x_5$ 中的每一个数都小于 $1$，}% 原卷 x_3 误印作 x_2，照录
\step{得 $x_1+x_2+x_2+x_4+x_5<5$，这与 $x_1+x_2+x_3+x_4+x_5=5$ 矛盾，}% 前半句原卷 x_3 误印作 x_2，照录
\step{故 $x_1$、$x_2$、$x_2$、$x_4$、$x_5$ 中至少有一个实数不小于 $1$；}% 原卷 x_3 误印作 x_2，照录
\stepm{（2）}{集合 $A=\{x_1,x_2,x_3,x_4,x_5\}$ 的非空子集个数为 $32-1=31$，}
\step{由于 $G(B)$ 是 $B$ 中所有元素之和，}
\step{得 $G(B)=16(x_1+x_2+x_3+x_4+x_5)=16\times5=80$，}
\step{则 $G(B)$ 的平均值为 $\dfrac{80}{31}$.}
\solend
\fi

\end{enumerate}

\end{document}
"
% 紧凑版：xelatex "\def\iscompact{1}% 高一98_市西中学_10月月考.tex
%   —— 高一上数学“2026-2027 年市西中学高一上 10 月月考”（讲评版）
% 试卷版：xelatex 高一98_市西中学_10月月考.tex
% 答案版：xelatex "\def\isanswer{1}\input{高一98_市西中学_10月月考.tex}"
% 紧凑版：xelatex "\def\iscompact{1}\input{高一98_市西中学_10月月考.tex}"
%
% 原始材料：微信公众号“魔都数学加油站”文章，连载“27学年高一98”；
% 原文标题“27学年高一98——2026-2027年上海市市西中学高一上10月月考”；
% 原文链接：https://mp.weixin.qq.com/s/duK4wz-HT3W9tjOeYQaVfg。
% 正文为 8 张 PNG 页（`原图/高一98_市西中学_10月月考/01..08.png`），
% 尺寸混排（第 1—4、6、7 页为 4762×6735，第 5、8 页为 2560×3620），按页序读取。
%
% 篇幅与题号：原卷自第 3 题起编号（第 1、2 题不在原文内）。填空题 3—12（10 题）、
% 选择题 13—16（4 题）、解答题 17—21（5 题），共 19 题。
%
% 讲评版：全部 19 题均印【解析】，其中第 13—16 题的答案在解析末句
% “故选 $B$／$D$／$A$／$D$”（题干括号内未印答案）。全卷无插图。
%
% 副标题：原卷未印副标题。本稿按内容标作“集合与常用逻辑用语、不等式”。
%
% 原卷笔误/存疑（均照录未改，正文对应处另有 % 注）：
%   ① 第 21 题 (1) 求证行、解析“假设”行与“故”行三处把 $x_3$ 误印作 $x_2$
%      （“$x_1$、$x_2$、$x_2$、$x_4$、$x_5$”）；解析同一行后半句
%      “这与 $x_1+x_2+x_3+x_4+x_5=5$ 矛盾”却用 $x_3$，显系笔误，照录；
%   ② 第 6 题解析的判别式行印作“$\Delta=(-4)^2-4a\geqslant0$”，与
%      $f(1)>0$ 并列，照录。
%
% 与原卷的排版差异：
%   ① 原卷每题下直接印【解析】，本稿按全稿惯例将答案与解析置于答案版；
%   ② 原卷未印副标题，本稿补出；填空横线统一排 \blankline[4.4em]；
%   ③ 原卷数集记号排斜体 $R$、$\Z$，本稿按全稿惯例改用 $\R$、$\Z$；
%   ④ 解答题按惯例加 \ansspace；选择题括号排 \mbox{（\quad）}；
%   ⑤ 第 12 题定义的运算 $\odot$ 原卷排实心小圆，本稿排 $\odot$。
%
% 行数核对：逐题数 \step/\stepm 条数，合计 82 条（Q3：1；Q4：2；Q5：1；Q6：4；Q7：3；
% Q8：2；Q9：5；Q10：2；Q11：5；Q12：4；Q13：3；Q14：2；Q15：2；Q16：9；Q17：2；
% Q18：11；Q19：6；Q20：11；Q21：7）。全卷 19 题均印【解析】。
% 原卷解析的换行按语意合并，未逐行强拆。
%
\documentclass[a4paper,11pt]{article}
\usepackage{common}
\begin{document}

\examtitlest[7]{2026--2027 年市西中学高一上 10 月月考}{集合与常用逻辑用语、不等式}

\bigsec{一、填空题}

\begin{enumerate}
\setcounter{enumi}{2}

\item 集合 $A=\{x\mid y=\sqrt{1-x^2}\},B=\{y\mid y=x^2+1,x\in A\}$，则 $A\cap\overline B=\blankline[4.4em]$.

\ifanswer
\solbegin
\step{$1-x^2\geqslant0\Rightarrow A=[-1,1]$，从而 $B=[1,2]$，所以 $A\cap\overline B=[-1,1)$.}
\solend
\fi

\item 已知 $b,c\in R$，关于 $x$ 的不等式 $x^2-bx+c<0$ 的解集为 $(-3,2)$，则 $b+c=\blankline[4.4em]$.

\ifanswer
\solbegin
\step{由题意得 $x=-3$ 与 $x=2$ 是 $x^2-bx+c=0$ 的两根，}
\step{所以 $\begin{cases}-3+2=b\\ -3\times2=c\end{cases}$，解得 $\begin{cases}b=-1\\ c=-6\end{cases}$，所以 $b+c=-7$.}
\solend
\fi

\item 若 $\alpha:\lvert x-m\rvert<1$ 是 $\beta:x^2-4x<0$ 的充分条件，则实数 $m$ 的取值范围是\blankline[4.4em].

\ifanswer
\solbegin
\step{$(m-1,m+1)\subseteq(0,4)\Leftrightarrow m-1\geqslant0$ 且 $m+1\leqslant4\Leftrightarrow m\in[1,3]$.}
\solend
\fi

\item 已知方程 $x^2-4x+a=0$ 的两根都大于 $1$，则 $a$ 的取值范围为\blankline[4.4em].

\ifanswer
\solbegin
\step{设函数 $f(x)=x^2-4x+a$，对称轴为 $x=2$，}
\step{因为方程 $x^2-4x+a=0$ 的两根都大于 $1$，}
\step{所以 $\begin{cases}\Delta=(-4)^2-4a\geqslant0\\ f(1)>0\end{cases}$，即 $\begin{cases}16-4a\geqslant0\\ 1-4+a>0\end{cases}$，解得 $3<a\leqslant4$，}
\step{即 $a$ 的取值范围为 $(3,4]$.}
\solend
\fi

\item 若不等式 $(a-2)x^2+2(a-2)x-4<0$ 对 $x\in R$ 恒成立，则实数 $a$ 的取值范围为\blankline[4.4em].

\ifanswer
\solbegin
\step{①当 $a-2=0$，即 $a=2$ 时，不等式为 $-4<0$ 恒成立，符合题意；}
\step{②当 $a\neq2$ 时，则 $\begin{cases}a-2<0\\ \Delta=4(a-2)^2+16(a-2)<0\end{cases}$，解得 $-2<a<2$.}
\step{综上所述，实数 $a$ 的取值范围是 $(-2,2]$.}
\solend
\fi

\item 已知 $-1\leqslant a+b\leqslant4,2\leqslant a-b\leqslant3$，则 $3a-2b$ 的取值范围为\blankline[4.4em].

\ifanswer
\solbegin
\step{设 $3a-2b=m(a+b)+n(a-b)$，则 $m+n=3,m-n=-2$，}
\step{所以 $m=\dfrac12,n=\dfrac52$，所以 $3a-2b=\dfrac12(a+b)+\dfrac52(a-b)\in\left[\dfrac92,\dfrac{19}2\right]$.}
\solend
\fi

\item 设 $x_1$、$x_2$ 是关于 $x$ 的方程 $2x^2-4mx+2m^2+3m-2=0$ 的两个实数根，则 $\lvert x_1-x_2\rvert$ 的最小值为\blankline[4.4em].

\ifanswer
\solbegin
\step{$x_1$、$x_2$ 是方程 $2x^2-4mx+2m^2+3m-2=0$ 的两个实数根，}
\step{所以 $\Delta=(4m)^2-8(2m^2+3m-2)\geqslant0$，解得 $m\leqslant\dfrac23$，}
\step{由韦达定理得 $x_1+x_2=2m,x_1\cdot x_2=\dfrac{2m^2+3m-2}2$，}
\step{则 $\lvert x_1-x_2\rvert=\sqrt{(x_1+x_2)^2-4x_1\cdot x_2}=\sqrt{4-6m}$，}
\step{当 $m=\dfrac23$ 时，$\lvert x_1-x_2\rvert$ 的最小值为 $0$.}
\solend
\fi

\item 设 $x\in R$，$[x]$ 表示不大于 $x$ 的最大整数，如 $[\pi]=3$，$[-1.2]=-2$，$\left[\dfrac12\right]=0$，则使 $\lvert\lvert x^2-1\rvert\rvert=3$ 成立的 $x$ 的取值范围是\blankline[4.4em].

\ifanswer
\solbegin
\step{若 $\lvert\lvert x^2-1\rvert\rvert=3$，则 $3\leqslant\lvert x^2-1\rvert<4$，即 $4\leqslant x^2<5$，}
\step{解得 $2\leqslant x<\sqrt5$ 或 $-\sqrt5<x\leqslant-2$，故 $x$ 的取值范围是 $(-\sqrt5,-2]\cup[2,\sqrt5)$.}
\solend
\fi

\item 设 $a$ 为实数，若关于 $x$ 的一元一次不等式组 $\begin{cases}2x+a>0\\ 3x-6a<0\end{cases}$ 的解集中有且仅有 $4$ 个整数，则 $a$ 的取值范围是\blankline[4.4em].

\ifanswer
\solbegin
\step{一元一次不等式组 $\begin{cases}2x+a>0\\ 3x-6a<0\end{cases}\Leftrightarrow-\dfrac a2<x<2a$，且 $a>0$，}
\step{关于 $x$ 的一元一次不等式组 $\begin{cases}2x+a>0\\ 3x-6a<0\end{cases}$ 的解集中有且仅有 $4$ 个整数，}
\step{①当 $4$ 个整数为 $0$，$1$，$2$，$3$ 时，则 $\begin{cases}-1\leqslant-\dfrac a2<0\\ 3<2a\leqslant4\end{cases}$，所以 $\dfrac32<a\leqslant2$，}
\step{②当 $4$ 个整数为 $-1$，$0$，$1$，$2$ 时，则 $\begin{cases}-2\leqslant-\dfrac a2<-1\\ 2<2a\leqslant3\end{cases}$，无解，}
\step{综上，$a$ 的取值范围是 $\left(\dfrac32,2\right]$.}
\solend
\fi

\item 集合 $M$、$N$、$S$ 都是非空集合，现规定如下运算：

$M\odot N\odot S=\{x\mid x\in(M\cap N)\cup(N\cap S)\cup(S\cap M)$ 且 $x\notin M\cap N\cap S\}$．假设集合 $A=\{x\mid a<x<b\}$，$B=\{x\mid c<x<d\}$，$C=\{x\mid e<x<f\}$，其中实数 $a$、$b$、$c$、$d$、$e$、$f$ 满足：$a<c<e<0<b<d<f$．计算 $A\odot B\odot C=\blankline[4.4em]$.

\ifanswer
\solbegin
\step{因为集合 $A=\{x\mid a<x<b\}$，$B=\{x\mid c<x<d\}$，$C=\{x\mid e<x<f\}$，}
\step{又因为 $a<c<e<0<b<d<f$，}
\step{所以 $A\cap B=\{x\mid c<x<b\}$，$B\cap C=\{x\mid e<x<d\}$，$C\cap A=\{x\mid e<x<b\}$，}
\step{$A\cap B\cap C=\{x\mid e<x<b\}$，故 $A\odot B\odot C=\{x\mid c<x\leqslant e$ 或 $b\leqslant x<d\}$.}
\solend
\fi

\end{enumerate}

\bigsec{二、选择题}

\begin{enumerate}
\setcounter{enumi}{12}

\item “$x\neq1$”是“$x^2-1\neq0$”的\mbox{（\quad）}条件.

\keepwithprev
A. 充分非必要\qquad B. 必要非充分

C. 充要\qquad D. 既非充分也非必要

\ifanswer
\solbegin
\step{$x^2-1\neq0\Leftrightarrow x\neq1$ 且 $x\neq-1\Rightarrow x\neq1$.}
\step{反之不成立，如 $x=-1$ 时，$x\neq1$，但 $x^2-1=0$.}
\step{因此“$x\neq1$”是“$x^2-1\neq0$”的必要非充分条件，故选 $B$.}
\solend
\fi

\item 若 $\dfrac1a<\dfrac1b<0$，则下列结论中不正确的是\mbox{（\quad）}

\keepwithprev
A. $a+b<0$\qquad B. $a^2<b^2$\qquad C. $ab<b^2$\qquad D. $ab<a^2$

\ifanswer
\solbegin
\step{由 $\dfrac1a<\dfrac1b<0$ 得 $b<a<0$，因此 $a+b<0$，且 $a^2<ab<b^2$，}
\step{所以 A、B、C 正确，D 不正确．故选 $D$.}
\solend
\fi

\item 非空集合 $X=\{x\mid a+1\leqslant x\leqslant3a-5\}$，$Y=\{x\mid1\leqslant x\leqslant16\}$，使得 $X\subseteq(X\cap Y)$ 成立的所有 $a$ 的集合是\mbox{（\quad）}

\keepwithprev
A. $\{a\mid3\leqslant a\leqslant7\}$\qquad B. $\{a\mid0\leqslant a\leqslant7\}$\qquad C. $\{a\mid3<a\leqslant7\}$\qquad D. $\{a\mid a\leqslant7\}$

\ifanswer
\solbegin
\step{因为 $X\subseteq(X\cap Y)$ 成立，所以 $X\subseteq Y$ 成立.}
\step{$X\neq\varnothing$，要使 $X\subseteq Y$，则 $\begin{cases}a\geqslant3\\ a+1\geqslant1\\ 3a-5\leqslant16\end{cases}$，即 $\begin{cases}a\geqslant3\\ a\geqslant0\\ a\leqslant7\end{cases}$，所以 $3\leqslant a\leqslant7$，故选 $A$.}
\solend
\fi

\item 已知 $M$ 是满足下列条件的集合：①$0\in M$，$1\in M$；②若 $x$、$y\in M$，则 $x-y\in M$；③若 $x\in M$ 且 $x\neq0$，则 $\dfrac1x\in M$．则下列说法正确的个数为\mbox{（\quad）}

（1）$\dfrac12\in M$；（2）$x+y\in M$；（3）$x^2\in M$．

\keepwithprev
A. $0$\qquad B. $1$\qquad C. $2$\qquad D. $3$

\ifanswer
\solbegin
\step{（1）$\dfrac12\in M$ 正确，理由如下：}
\step{由①得 $0\in M$，$1\in M$，由②得 $0-1=-1\in M$，$1-(-1)=2\in M$，}
\step{由③得 $\dfrac12\in M$；}
\step{（2）由①得 $0\in M$，由题意 $y\in M$，}
\step{所以由②得 $0-y=-y\in M$，又 $x\in M$，所以 $x-(-y)=x+y\in M$；}
\step{（3）$x\in M$，由②得 $x-1\in M$，由③得 $\dfrac1x\in M$，$\dfrac1{x-1}\in M$，}
\step{所以 $\dfrac1x-\dfrac1{x-1}\in M$，即 $\dfrac1{x(x-1)}\in M$，所以 $x(x-1)\in M$，}
\step{由（2）结论得 $x(x-1)+x\in M$，即 $x^2\in M$，}
\step{所以说法正确的个数为 $3$．故选 $D$.}
\solend
\fi

\end{enumerate}

\bigsec{三、解答题}

\begin{enumerate}
\setcounter{enumi}{16}

\item 已知 $a>b$，用比较法证明：$a^3>b^3$.

\ansspace*[8]

\ifanswer
\solbegin
\step{$a>b$，$a^3-b^3=(a-b)(a^2+ab+b^2)=(a-b)\left[\left(a+\dfrac b2\right)^2+\dfrac34b^2\right]>0$，}
\step{故 $a^3>b^3$.}
\solend
\fi

\item 求下列关于 $x$ 的不等式的解集：

（1）$\dfrac{x^2(x+3)(x^2-x+1)}{x^2-5x+4}>0$；

（2）$\dfrac{ax-1}{x-2}\geqslant0\left(a<\dfrac12\right)$.

\ansspace*[12]

\ifanswer
\solbegin
\stepm{（1）}{分母 $x^2-5x+4=(x-1)(x-4)$，要求 $x\neq1,4$；}
\step{又 $x=0$ 时原不等式不成立，且 $x^2-x+1=\left(x-\dfrac12\right)^2+\dfrac34>0$，}
\step{因此原不等式等价于 $\dfrac{x+3}{(x-1)(x-4)}>0,x\neq0$.}
\step{在数轴上标出 $-3,1,4$，分段判号得 $-3<x<1$ 或 $x>4$，}
\step{再排除 $x=0$，解集为 $(-3,0)\cup(0,1)\cup(4,+\infty)$.}
\stepm{（2）}{原不等式要求 $x\neq2$，按 $a$ 的符号分类讨论.}
\step{①当 $a<0$ 时，$\dfrac1a<2$，原不等式等价于 $\dfrac{x-\dfrac1a}{x-2}\leqslant0\Leftrightarrow\dfrac1a\leqslant x<2$，}
\step{解集为 $\left[\dfrac1a,2\right)$.}
\step{②当 $a=0$ 时，原不等式为 $\dfrac{-1}{x-2}\geqslant0\Leftrightarrow x<2$，解集为 $(-\infty,2)$.}
\step{③当 $0<a<\dfrac12$ 时，$\dfrac1a>2$，}
\step{原不等式等价于 $\dfrac{x-\dfrac1a}{x-2}\geqslant0\Leftrightarrow x<2$ 或 $x\geqslant\dfrac1a$，解集为 $(-\infty,2)\cup\left[\dfrac1a,+\infty\right)$.}
\solend
\fi

\item 集合 $A=\{x\mid-2\leqslant x\leqslant5\}$，$B=\{x\mid m+1\leqslant x\leqslant2m-1\}$.

（1）若 $B\subseteq A$，求实数 $m$ 的取值范围；

（2）若 $A\cap B\neq\varnothing$，求实数 $m$ 的取值范围.

\ansspace*[10]

\ifanswer
\solbegin
\stepm{（1）}{当 $B=\varnothing$ 时，$m+1>2m-1$，解得 $m<2$，}
\step{当 $B\neq\varnothing$ 时，$\begin{cases}m+1\leqslant2m-1\\ m+1\geqslant-2\\ 2m-1\leqslant5\end{cases}$，解得 $2\leqslant m\leqslant3$，}
\step{综上，实数 $m$ 的取值范围是 $(-\infty,3]$；}
\stepm{（2）}{当 $B=\varnothing$ 时，$m+1>2m-1$，解得 $m<2$，}
\step{当 $B\neq\varnothing$ 时，$\begin{cases}m+1\leqslant2m-1\\ 2m-1<-2\end{cases}$ 或 $\begin{cases}m+1\leqslant2m-1\\ m+1>5\end{cases}$，解得 $m\in\varnothing$ 或 $m>4$.}
\step{所以当 $A\cap B\neq\varnothing$ 时，实数 $m$ 的取值范围是 $[2,4]$.}
\solend
\fi

\item 在解决实际问题时，往往会有不同的思路和方法，这些方法有些正确，有些错误；有些简洁，有些复杂.

问题①：设 $a\in R$，集合 $M=\{x\mid(x-1)(x-a)=0\}$，若 $x\in M$ 是 $x\in\{1,2,3\}$ 的充分条件，求：$a$ 的取值集合.

问题②：设 $x\in Z$，$a=x+3$，$b=18-x$，若 $c\in Z$，求证：$c-a$ 和 $c-b$ 至少有一个数是奇数

（1）小明在解决问题①，他认为原问题等价于 $\{1,a\}\subseteq\{1,2,3\}$，解得 $a$ 的取值集合为 $\{2,3\}$，张老师判断小明解题错误，请计算正确的 $a$ 的取值集合；

（2）小红认为既然 $c\in Z$，只需根据 $c$ 是奇数还是偶数，分类讨论即可；小华则认为可以使用反证法解决问题，请你选一种你认为更好的方法并证明.

\ansspace*[12]

\ifanswer
\solbegin
\stepm{（1）}{当 $a=1$ 时，集合 $M=\{x\mid(x-1)^2=0\}=\{1\}$，}
\step{此时 $x\in M\Rightarrow x\in\{1,2,3\}$，即 $x\in M$ 是 $x\in\{1,2,3\}$ 的充分条件.}
\step{当 $a\neq1$ 时，集合 $M=\{x\mid(x-1)(x-a)=0\}=\{1,a\}$，}
\step{由 $x\in M$ 是 $x\in\{1,2,3\}$ 的充分条件，得 $\{1,a\}\subseteq\{1,2,3\}$，}
\step{所以 $a=2$ 或 $a=3$.}
\step{综上，实数 $a$ 的取值集合为 $\{1,2,3\}$；}
\stepm{（2）}{利用反证法证明．证明如下：}
\step{假设 $c-a$、$c-b$ 均为偶数，则 $(c-b)-(c-a)=a-b$ 为偶数.}
\step{因为 $a-b=x+3-18+x=2x-15$ 为奇数，}
\step{与 $b-a$ 为偶数相矛盾，所以原假设不成立.}
\step{故 $c-a$、$c-b$ 至少有一个数是奇数.}
\solend
\fi

\item 已知实数 $x_1$、$x_2$、$x_3$、$x_4$、$x_5$，满足 $x_1+x_2+x_3+x_4+x_5=5$.

（1）求证：$x_1$、$x_2$、$x_2$、$x_4$、$x_5$ 中至少有一个实数不小于 $1$；% 原卷把 x_3 误印作 x_2，照录

（2）设 $x_1$、$x_2$、$x_3$、$x_4$、$x_5$ 这五个实数两两不等，集合 $A=\{x_1,x_2,x_3,x_4,x_5\}$，若 $B\neq\varnothing$ 且 $B\subseteq A$，记 $G(B)$ 是 $B$ 中所有元素之和，对所有的 $B$，求 $G(B)$ 的平均值.

\ansspace*[13]

\ifanswer
\solbegin
\stepm{（1）}{假设 $x_1$、$x_2$、$x_2$、$x_4$、$x_5$ 中的每一个数都小于 $1$，}% 原卷 x_3 误印作 x_2，照录
\step{得 $x_1+x_2+x_2+x_4+x_5<5$，这与 $x_1+x_2+x_3+x_4+x_5=5$ 矛盾，}% 前半句原卷 x_3 误印作 x_2，照录
\step{故 $x_1$、$x_2$、$x_2$、$x_4$、$x_5$ 中至少有一个实数不小于 $1$；}% 原卷 x_3 误印作 x_2，照录
\stepm{（2）}{集合 $A=\{x_1,x_2,x_3,x_4,x_5\}$ 的非空子集个数为 $32-1=31$，}
\step{由于 $G(B)$ 是 $B$ 中所有元素之和，}
\step{得 $G(B)=16(x_1+x_2+x_3+x_4+x_5)=16\times5=80$，}
\step{则 $G(B)$ 的平均值为 $\dfrac{80}{31}$.}
\solend
\fi

\end{enumerate}

\end{document}
"
%
% 原始材料：微信公众号“魔都数学加油站”文章，连载“27学年高一98”；
% 原文标题“27学年高一98——2026-2027年上海市市西中学高一上10月月考”；
% 原文链接：https://mp.weixin.qq.com/s/duK4wz-HT3W9tjOeYQaVfg。
% 正文为 8 张 PNG 页（`原图/高一98_市西中学_10月月考/01..08.png`），
% 尺寸混排（第 1—4、6、7 页为 4762×6735，第 5、8 页为 2560×3620），按页序读取。
%
% 篇幅与题号：原卷自第 3 题起编号（第 1、2 题不在原文内）。填空题 3—12（10 题）、
% 选择题 13—16（4 题）、解答题 17—21（5 题），共 19 题。
%
% 讲评版：全部 19 题均印【解析】，其中第 13—16 题的答案在解析末句
% “故选 $B$／$D$／$A$／$D$”（题干括号内未印答案）。全卷无插图。
%
% 副标题：原卷未印副标题。本稿按内容标作“集合与常用逻辑用语、不等式”。
%
% 原卷笔误/存疑（均照录未改，正文对应处另有 % 注）：
%   ① 第 21 题 (1) 求证行、解析“假设”行与“故”行三处把 $x_3$ 误印作 $x_2$
%      （“$x_1$、$x_2$、$x_2$、$x_4$、$x_5$”）；解析同一行后半句
%      “这与 $x_1+x_2+x_3+x_4+x_5=5$ 矛盾”却用 $x_3$，显系笔误，照录；
%   ② 第 6 题解析的判别式行印作“$\Delta=(-4)^2-4a\geqslant0$”，与
%      $f(1)>0$ 并列，照录。
%
% 与原卷的排版差异：
%   ① 原卷每题下直接印【解析】，本稿按全稿惯例将答案与解析置于答案版；
%   ② 原卷未印副标题，本稿补出；填空横线统一排 \blankline[4.4em]；
%   ③ 原卷数集记号排斜体 $R$、$\Z$，本稿按全稿惯例改用 $\R$、$\Z$；
%   ④ 解答题按惯例加 \ansspace；选择题括号排 \mbox{（\quad）}；
%   ⑤ 第 12 题定义的运算 $\odot$ 原卷排实心小圆，本稿排 $\odot$。
%
% 行数核对：逐题数 \step/\stepm 条数，合计 82 条（Q3：1；Q4：2；Q5：1；Q6：4；Q7：3；
% Q8：2；Q9：5；Q10：2；Q11：5；Q12：4；Q13：3；Q14：2；Q15：2；Q16：9；Q17：2；
% Q18：11；Q19：6；Q20：11；Q21：7）。全卷 19 题均印【解析】。
% 原卷解析的换行按语意合并，未逐行强拆。
%
\documentclass[a4paper,11pt]{article}
\usepackage{common}
\begin{document}

\examtitlest[7]{2026--2027 年市西中学高一上 10 月月考}{集合与常用逻辑用语、不等式}

\bigsec{一、填空题}

\begin{enumerate}
\setcounter{enumi}{2}

\item 集合 $A=\{x\mid y=\sqrt{1-x^2}\},B=\{y\mid y=x^2+1,x\in A\}$，则 $A\cap\overline B=\blankline[4.4em]$.

\ifanswer
\solbegin
\step{$1-x^2\geqslant0\Rightarrow A=[-1,1]$，从而 $B=[1,2]$，所以 $A\cap\overline B=[-1,1)$.}
\solend
\fi

\item 已知 $b,c\in R$，关于 $x$ 的不等式 $x^2-bx+c<0$ 的解集为 $(-3,2)$，则 $b+c=\blankline[4.4em]$.

\ifanswer
\solbegin
\step{由题意得 $x=-3$ 与 $x=2$ 是 $x^2-bx+c=0$ 的两根，}
\step{所以 $\begin{cases}-3+2=b\\ -3\times2=c\end{cases}$，解得 $\begin{cases}b=-1\\ c=-6\end{cases}$，所以 $b+c=-7$.}
\solend
\fi

\item 若 $\alpha:\lvert x-m\rvert<1$ 是 $\beta:x^2-4x<0$ 的充分条件，则实数 $m$ 的取值范围是\blankline[4.4em].

\ifanswer
\solbegin
\step{$(m-1,m+1)\subseteq(0,4)\Leftrightarrow m-1\geqslant0$ 且 $m+1\leqslant4\Leftrightarrow m\in[1,3]$.}
\solend
\fi

\item 已知方程 $x^2-4x+a=0$ 的两根都大于 $1$，则 $a$ 的取值范围为\blankline[4.4em].

\ifanswer
\solbegin
\step{设函数 $f(x)=x^2-4x+a$，对称轴为 $x=2$，}
\step{因为方程 $x^2-4x+a=0$ 的两根都大于 $1$，}
\step{所以 $\begin{cases}\Delta=(-4)^2-4a\geqslant0\\ f(1)>0\end{cases}$，即 $\begin{cases}16-4a\geqslant0\\ 1-4+a>0\end{cases}$，解得 $3<a\leqslant4$，}
\step{即 $a$ 的取值范围为 $(3,4]$.}
\solend
\fi

\item 若不等式 $(a-2)x^2+2(a-2)x-4<0$ 对 $x\in R$ 恒成立，则实数 $a$ 的取值范围为\blankline[4.4em].

\ifanswer
\solbegin
\step{①当 $a-2=0$，即 $a=2$ 时，不等式为 $-4<0$ 恒成立，符合题意；}
\step{②当 $a\neq2$ 时，则 $\begin{cases}a-2<0\\ \Delta=4(a-2)^2+16(a-2)<0\end{cases}$，解得 $-2<a<2$.}
\step{综上所述，实数 $a$ 的取值范围是 $(-2,2]$.}
\solend
\fi

\item 已知 $-1\leqslant a+b\leqslant4,2\leqslant a-b\leqslant3$，则 $3a-2b$ 的取值范围为\blankline[4.4em].

\ifanswer
\solbegin
\step{设 $3a-2b=m(a+b)+n(a-b)$，则 $m+n=3,m-n=-2$，}
\step{所以 $m=\dfrac12,n=\dfrac52$，所以 $3a-2b=\dfrac12(a+b)+\dfrac52(a-b)\in\left[\dfrac92,\dfrac{19}2\right]$.}
\solend
\fi

\item 设 $x_1$、$x_2$ 是关于 $x$ 的方程 $2x^2-4mx+2m^2+3m-2=0$ 的两个实数根，则 $\lvert x_1-x_2\rvert$ 的最小值为\blankline[4.4em].

\ifanswer
\solbegin
\step{$x_1$、$x_2$ 是方程 $2x^2-4mx+2m^2+3m-2=0$ 的两个实数根，}
\step{所以 $\Delta=(4m)^2-8(2m^2+3m-2)\geqslant0$，解得 $m\leqslant\dfrac23$，}
\step{由韦达定理得 $x_1+x_2=2m,x_1\cdot x_2=\dfrac{2m^2+3m-2}2$，}
\step{则 $\lvert x_1-x_2\rvert=\sqrt{(x_1+x_2)^2-4x_1\cdot x_2}=\sqrt{4-6m}$，}
\step{当 $m=\dfrac23$ 时，$\lvert x_1-x_2\rvert$ 的最小值为 $0$.}
\solend
\fi

\item 设 $x\in R$，$[x]$ 表示不大于 $x$ 的最大整数，如 $[\pi]=3$，$[-1.2]=-2$，$\left[\dfrac12\right]=0$，则使 $\lvert\lvert x^2-1\rvert\rvert=3$ 成立的 $x$ 的取值范围是\blankline[4.4em].

\ifanswer
\solbegin
\step{若 $\lvert\lvert x^2-1\rvert\rvert=3$，则 $3\leqslant\lvert x^2-1\rvert<4$，即 $4\leqslant x^2<5$，}
\step{解得 $2\leqslant x<\sqrt5$ 或 $-\sqrt5<x\leqslant-2$，故 $x$ 的取值范围是 $(-\sqrt5,-2]\cup[2,\sqrt5)$.}
\solend
\fi

\item 设 $a$ 为实数，若关于 $x$ 的一元一次不等式组 $\begin{cases}2x+a>0\\ 3x-6a<0\end{cases}$ 的解集中有且仅有 $4$ 个整数，则 $a$ 的取值范围是\blankline[4.4em].

\ifanswer
\solbegin
\step{一元一次不等式组 $\begin{cases}2x+a>0\\ 3x-6a<0\end{cases}\Leftrightarrow-\dfrac a2<x<2a$，且 $a>0$，}
\step{关于 $x$ 的一元一次不等式组 $\begin{cases}2x+a>0\\ 3x-6a<0\end{cases}$ 的解集中有且仅有 $4$ 个整数，}
\step{①当 $4$ 个整数为 $0$，$1$，$2$，$3$ 时，则 $\begin{cases}-1\leqslant-\dfrac a2<0\\ 3<2a\leqslant4\end{cases}$，所以 $\dfrac32<a\leqslant2$，}
\step{②当 $4$ 个整数为 $-1$，$0$，$1$，$2$ 时，则 $\begin{cases}-2\leqslant-\dfrac a2<-1\\ 2<2a\leqslant3\end{cases}$，无解，}
\step{综上，$a$ 的取值范围是 $\left(\dfrac32,2\right]$.}
\solend
\fi

\item 集合 $M$、$N$、$S$ 都是非空集合，现规定如下运算：

$M\odot N\odot S=\{x\mid x\in(M\cap N)\cup(N\cap S)\cup(S\cap M)$ 且 $x\notin M\cap N\cap S\}$．假设集合 $A=\{x\mid a<x<b\}$，$B=\{x\mid c<x<d\}$，$C=\{x\mid e<x<f\}$，其中实数 $a$、$b$、$c$、$d$、$e$、$f$ 满足：$a<c<e<0<b<d<f$．计算 $A\odot B\odot C=\blankline[4.4em]$.

\ifanswer
\solbegin
\step{因为集合 $A=\{x\mid a<x<b\}$，$B=\{x\mid c<x<d\}$，$C=\{x\mid e<x<f\}$，}
\step{又因为 $a<c<e<0<b<d<f$，}
\step{所以 $A\cap B=\{x\mid c<x<b\}$，$B\cap C=\{x\mid e<x<d\}$，$C\cap A=\{x\mid e<x<b\}$，}
\step{$A\cap B\cap C=\{x\mid e<x<b\}$，故 $A\odot B\odot C=\{x\mid c<x\leqslant e$ 或 $b\leqslant x<d\}$.}
\solend
\fi

\end{enumerate}

\bigsec{二、选择题}

\begin{enumerate}
\setcounter{enumi}{12}

\item “$x\neq1$”是“$x^2-1\neq0$”的\mbox{（\quad）}条件.

\keepwithprev
A. 充分非必要\qquad B. 必要非充分

C. 充要\qquad D. 既非充分也非必要

\ifanswer
\solbegin
\step{$x^2-1\neq0\Leftrightarrow x\neq1$ 且 $x\neq-1\Rightarrow x\neq1$.}
\step{反之不成立，如 $x=-1$ 时，$x\neq1$，但 $x^2-1=0$.}
\step{因此“$x\neq1$”是“$x^2-1\neq0$”的必要非充分条件，故选 $B$.}
\solend
\fi

\item 若 $\dfrac1a<\dfrac1b<0$，则下列结论中不正确的是\mbox{（\quad）}

\keepwithprev
A. $a+b<0$\qquad B. $a^2<b^2$\qquad C. $ab<b^2$\qquad D. $ab<a^2$

\ifanswer
\solbegin
\step{由 $\dfrac1a<\dfrac1b<0$ 得 $b<a<0$，因此 $a+b<0$，且 $a^2<ab<b^2$，}
\step{所以 A、B、C 正确，D 不正确．故选 $D$.}
\solend
\fi

\item 非空集合 $X=\{x\mid a+1\leqslant x\leqslant3a-5\}$，$Y=\{x\mid1\leqslant x\leqslant16\}$，使得 $X\subseteq(X\cap Y)$ 成立的所有 $a$ 的集合是\mbox{（\quad）}

\keepwithprev
A. $\{a\mid3\leqslant a\leqslant7\}$\qquad B. $\{a\mid0\leqslant a\leqslant7\}$\qquad C. $\{a\mid3<a\leqslant7\}$\qquad D. $\{a\mid a\leqslant7\}$

\ifanswer
\solbegin
\step{因为 $X\subseteq(X\cap Y)$ 成立，所以 $X\subseteq Y$ 成立.}
\step{$X\neq\varnothing$，要使 $X\subseteq Y$，则 $\begin{cases}a\geqslant3\\ a+1\geqslant1\\ 3a-5\leqslant16\end{cases}$，即 $\begin{cases}a\geqslant3\\ a\geqslant0\\ a\leqslant7\end{cases}$，所以 $3\leqslant a\leqslant7$，故选 $A$.}
\solend
\fi

\item 已知 $M$ 是满足下列条件的集合：①$0\in M$，$1\in M$；②若 $x$、$y\in M$，则 $x-y\in M$；③若 $x\in M$ 且 $x\neq0$，则 $\dfrac1x\in M$．则下列说法正确的个数为\mbox{（\quad）}

（1）$\dfrac12\in M$；（2）$x+y\in M$；（3）$x^2\in M$．

\keepwithprev
A. $0$\qquad B. $1$\qquad C. $2$\qquad D. $3$

\ifanswer
\solbegin
\step{（1）$\dfrac12\in M$ 正确，理由如下：}
\step{由①得 $0\in M$，$1\in M$，由②得 $0-1=-1\in M$，$1-(-1)=2\in M$，}
\step{由③得 $\dfrac12\in M$；}
\step{（2）由①得 $0\in M$，由题意 $y\in M$，}
\step{所以由②得 $0-y=-y\in M$，又 $x\in M$，所以 $x-(-y)=x+y\in M$；}
\step{（3）$x\in M$，由②得 $x-1\in M$，由③得 $\dfrac1x\in M$，$\dfrac1{x-1}\in M$，}
\step{所以 $\dfrac1x-\dfrac1{x-1}\in M$，即 $\dfrac1{x(x-1)}\in M$，所以 $x(x-1)\in M$，}
\step{由（2）结论得 $x(x-1)+x\in M$，即 $x^2\in M$，}
\step{所以说法正确的个数为 $3$．故选 $D$.}
\solend
\fi

\end{enumerate}

\bigsec{三、解答题}

\begin{enumerate}
\setcounter{enumi}{16}

\item 已知 $a>b$，用比较法证明：$a^3>b^3$.

\ansspace*[8]

\ifanswer
\solbegin
\step{$a>b$，$a^3-b^3=(a-b)(a^2+ab+b^2)=(a-b)\left[\left(a+\dfrac b2\right)^2+\dfrac34b^2\right]>0$，}
\step{故 $a^3>b^3$.}
\solend
\fi

\item 求下列关于 $x$ 的不等式的解集：

（1）$\dfrac{x^2(x+3)(x^2-x+1)}{x^2-5x+4}>0$；

（2）$\dfrac{ax-1}{x-2}\geqslant0\left(a<\dfrac12\right)$.

\ansspace*[12]

\ifanswer
\solbegin
\stepm{（1）}{分母 $x^2-5x+4=(x-1)(x-4)$，要求 $x\neq1,4$；}
\step{又 $x=0$ 时原不等式不成立，且 $x^2-x+1=\left(x-\dfrac12\right)^2+\dfrac34>0$，}
\step{因此原不等式等价于 $\dfrac{x+3}{(x-1)(x-4)}>0,x\neq0$.}
\step{在数轴上标出 $-3,1,4$，分段判号得 $-3<x<1$ 或 $x>4$，}
\step{再排除 $x=0$，解集为 $(-3,0)\cup(0,1)\cup(4,+\infty)$.}
\stepm{（2）}{原不等式要求 $x\neq2$，按 $a$ 的符号分类讨论.}
\step{①当 $a<0$ 时，$\dfrac1a<2$，原不等式等价于 $\dfrac{x-\dfrac1a}{x-2}\leqslant0\Leftrightarrow\dfrac1a\leqslant x<2$，}
\step{解集为 $\left[\dfrac1a,2\right)$.}
\step{②当 $a=0$ 时，原不等式为 $\dfrac{-1}{x-2}\geqslant0\Leftrightarrow x<2$，解集为 $(-\infty,2)$.}
\step{③当 $0<a<\dfrac12$ 时，$\dfrac1a>2$，}
\step{原不等式等价于 $\dfrac{x-\dfrac1a}{x-2}\geqslant0\Leftrightarrow x<2$ 或 $x\geqslant\dfrac1a$，解集为 $(-\infty,2)\cup\left[\dfrac1a,+\infty\right)$.}
\solend
\fi

\item 集合 $A=\{x\mid-2\leqslant x\leqslant5\}$，$B=\{x\mid m+1\leqslant x\leqslant2m-1\}$.

（1）若 $B\subseteq A$，求实数 $m$ 的取值范围；

（2）若 $A\cap B\neq\varnothing$，求实数 $m$ 的取值范围.

\ansspace*[10]

\ifanswer
\solbegin
\stepm{（1）}{当 $B=\varnothing$ 时，$m+1>2m-1$，解得 $m<2$，}
\step{当 $B\neq\varnothing$ 时，$\begin{cases}m+1\leqslant2m-1\\ m+1\geqslant-2\\ 2m-1\leqslant5\end{cases}$，解得 $2\leqslant m\leqslant3$，}
\step{综上，实数 $m$ 的取值范围是 $(-\infty,3]$；}
\stepm{（2）}{当 $B=\varnothing$ 时，$m+1>2m-1$，解得 $m<2$，}
\step{当 $B\neq\varnothing$ 时，$\begin{cases}m+1\leqslant2m-1\\ 2m-1<-2\end{cases}$ 或 $\begin{cases}m+1\leqslant2m-1\\ m+1>5\end{cases}$，解得 $m\in\varnothing$ 或 $m>4$.}
\step{所以当 $A\cap B\neq\varnothing$ 时，实数 $m$ 的取值范围是 $[2,4]$.}
\solend
\fi

\item 在解决实际问题时，往往会有不同的思路和方法，这些方法有些正确，有些错误；有些简洁，有些复杂.

问题①：设 $a\in R$，集合 $M=\{x\mid(x-1)(x-a)=0\}$，若 $x\in M$ 是 $x\in\{1,2,3\}$ 的充分条件，求：$a$ 的取值集合.

问题②：设 $x\in Z$，$a=x+3$，$b=18-x$，若 $c\in Z$，求证：$c-a$ 和 $c-b$ 至少有一个数是奇数

（1）小明在解决问题①，他认为原问题等价于 $\{1,a\}\subseteq\{1,2,3\}$，解得 $a$ 的取值集合为 $\{2,3\}$，张老师判断小明解题错误，请计算正确的 $a$ 的取值集合；

（2）小红认为既然 $c\in Z$，只需根据 $c$ 是奇数还是偶数，分类讨论即可；小华则认为可以使用反证法解决问题，请你选一种你认为更好的方法并证明.

\ansspace*[12]

\ifanswer
\solbegin
\stepm{（1）}{当 $a=1$ 时，集合 $M=\{x\mid(x-1)^2=0\}=\{1\}$，}
\step{此时 $x\in M\Rightarrow x\in\{1,2,3\}$，即 $x\in M$ 是 $x\in\{1,2,3\}$ 的充分条件.}
\step{当 $a\neq1$ 时，集合 $M=\{x\mid(x-1)(x-a)=0\}=\{1,a\}$，}
\step{由 $x\in M$ 是 $x\in\{1,2,3\}$ 的充分条件，得 $\{1,a\}\subseteq\{1,2,3\}$，}
\step{所以 $a=2$ 或 $a=3$.}
\step{综上，实数 $a$ 的取值集合为 $\{1,2,3\}$；}
\stepm{（2）}{利用反证法证明．证明如下：}
\step{假设 $c-a$、$c-b$ 均为偶数，则 $(c-b)-(c-a)=a-b$ 为偶数.}
\step{因为 $a-b=x+3-18+x=2x-15$ 为奇数，}
\step{与 $b-a$ 为偶数相矛盾，所以原假设不成立.}
\step{故 $c-a$、$c-b$ 至少有一个数是奇数.}
\solend
\fi

\item 已知实数 $x_1$、$x_2$、$x_3$、$x_4$、$x_5$，满足 $x_1+x_2+x_3+x_4+x_5=5$.

（1）求证：$x_1$、$x_2$、$x_2$、$x_4$、$x_5$ 中至少有一个实数不小于 $1$；% 原卷把 x_3 误印作 x_2，照录

（2）设 $x_1$、$x_2$、$x_3$、$x_4$、$x_5$ 这五个实数两两不等，集合 $A=\{x_1,x_2,x_3,x_4,x_5\}$，若 $B\neq\varnothing$ 且 $B\subseteq A$，记 $G(B)$ 是 $B$ 中所有元素之和，对所有的 $B$，求 $G(B)$ 的平均值.

\ansspace*[13]

\ifanswer
\solbegin
\stepm{（1）}{假设 $x_1$、$x_2$、$x_2$、$x_4$、$x_5$ 中的每一个数都小于 $1$，}% 原卷 x_3 误印作 x_2，照录
\step{得 $x_1+x_2+x_2+x_4+x_5<5$，这与 $x_1+x_2+x_3+x_4+x_5=5$ 矛盾，}% 前半句原卷 x_3 误印作 x_2，照录
\step{故 $x_1$、$x_2$、$x_2$、$x_4$、$x_5$ 中至少有一个实数不小于 $1$；}% 原卷 x_3 误印作 x_2，照录
\stepm{（2）}{集合 $A=\{x_1,x_2,x_3,x_4,x_5\}$ 的非空子集个数为 $32-1=31$，}
\step{由于 $G(B)$ 是 $B$ 中所有元素之和，}
\step{得 $G(B)=16(x_1+x_2+x_3+x_4+x_5)=16\times5=80$，}
\step{则 $G(B)$ 的平均值为 $\dfrac{80}{31}$.}
\solend
\fi

\end{enumerate}

\end{document}
"
% 紧凑版：xelatex "\def\iscompact{1}% 高一98_市西中学_10月月考.tex
%   —— 高一上数学“2026-2027 年市西中学高一上 10 月月考”（讲评版）
% 试卷版：xelatex 高一98_市西中学_10月月考.tex
% 答案版：xelatex "\def\isanswer{1}% 高一98_市西中学_10月月考.tex
%   —— 高一上数学“2026-2027 年市西中学高一上 10 月月考”（讲评版）
% 试卷版：xelatex 高一98_市西中学_10月月考.tex
% 答案版：xelatex "\def\isanswer{1}\input{高一98_市西中学_10月月考.tex}"
% 紧凑版：xelatex "\def\iscompact{1}\input{高一98_市西中学_10月月考.tex}"
%
% 原始材料：微信公众号“魔都数学加油站”文章，连载“27学年高一98”；
% 原文标题“27学年高一98——2026-2027年上海市市西中学高一上10月月考”；
% 原文链接：https://mp.weixin.qq.com/s/duK4wz-HT3W9tjOeYQaVfg。
% 正文为 8 张 PNG 页（`原图/高一98_市西中学_10月月考/01..08.png`），
% 尺寸混排（第 1—4、6、7 页为 4762×6735，第 5、8 页为 2560×3620），按页序读取。
%
% 篇幅与题号：原卷自第 3 题起编号（第 1、2 题不在原文内）。填空题 3—12（10 题）、
% 选择题 13—16（4 题）、解答题 17—21（5 题），共 19 题。
%
% 讲评版：全部 19 题均印【解析】，其中第 13—16 题的答案在解析末句
% “故选 $B$／$D$／$A$／$D$”（题干括号内未印答案）。全卷无插图。
%
% 副标题：原卷未印副标题。本稿按内容标作“集合与常用逻辑用语、不等式”。
%
% 原卷笔误/存疑（均照录未改，正文对应处另有 % 注）：
%   ① 第 21 题 (1) 求证行、解析“假设”行与“故”行三处把 $x_3$ 误印作 $x_2$
%      （“$x_1$、$x_2$、$x_2$、$x_4$、$x_5$”）；解析同一行后半句
%      “这与 $x_1+x_2+x_3+x_4+x_5=5$ 矛盾”却用 $x_3$，显系笔误，照录；
%   ② 第 6 题解析的判别式行印作“$\Delta=(-4)^2-4a\geqslant0$”，与
%      $f(1)>0$ 并列，照录。
%
% 与原卷的排版差异：
%   ① 原卷每题下直接印【解析】，本稿按全稿惯例将答案与解析置于答案版；
%   ② 原卷未印副标题，本稿补出；填空横线统一排 \blankline[4.4em]；
%   ③ 原卷数集记号排斜体 $R$、$\Z$，本稿按全稿惯例改用 $\R$、$\Z$；
%   ④ 解答题按惯例加 \ansspace；选择题括号排 \mbox{（\quad）}；
%   ⑤ 第 12 题定义的运算 $\odot$ 原卷排实心小圆，本稿排 $\odot$。
%
% 行数核对：逐题数 \step/\stepm 条数，合计 82 条（Q3：1；Q4：2；Q5：1；Q6：4；Q7：3；
% Q8：2；Q9：5；Q10：2；Q11：5；Q12：4；Q13：3；Q14：2；Q15：2；Q16：9；Q17：2；
% Q18：11；Q19：6；Q20：11；Q21：7）。全卷 19 题均印【解析】。
% 原卷解析的换行按语意合并，未逐行强拆。
%
\documentclass[a4paper,11pt]{article}
\usepackage{common}
\begin{document}

\examtitlest[7]{2026--2027 年市西中学高一上 10 月月考}{集合与常用逻辑用语、不等式}

\bigsec{一、填空题}

\begin{enumerate}
\setcounter{enumi}{2}

\item 集合 $A=\{x\mid y=\sqrt{1-x^2}\},B=\{y\mid y=x^2+1,x\in A\}$，则 $A\cap\overline B=\blankline[4.4em]$.

\ifanswer
\solbegin
\step{$1-x^2\geqslant0\Rightarrow A=[-1,1]$，从而 $B=[1,2]$，所以 $A\cap\overline B=[-1,1)$.}
\solend
\fi

\item 已知 $b,c\in R$，关于 $x$ 的不等式 $x^2-bx+c<0$ 的解集为 $(-3,2)$，则 $b+c=\blankline[4.4em]$.

\ifanswer
\solbegin
\step{由题意得 $x=-3$ 与 $x=2$ 是 $x^2-bx+c=0$ 的两根，}
\step{所以 $\begin{cases}-3+2=b\\ -3\times2=c\end{cases}$，解得 $\begin{cases}b=-1\\ c=-6\end{cases}$，所以 $b+c=-7$.}
\solend
\fi

\item 若 $\alpha:\lvert x-m\rvert<1$ 是 $\beta:x^2-4x<0$ 的充分条件，则实数 $m$ 的取值范围是\blankline[4.4em].

\ifanswer
\solbegin
\step{$(m-1,m+1)\subseteq(0,4)\Leftrightarrow m-1\geqslant0$ 且 $m+1\leqslant4\Leftrightarrow m\in[1,3]$.}
\solend
\fi

\item 已知方程 $x^2-4x+a=0$ 的两根都大于 $1$，则 $a$ 的取值范围为\blankline[4.4em].

\ifanswer
\solbegin
\step{设函数 $f(x)=x^2-4x+a$，对称轴为 $x=2$，}
\step{因为方程 $x^2-4x+a=0$ 的两根都大于 $1$，}
\step{所以 $\begin{cases}\Delta=(-4)^2-4a\geqslant0\\ f(1)>0\end{cases}$，即 $\begin{cases}16-4a\geqslant0\\ 1-4+a>0\end{cases}$，解得 $3<a\leqslant4$，}
\step{即 $a$ 的取值范围为 $(3,4]$.}
\solend
\fi

\item 若不等式 $(a-2)x^2+2(a-2)x-4<0$ 对 $x\in R$ 恒成立，则实数 $a$ 的取值范围为\blankline[4.4em].

\ifanswer
\solbegin
\step{①当 $a-2=0$，即 $a=2$ 时，不等式为 $-4<0$ 恒成立，符合题意；}
\step{②当 $a\neq2$ 时，则 $\begin{cases}a-2<0\\ \Delta=4(a-2)^2+16(a-2)<0\end{cases}$，解得 $-2<a<2$.}
\step{综上所述，实数 $a$ 的取值范围是 $(-2,2]$.}
\solend
\fi

\item 已知 $-1\leqslant a+b\leqslant4,2\leqslant a-b\leqslant3$，则 $3a-2b$ 的取值范围为\blankline[4.4em].

\ifanswer
\solbegin
\step{设 $3a-2b=m(a+b)+n(a-b)$，则 $m+n=3,m-n=-2$，}
\step{所以 $m=\dfrac12,n=\dfrac52$，所以 $3a-2b=\dfrac12(a+b)+\dfrac52(a-b)\in\left[\dfrac92,\dfrac{19}2\right]$.}
\solend
\fi

\item 设 $x_1$、$x_2$ 是关于 $x$ 的方程 $2x^2-4mx+2m^2+3m-2=0$ 的两个实数根，则 $\lvert x_1-x_2\rvert$ 的最小值为\blankline[4.4em].

\ifanswer
\solbegin
\step{$x_1$、$x_2$ 是方程 $2x^2-4mx+2m^2+3m-2=0$ 的两个实数根，}
\step{所以 $\Delta=(4m)^2-8(2m^2+3m-2)\geqslant0$，解得 $m\leqslant\dfrac23$，}
\step{由韦达定理得 $x_1+x_2=2m,x_1\cdot x_2=\dfrac{2m^2+3m-2}2$，}
\step{则 $\lvert x_1-x_2\rvert=\sqrt{(x_1+x_2)^2-4x_1\cdot x_2}=\sqrt{4-6m}$，}
\step{当 $m=\dfrac23$ 时，$\lvert x_1-x_2\rvert$ 的最小值为 $0$.}
\solend
\fi

\item 设 $x\in R$，$[x]$ 表示不大于 $x$ 的最大整数，如 $[\pi]=3$，$[-1.2]=-2$，$\left[\dfrac12\right]=0$，则使 $\lvert\lvert x^2-1\rvert\rvert=3$ 成立的 $x$ 的取值范围是\blankline[4.4em].

\ifanswer
\solbegin
\step{若 $\lvert\lvert x^2-1\rvert\rvert=3$，则 $3\leqslant\lvert x^2-1\rvert<4$，即 $4\leqslant x^2<5$，}
\step{解得 $2\leqslant x<\sqrt5$ 或 $-\sqrt5<x\leqslant-2$，故 $x$ 的取值范围是 $(-\sqrt5,-2]\cup[2,\sqrt5)$.}
\solend
\fi

\item 设 $a$ 为实数，若关于 $x$ 的一元一次不等式组 $\begin{cases}2x+a>0\\ 3x-6a<0\end{cases}$ 的解集中有且仅有 $4$ 个整数，则 $a$ 的取值范围是\blankline[4.4em].

\ifanswer
\solbegin
\step{一元一次不等式组 $\begin{cases}2x+a>0\\ 3x-6a<0\end{cases}\Leftrightarrow-\dfrac a2<x<2a$，且 $a>0$，}
\step{关于 $x$ 的一元一次不等式组 $\begin{cases}2x+a>0\\ 3x-6a<0\end{cases}$ 的解集中有且仅有 $4$ 个整数，}
\step{①当 $4$ 个整数为 $0$，$1$，$2$，$3$ 时，则 $\begin{cases}-1\leqslant-\dfrac a2<0\\ 3<2a\leqslant4\end{cases}$，所以 $\dfrac32<a\leqslant2$，}
\step{②当 $4$ 个整数为 $-1$，$0$，$1$，$2$ 时，则 $\begin{cases}-2\leqslant-\dfrac a2<-1\\ 2<2a\leqslant3\end{cases}$，无解，}
\step{综上，$a$ 的取值范围是 $\left(\dfrac32,2\right]$.}
\solend
\fi

\item 集合 $M$、$N$、$S$ 都是非空集合，现规定如下运算：

$M\odot N\odot S=\{x\mid x\in(M\cap N)\cup(N\cap S)\cup(S\cap M)$ 且 $x\notin M\cap N\cap S\}$．假设集合 $A=\{x\mid a<x<b\}$，$B=\{x\mid c<x<d\}$，$C=\{x\mid e<x<f\}$，其中实数 $a$、$b$、$c$、$d$、$e$、$f$ 满足：$a<c<e<0<b<d<f$．计算 $A\odot B\odot C=\blankline[4.4em]$.

\ifanswer
\solbegin
\step{因为集合 $A=\{x\mid a<x<b\}$，$B=\{x\mid c<x<d\}$，$C=\{x\mid e<x<f\}$，}
\step{又因为 $a<c<e<0<b<d<f$，}
\step{所以 $A\cap B=\{x\mid c<x<b\}$，$B\cap C=\{x\mid e<x<d\}$，$C\cap A=\{x\mid e<x<b\}$，}
\step{$A\cap B\cap C=\{x\mid e<x<b\}$，故 $A\odot B\odot C=\{x\mid c<x\leqslant e$ 或 $b\leqslant x<d\}$.}
\solend
\fi

\end{enumerate}

\bigsec{二、选择题}

\begin{enumerate}
\setcounter{enumi}{12}

\item “$x\neq1$”是“$x^2-1\neq0$”的\mbox{（\quad）}条件.

\keepwithprev
A. 充分非必要\qquad B. 必要非充分

C. 充要\qquad D. 既非充分也非必要

\ifanswer
\solbegin
\step{$x^2-1\neq0\Leftrightarrow x\neq1$ 且 $x\neq-1\Rightarrow x\neq1$.}
\step{反之不成立，如 $x=-1$ 时，$x\neq1$，但 $x^2-1=0$.}
\step{因此“$x\neq1$”是“$x^2-1\neq0$”的必要非充分条件，故选 $B$.}
\solend
\fi

\item 若 $\dfrac1a<\dfrac1b<0$，则下列结论中不正确的是\mbox{（\quad）}

\keepwithprev
A. $a+b<0$\qquad B. $a^2<b^2$\qquad C. $ab<b^2$\qquad D. $ab<a^2$

\ifanswer
\solbegin
\step{由 $\dfrac1a<\dfrac1b<0$ 得 $b<a<0$，因此 $a+b<0$，且 $a^2<ab<b^2$，}
\step{所以 A、B、C 正确，D 不正确．故选 $D$.}
\solend
\fi

\item 非空集合 $X=\{x\mid a+1\leqslant x\leqslant3a-5\}$，$Y=\{x\mid1\leqslant x\leqslant16\}$，使得 $X\subseteq(X\cap Y)$ 成立的所有 $a$ 的集合是\mbox{（\quad）}

\keepwithprev
A. $\{a\mid3\leqslant a\leqslant7\}$\qquad B. $\{a\mid0\leqslant a\leqslant7\}$\qquad C. $\{a\mid3<a\leqslant7\}$\qquad D. $\{a\mid a\leqslant7\}$

\ifanswer
\solbegin
\step{因为 $X\subseteq(X\cap Y)$ 成立，所以 $X\subseteq Y$ 成立.}
\step{$X\neq\varnothing$，要使 $X\subseteq Y$，则 $\begin{cases}a\geqslant3\\ a+1\geqslant1\\ 3a-5\leqslant16\end{cases}$，即 $\begin{cases}a\geqslant3\\ a\geqslant0\\ a\leqslant7\end{cases}$，所以 $3\leqslant a\leqslant7$，故选 $A$.}
\solend
\fi

\item 已知 $M$ 是满足下列条件的集合：①$0\in M$，$1\in M$；②若 $x$、$y\in M$，则 $x-y\in M$；③若 $x\in M$ 且 $x\neq0$，则 $\dfrac1x\in M$．则下列说法正确的个数为\mbox{（\quad）}

（1）$\dfrac12\in M$；（2）$x+y\in M$；（3）$x^2\in M$．

\keepwithprev
A. $0$\qquad B. $1$\qquad C. $2$\qquad D. $3$

\ifanswer
\solbegin
\step{（1）$\dfrac12\in M$ 正确，理由如下：}
\step{由①得 $0\in M$，$1\in M$，由②得 $0-1=-1\in M$，$1-(-1)=2\in M$，}
\step{由③得 $\dfrac12\in M$；}
\step{（2）由①得 $0\in M$，由题意 $y\in M$，}
\step{所以由②得 $0-y=-y\in M$，又 $x\in M$，所以 $x-(-y)=x+y\in M$；}
\step{（3）$x\in M$，由②得 $x-1\in M$，由③得 $\dfrac1x\in M$，$\dfrac1{x-1}\in M$，}
\step{所以 $\dfrac1x-\dfrac1{x-1}\in M$，即 $\dfrac1{x(x-1)}\in M$，所以 $x(x-1)\in M$，}
\step{由（2）结论得 $x(x-1)+x\in M$，即 $x^2\in M$，}
\step{所以说法正确的个数为 $3$．故选 $D$.}
\solend
\fi

\end{enumerate}

\bigsec{三、解答题}

\begin{enumerate}
\setcounter{enumi}{16}

\item 已知 $a>b$，用比较法证明：$a^3>b^3$.

\ansspace*[8]

\ifanswer
\solbegin
\step{$a>b$，$a^3-b^3=(a-b)(a^2+ab+b^2)=(a-b)\left[\left(a+\dfrac b2\right)^2+\dfrac34b^2\right]>0$，}
\step{故 $a^3>b^3$.}
\solend
\fi

\item 求下列关于 $x$ 的不等式的解集：

（1）$\dfrac{x^2(x+3)(x^2-x+1)}{x^2-5x+4}>0$；

（2）$\dfrac{ax-1}{x-2}\geqslant0\left(a<\dfrac12\right)$.

\ansspace*[12]

\ifanswer
\solbegin
\stepm{（1）}{分母 $x^2-5x+4=(x-1)(x-4)$，要求 $x\neq1,4$；}
\step{又 $x=0$ 时原不等式不成立，且 $x^2-x+1=\left(x-\dfrac12\right)^2+\dfrac34>0$，}
\step{因此原不等式等价于 $\dfrac{x+3}{(x-1)(x-4)}>0,x\neq0$.}
\step{在数轴上标出 $-3,1,4$，分段判号得 $-3<x<1$ 或 $x>4$，}
\step{再排除 $x=0$，解集为 $(-3,0)\cup(0,1)\cup(4,+\infty)$.}
\stepm{（2）}{原不等式要求 $x\neq2$，按 $a$ 的符号分类讨论.}
\step{①当 $a<0$ 时，$\dfrac1a<2$，原不等式等价于 $\dfrac{x-\dfrac1a}{x-2}\leqslant0\Leftrightarrow\dfrac1a\leqslant x<2$，}
\step{解集为 $\left[\dfrac1a,2\right)$.}
\step{②当 $a=0$ 时，原不等式为 $\dfrac{-1}{x-2}\geqslant0\Leftrightarrow x<2$，解集为 $(-\infty,2)$.}
\step{③当 $0<a<\dfrac12$ 时，$\dfrac1a>2$，}
\step{原不等式等价于 $\dfrac{x-\dfrac1a}{x-2}\geqslant0\Leftrightarrow x<2$ 或 $x\geqslant\dfrac1a$，解集为 $(-\infty,2)\cup\left[\dfrac1a,+\infty\right)$.}
\solend
\fi

\item 集合 $A=\{x\mid-2\leqslant x\leqslant5\}$，$B=\{x\mid m+1\leqslant x\leqslant2m-1\}$.

（1）若 $B\subseteq A$，求实数 $m$ 的取值范围；

（2）若 $A\cap B\neq\varnothing$，求实数 $m$ 的取值范围.

\ansspace*[10]

\ifanswer
\solbegin
\stepm{（1）}{当 $B=\varnothing$ 时，$m+1>2m-1$，解得 $m<2$，}
\step{当 $B\neq\varnothing$ 时，$\begin{cases}m+1\leqslant2m-1\\ m+1\geqslant-2\\ 2m-1\leqslant5\end{cases}$，解得 $2\leqslant m\leqslant3$，}
\step{综上，实数 $m$ 的取值范围是 $(-\infty,3]$；}
\stepm{（2）}{当 $B=\varnothing$ 时，$m+1>2m-1$，解得 $m<2$，}
\step{当 $B\neq\varnothing$ 时，$\begin{cases}m+1\leqslant2m-1\\ 2m-1<-2\end{cases}$ 或 $\begin{cases}m+1\leqslant2m-1\\ m+1>5\end{cases}$，解得 $m\in\varnothing$ 或 $m>4$.}
\step{所以当 $A\cap B\neq\varnothing$ 时，实数 $m$ 的取值范围是 $[2,4]$.}
\solend
\fi

\item 在解决实际问题时，往往会有不同的思路和方法，这些方法有些正确，有些错误；有些简洁，有些复杂.

问题①：设 $a\in R$，集合 $M=\{x\mid(x-1)(x-a)=0\}$，若 $x\in M$ 是 $x\in\{1,2,3\}$ 的充分条件，求：$a$ 的取值集合.

问题②：设 $x\in Z$，$a=x+3$，$b=18-x$，若 $c\in Z$，求证：$c-a$ 和 $c-b$ 至少有一个数是奇数

（1）小明在解决问题①，他认为原问题等价于 $\{1,a\}\subseteq\{1,2,3\}$，解得 $a$ 的取值集合为 $\{2,3\}$，张老师判断小明解题错误，请计算正确的 $a$ 的取值集合；

（2）小红认为既然 $c\in Z$，只需根据 $c$ 是奇数还是偶数，分类讨论即可；小华则认为可以使用反证法解决问题，请你选一种你认为更好的方法并证明.

\ansspace*[12]

\ifanswer
\solbegin
\stepm{（1）}{当 $a=1$ 时，集合 $M=\{x\mid(x-1)^2=0\}=\{1\}$，}
\step{此时 $x\in M\Rightarrow x\in\{1,2,3\}$，即 $x\in M$ 是 $x\in\{1,2,3\}$ 的充分条件.}
\step{当 $a\neq1$ 时，集合 $M=\{x\mid(x-1)(x-a)=0\}=\{1,a\}$，}
\step{由 $x\in M$ 是 $x\in\{1,2,3\}$ 的充分条件，得 $\{1,a\}\subseteq\{1,2,3\}$，}
\step{所以 $a=2$ 或 $a=3$.}
\step{综上，实数 $a$ 的取值集合为 $\{1,2,3\}$；}
\stepm{（2）}{利用反证法证明．证明如下：}
\step{假设 $c-a$、$c-b$ 均为偶数，则 $(c-b)-(c-a)=a-b$ 为偶数.}
\step{因为 $a-b=x+3-18+x=2x-15$ 为奇数，}
\step{与 $b-a$ 为偶数相矛盾，所以原假设不成立.}
\step{故 $c-a$、$c-b$ 至少有一个数是奇数.}
\solend
\fi

\item 已知实数 $x_1$、$x_2$、$x_3$、$x_4$、$x_5$，满足 $x_1+x_2+x_3+x_4+x_5=5$.

（1）求证：$x_1$、$x_2$、$x_2$、$x_4$、$x_5$ 中至少有一个实数不小于 $1$；% 原卷把 x_3 误印作 x_2，照录

（2）设 $x_1$、$x_2$、$x_3$、$x_4$、$x_5$ 这五个实数两两不等，集合 $A=\{x_1,x_2,x_3,x_4,x_5\}$，若 $B\neq\varnothing$ 且 $B\subseteq A$，记 $G(B)$ 是 $B$ 中所有元素之和，对所有的 $B$，求 $G(B)$ 的平均值.

\ansspace*[13]

\ifanswer
\solbegin
\stepm{（1）}{假设 $x_1$、$x_2$、$x_2$、$x_4$、$x_5$ 中的每一个数都小于 $1$，}% 原卷 x_3 误印作 x_2，照录
\step{得 $x_1+x_2+x_2+x_4+x_5<5$，这与 $x_1+x_2+x_3+x_4+x_5=5$ 矛盾，}% 前半句原卷 x_3 误印作 x_2，照录
\step{故 $x_1$、$x_2$、$x_2$、$x_4$、$x_5$ 中至少有一个实数不小于 $1$；}% 原卷 x_3 误印作 x_2，照录
\stepm{（2）}{集合 $A=\{x_1,x_2,x_3,x_4,x_5\}$ 的非空子集个数为 $32-1=31$，}
\step{由于 $G(B)$ 是 $B$ 中所有元素之和，}
\step{得 $G(B)=16(x_1+x_2+x_3+x_4+x_5)=16\times5=80$，}
\step{则 $G(B)$ 的平均值为 $\dfrac{80}{31}$.}
\solend
\fi

\end{enumerate}

\end{document}
"
% 紧凑版：xelatex "\def\iscompact{1}% 高一98_市西中学_10月月考.tex
%   —— 高一上数学“2026-2027 年市西中学高一上 10 月月考”（讲评版）
% 试卷版：xelatex 高一98_市西中学_10月月考.tex
% 答案版：xelatex "\def\isanswer{1}\input{高一98_市西中学_10月月考.tex}"
% 紧凑版：xelatex "\def\iscompact{1}\input{高一98_市西中学_10月月考.tex}"
%
% 原始材料：微信公众号“魔都数学加油站”文章，连载“27学年高一98”；
% 原文标题“27学年高一98——2026-2027年上海市市西中学高一上10月月考”；
% 原文链接：https://mp.weixin.qq.com/s/duK4wz-HT3W9tjOeYQaVfg。
% 正文为 8 张 PNG 页（`原图/高一98_市西中学_10月月考/01..08.png`），
% 尺寸混排（第 1—4、6、7 页为 4762×6735，第 5、8 页为 2560×3620），按页序读取。
%
% 篇幅与题号：原卷自第 3 题起编号（第 1、2 题不在原文内）。填空题 3—12（10 题）、
% 选择题 13—16（4 题）、解答题 17—21（5 题），共 19 题。
%
% 讲评版：全部 19 题均印【解析】，其中第 13—16 题的答案在解析末句
% “故选 $B$／$D$／$A$／$D$”（题干括号内未印答案）。全卷无插图。
%
% 副标题：原卷未印副标题。本稿按内容标作“集合与常用逻辑用语、不等式”。
%
% 原卷笔误/存疑（均照录未改，正文对应处另有 % 注）：
%   ① 第 21 题 (1) 求证行、解析“假设”行与“故”行三处把 $x_3$ 误印作 $x_2$
%      （“$x_1$、$x_2$、$x_2$、$x_4$、$x_5$”）；解析同一行后半句
%      “这与 $x_1+x_2+x_3+x_4+x_5=5$ 矛盾”却用 $x_3$，显系笔误，照录；
%   ② 第 6 题解析的判别式行印作“$\Delta=(-4)^2-4a\geqslant0$”，与
%      $f(1)>0$ 并列，照录。
%
% 与原卷的排版差异：
%   ① 原卷每题下直接印【解析】，本稿按全稿惯例将答案与解析置于答案版；
%   ② 原卷未印副标题，本稿补出；填空横线统一排 \blankline[4.4em]；
%   ③ 原卷数集记号排斜体 $R$、$\Z$，本稿按全稿惯例改用 $\R$、$\Z$；
%   ④ 解答题按惯例加 \ansspace；选择题括号排 \mbox{（\quad）}；
%   ⑤ 第 12 题定义的运算 $\odot$ 原卷排实心小圆，本稿排 $\odot$。
%
% 行数核对：逐题数 \step/\stepm 条数，合计 82 条（Q3：1；Q4：2；Q5：1；Q6：4；Q7：3；
% Q8：2；Q9：5；Q10：2；Q11：5；Q12：4；Q13：3；Q14：2；Q15：2；Q16：9；Q17：2；
% Q18：11；Q19：6；Q20：11；Q21：7）。全卷 19 题均印【解析】。
% 原卷解析的换行按语意合并，未逐行强拆。
%
\documentclass[a4paper,11pt]{article}
\usepackage{common}
\begin{document}

\examtitlest[7]{2026--2027 年市西中学高一上 10 月月考}{集合与常用逻辑用语、不等式}

\bigsec{一、填空题}

\begin{enumerate}
\setcounter{enumi}{2}

\item 集合 $A=\{x\mid y=\sqrt{1-x^2}\},B=\{y\mid y=x^2+1,x\in A\}$，则 $A\cap\overline B=\blankline[4.4em]$.

\ifanswer
\solbegin
\step{$1-x^2\geqslant0\Rightarrow A=[-1,1]$，从而 $B=[1,2]$，所以 $A\cap\overline B=[-1,1)$.}
\solend
\fi

\item 已知 $b,c\in R$，关于 $x$ 的不等式 $x^2-bx+c<0$ 的解集为 $(-3,2)$，则 $b+c=\blankline[4.4em]$.

\ifanswer
\solbegin
\step{由题意得 $x=-3$ 与 $x=2$ 是 $x^2-bx+c=0$ 的两根，}
\step{所以 $\begin{cases}-3+2=b\\ -3\times2=c\end{cases}$，解得 $\begin{cases}b=-1\\ c=-6\end{cases}$，所以 $b+c=-7$.}
\solend
\fi

\item 若 $\alpha:\lvert x-m\rvert<1$ 是 $\beta:x^2-4x<0$ 的充分条件，则实数 $m$ 的取值范围是\blankline[4.4em].

\ifanswer
\solbegin
\step{$(m-1,m+1)\subseteq(0,4)\Leftrightarrow m-1\geqslant0$ 且 $m+1\leqslant4\Leftrightarrow m\in[1,3]$.}
\solend
\fi

\item 已知方程 $x^2-4x+a=0$ 的两根都大于 $1$，则 $a$ 的取值范围为\blankline[4.4em].

\ifanswer
\solbegin
\step{设函数 $f(x)=x^2-4x+a$，对称轴为 $x=2$，}
\step{因为方程 $x^2-4x+a=0$ 的两根都大于 $1$，}
\step{所以 $\begin{cases}\Delta=(-4)^2-4a\geqslant0\\ f(1)>0\end{cases}$，即 $\begin{cases}16-4a\geqslant0\\ 1-4+a>0\end{cases}$，解得 $3<a\leqslant4$，}
\step{即 $a$ 的取值范围为 $(3,4]$.}
\solend
\fi

\item 若不等式 $(a-2)x^2+2(a-2)x-4<0$ 对 $x\in R$ 恒成立，则实数 $a$ 的取值范围为\blankline[4.4em].

\ifanswer
\solbegin
\step{①当 $a-2=0$，即 $a=2$ 时，不等式为 $-4<0$ 恒成立，符合题意；}
\step{②当 $a\neq2$ 时，则 $\begin{cases}a-2<0\\ \Delta=4(a-2)^2+16(a-2)<0\end{cases}$，解得 $-2<a<2$.}
\step{综上所述，实数 $a$ 的取值范围是 $(-2,2]$.}
\solend
\fi

\item 已知 $-1\leqslant a+b\leqslant4,2\leqslant a-b\leqslant3$，则 $3a-2b$ 的取值范围为\blankline[4.4em].

\ifanswer
\solbegin
\step{设 $3a-2b=m(a+b)+n(a-b)$，则 $m+n=3,m-n=-2$，}
\step{所以 $m=\dfrac12,n=\dfrac52$，所以 $3a-2b=\dfrac12(a+b)+\dfrac52(a-b)\in\left[\dfrac92,\dfrac{19}2\right]$.}
\solend
\fi

\item 设 $x_1$、$x_2$ 是关于 $x$ 的方程 $2x^2-4mx+2m^2+3m-2=0$ 的两个实数根，则 $\lvert x_1-x_2\rvert$ 的最小值为\blankline[4.4em].

\ifanswer
\solbegin
\step{$x_1$、$x_2$ 是方程 $2x^2-4mx+2m^2+3m-2=0$ 的两个实数根，}
\step{所以 $\Delta=(4m)^2-8(2m^2+3m-2)\geqslant0$，解得 $m\leqslant\dfrac23$，}
\step{由韦达定理得 $x_1+x_2=2m,x_1\cdot x_2=\dfrac{2m^2+3m-2}2$，}
\step{则 $\lvert x_1-x_2\rvert=\sqrt{(x_1+x_2)^2-4x_1\cdot x_2}=\sqrt{4-6m}$，}
\step{当 $m=\dfrac23$ 时，$\lvert x_1-x_2\rvert$ 的最小值为 $0$.}
\solend
\fi

\item 设 $x\in R$，$[x]$ 表示不大于 $x$ 的最大整数，如 $[\pi]=3$，$[-1.2]=-2$，$\left[\dfrac12\right]=0$，则使 $\lvert\lvert x^2-1\rvert\rvert=3$ 成立的 $x$ 的取值范围是\blankline[4.4em].

\ifanswer
\solbegin
\step{若 $\lvert\lvert x^2-1\rvert\rvert=3$，则 $3\leqslant\lvert x^2-1\rvert<4$，即 $4\leqslant x^2<5$，}
\step{解得 $2\leqslant x<\sqrt5$ 或 $-\sqrt5<x\leqslant-2$，故 $x$ 的取值范围是 $(-\sqrt5,-2]\cup[2,\sqrt5)$.}
\solend
\fi

\item 设 $a$ 为实数，若关于 $x$ 的一元一次不等式组 $\begin{cases}2x+a>0\\ 3x-6a<0\end{cases}$ 的解集中有且仅有 $4$ 个整数，则 $a$ 的取值范围是\blankline[4.4em].

\ifanswer
\solbegin
\step{一元一次不等式组 $\begin{cases}2x+a>0\\ 3x-6a<0\end{cases}\Leftrightarrow-\dfrac a2<x<2a$，且 $a>0$，}
\step{关于 $x$ 的一元一次不等式组 $\begin{cases}2x+a>0\\ 3x-6a<0\end{cases}$ 的解集中有且仅有 $4$ 个整数，}
\step{①当 $4$ 个整数为 $0$，$1$，$2$，$3$ 时，则 $\begin{cases}-1\leqslant-\dfrac a2<0\\ 3<2a\leqslant4\end{cases}$，所以 $\dfrac32<a\leqslant2$，}
\step{②当 $4$ 个整数为 $-1$，$0$，$1$，$2$ 时，则 $\begin{cases}-2\leqslant-\dfrac a2<-1\\ 2<2a\leqslant3\end{cases}$，无解，}
\step{综上，$a$ 的取值范围是 $\left(\dfrac32,2\right]$.}
\solend
\fi

\item 集合 $M$、$N$、$S$ 都是非空集合，现规定如下运算：

$M\odot N\odot S=\{x\mid x\in(M\cap N)\cup(N\cap S)\cup(S\cap M)$ 且 $x\notin M\cap N\cap S\}$．假设集合 $A=\{x\mid a<x<b\}$，$B=\{x\mid c<x<d\}$，$C=\{x\mid e<x<f\}$，其中实数 $a$、$b$、$c$、$d$、$e$、$f$ 满足：$a<c<e<0<b<d<f$．计算 $A\odot B\odot C=\blankline[4.4em]$.

\ifanswer
\solbegin
\step{因为集合 $A=\{x\mid a<x<b\}$，$B=\{x\mid c<x<d\}$，$C=\{x\mid e<x<f\}$，}
\step{又因为 $a<c<e<0<b<d<f$，}
\step{所以 $A\cap B=\{x\mid c<x<b\}$，$B\cap C=\{x\mid e<x<d\}$，$C\cap A=\{x\mid e<x<b\}$，}
\step{$A\cap B\cap C=\{x\mid e<x<b\}$，故 $A\odot B\odot C=\{x\mid c<x\leqslant e$ 或 $b\leqslant x<d\}$.}
\solend
\fi

\end{enumerate}

\bigsec{二、选择题}

\begin{enumerate}
\setcounter{enumi}{12}

\item “$x\neq1$”是“$x^2-1\neq0$”的\mbox{（\quad）}条件.

\keepwithprev
A. 充分非必要\qquad B. 必要非充分

C. 充要\qquad D. 既非充分也非必要

\ifanswer
\solbegin
\step{$x^2-1\neq0\Leftrightarrow x\neq1$ 且 $x\neq-1\Rightarrow x\neq1$.}
\step{反之不成立，如 $x=-1$ 时，$x\neq1$，但 $x^2-1=0$.}
\step{因此“$x\neq1$”是“$x^2-1\neq0$”的必要非充分条件，故选 $B$.}
\solend
\fi

\item 若 $\dfrac1a<\dfrac1b<0$，则下列结论中不正确的是\mbox{（\quad）}

\keepwithprev
A. $a+b<0$\qquad B. $a^2<b^2$\qquad C. $ab<b^2$\qquad D. $ab<a^2$

\ifanswer
\solbegin
\step{由 $\dfrac1a<\dfrac1b<0$ 得 $b<a<0$，因此 $a+b<0$，且 $a^2<ab<b^2$，}
\step{所以 A、B、C 正确，D 不正确．故选 $D$.}
\solend
\fi

\item 非空集合 $X=\{x\mid a+1\leqslant x\leqslant3a-5\}$，$Y=\{x\mid1\leqslant x\leqslant16\}$，使得 $X\subseteq(X\cap Y)$ 成立的所有 $a$ 的集合是\mbox{（\quad）}

\keepwithprev
A. $\{a\mid3\leqslant a\leqslant7\}$\qquad B. $\{a\mid0\leqslant a\leqslant7\}$\qquad C. $\{a\mid3<a\leqslant7\}$\qquad D. $\{a\mid a\leqslant7\}$

\ifanswer
\solbegin
\step{因为 $X\subseteq(X\cap Y)$ 成立，所以 $X\subseteq Y$ 成立.}
\step{$X\neq\varnothing$，要使 $X\subseteq Y$，则 $\begin{cases}a\geqslant3\\ a+1\geqslant1\\ 3a-5\leqslant16\end{cases}$，即 $\begin{cases}a\geqslant3\\ a\geqslant0\\ a\leqslant7\end{cases}$，所以 $3\leqslant a\leqslant7$，故选 $A$.}
\solend
\fi

\item 已知 $M$ 是满足下列条件的集合：①$0\in M$，$1\in M$；②若 $x$、$y\in M$，则 $x-y\in M$；③若 $x\in M$ 且 $x\neq0$，则 $\dfrac1x\in M$．则下列说法正确的个数为\mbox{（\quad）}

（1）$\dfrac12\in M$；（2）$x+y\in M$；（3）$x^2\in M$．

\keepwithprev
A. $0$\qquad B. $1$\qquad C. $2$\qquad D. $3$

\ifanswer
\solbegin
\step{（1）$\dfrac12\in M$ 正确，理由如下：}
\step{由①得 $0\in M$，$1\in M$，由②得 $0-1=-1\in M$，$1-(-1)=2\in M$，}
\step{由③得 $\dfrac12\in M$；}
\step{（2）由①得 $0\in M$，由题意 $y\in M$，}
\step{所以由②得 $0-y=-y\in M$，又 $x\in M$，所以 $x-(-y)=x+y\in M$；}
\step{（3）$x\in M$，由②得 $x-1\in M$，由③得 $\dfrac1x\in M$，$\dfrac1{x-1}\in M$，}
\step{所以 $\dfrac1x-\dfrac1{x-1}\in M$，即 $\dfrac1{x(x-1)}\in M$，所以 $x(x-1)\in M$，}
\step{由（2）结论得 $x(x-1)+x\in M$，即 $x^2\in M$，}
\step{所以说法正确的个数为 $3$．故选 $D$.}
\solend
\fi

\end{enumerate}

\bigsec{三、解答题}

\begin{enumerate}
\setcounter{enumi}{16}

\item 已知 $a>b$，用比较法证明：$a^3>b^3$.

\ansspace*[8]

\ifanswer
\solbegin
\step{$a>b$，$a^3-b^3=(a-b)(a^2+ab+b^2)=(a-b)\left[\left(a+\dfrac b2\right)^2+\dfrac34b^2\right]>0$，}
\step{故 $a^3>b^3$.}
\solend
\fi

\item 求下列关于 $x$ 的不等式的解集：

（1）$\dfrac{x^2(x+3)(x^2-x+1)}{x^2-5x+4}>0$；

（2）$\dfrac{ax-1}{x-2}\geqslant0\left(a<\dfrac12\right)$.

\ansspace*[12]

\ifanswer
\solbegin
\stepm{（1）}{分母 $x^2-5x+4=(x-1)(x-4)$，要求 $x\neq1,4$；}
\step{又 $x=0$ 时原不等式不成立，且 $x^2-x+1=\left(x-\dfrac12\right)^2+\dfrac34>0$，}
\step{因此原不等式等价于 $\dfrac{x+3}{(x-1)(x-4)}>0,x\neq0$.}
\step{在数轴上标出 $-3,1,4$，分段判号得 $-3<x<1$ 或 $x>4$，}
\step{再排除 $x=0$，解集为 $(-3,0)\cup(0,1)\cup(4,+\infty)$.}
\stepm{（2）}{原不等式要求 $x\neq2$，按 $a$ 的符号分类讨论.}
\step{①当 $a<0$ 时，$\dfrac1a<2$，原不等式等价于 $\dfrac{x-\dfrac1a}{x-2}\leqslant0\Leftrightarrow\dfrac1a\leqslant x<2$，}
\step{解集为 $\left[\dfrac1a,2\right)$.}
\step{②当 $a=0$ 时，原不等式为 $\dfrac{-1}{x-2}\geqslant0\Leftrightarrow x<2$，解集为 $(-\infty,2)$.}
\step{③当 $0<a<\dfrac12$ 时，$\dfrac1a>2$，}
\step{原不等式等价于 $\dfrac{x-\dfrac1a}{x-2}\geqslant0\Leftrightarrow x<2$ 或 $x\geqslant\dfrac1a$，解集为 $(-\infty,2)\cup\left[\dfrac1a,+\infty\right)$.}
\solend
\fi

\item 集合 $A=\{x\mid-2\leqslant x\leqslant5\}$，$B=\{x\mid m+1\leqslant x\leqslant2m-1\}$.

（1）若 $B\subseteq A$，求实数 $m$ 的取值范围；

（2）若 $A\cap B\neq\varnothing$，求实数 $m$ 的取值范围.

\ansspace*[10]

\ifanswer
\solbegin
\stepm{（1）}{当 $B=\varnothing$ 时，$m+1>2m-1$，解得 $m<2$，}
\step{当 $B\neq\varnothing$ 时，$\begin{cases}m+1\leqslant2m-1\\ m+1\geqslant-2\\ 2m-1\leqslant5\end{cases}$，解得 $2\leqslant m\leqslant3$，}
\step{综上，实数 $m$ 的取值范围是 $(-\infty,3]$；}
\stepm{（2）}{当 $B=\varnothing$ 时，$m+1>2m-1$，解得 $m<2$，}
\step{当 $B\neq\varnothing$ 时，$\begin{cases}m+1\leqslant2m-1\\ 2m-1<-2\end{cases}$ 或 $\begin{cases}m+1\leqslant2m-1\\ m+1>5\end{cases}$，解得 $m\in\varnothing$ 或 $m>4$.}
\step{所以当 $A\cap B\neq\varnothing$ 时，实数 $m$ 的取值范围是 $[2,4]$.}
\solend
\fi

\item 在解决实际问题时，往往会有不同的思路和方法，这些方法有些正确，有些错误；有些简洁，有些复杂.

问题①：设 $a\in R$，集合 $M=\{x\mid(x-1)(x-a)=0\}$，若 $x\in M$ 是 $x\in\{1,2,3\}$ 的充分条件，求：$a$ 的取值集合.

问题②：设 $x\in Z$，$a=x+3$，$b=18-x$，若 $c\in Z$，求证：$c-a$ 和 $c-b$ 至少有一个数是奇数

（1）小明在解决问题①，他认为原问题等价于 $\{1,a\}\subseteq\{1,2,3\}$，解得 $a$ 的取值集合为 $\{2,3\}$，张老师判断小明解题错误，请计算正确的 $a$ 的取值集合；

（2）小红认为既然 $c\in Z$，只需根据 $c$ 是奇数还是偶数，分类讨论即可；小华则认为可以使用反证法解决问题，请你选一种你认为更好的方法并证明.

\ansspace*[12]

\ifanswer
\solbegin
\stepm{（1）}{当 $a=1$ 时，集合 $M=\{x\mid(x-1)^2=0\}=\{1\}$，}
\step{此时 $x\in M\Rightarrow x\in\{1,2,3\}$，即 $x\in M$ 是 $x\in\{1,2,3\}$ 的充分条件.}
\step{当 $a\neq1$ 时，集合 $M=\{x\mid(x-1)(x-a)=0\}=\{1,a\}$，}
\step{由 $x\in M$ 是 $x\in\{1,2,3\}$ 的充分条件，得 $\{1,a\}\subseteq\{1,2,3\}$，}
\step{所以 $a=2$ 或 $a=3$.}
\step{综上，实数 $a$ 的取值集合为 $\{1,2,3\}$；}
\stepm{（2）}{利用反证法证明．证明如下：}
\step{假设 $c-a$、$c-b$ 均为偶数，则 $(c-b)-(c-a)=a-b$ 为偶数.}
\step{因为 $a-b=x+3-18+x=2x-15$ 为奇数，}
\step{与 $b-a$ 为偶数相矛盾，所以原假设不成立.}
\step{故 $c-a$、$c-b$ 至少有一个数是奇数.}
\solend
\fi

\item 已知实数 $x_1$、$x_2$、$x_3$、$x_4$、$x_5$，满足 $x_1+x_2+x_3+x_4+x_5=5$.

（1）求证：$x_1$、$x_2$、$x_2$、$x_4$、$x_5$ 中至少有一个实数不小于 $1$；% 原卷把 x_3 误印作 x_2，照录

（2）设 $x_1$、$x_2$、$x_3$、$x_4$、$x_5$ 这五个实数两两不等，集合 $A=\{x_1,x_2,x_3,x_4,x_5\}$，若 $B\neq\varnothing$ 且 $B\subseteq A$，记 $G(B)$ 是 $B$ 中所有元素之和，对所有的 $B$，求 $G(B)$ 的平均值.

\ansspace*[13]

\ifanswer
\solbegin
\stepm{（1）}{假设 $x_1$、$x_2$、$x_2$、$x_4$、$x_5$ 中的每一个数都小于 $1$，}% 原卷 x_3 误印作 x_2，照录
\step{得 $x_1+x_2+x_2+x_4+x_5<5$，这与 $x_1+x_2+x_3+x_4+x_5=5$ 矛盾，}% 前半句原卷 x_3 误印作 x_2，照录
\step{故 $x_1$、$x_2$、$x_2$、$x_4$、$x_5$ 中至少有一个实数不小于 $1$；}% 原卷 x_3 误印作 x_2，照录
\stepm{（2）}{集合 $A=\{x_1,x_2,x_3,x_4,x_5\}$ 的非空子集个数为 $32-1=31$，}
\step{由于 $G(B)$ 是 $B$ 中所有元素之和，}
\step{得 $G(B)=16(x_1+x_2+x_3+x_4+x_5)=16\times5=80$，}
\step{则 $G(B)$ 的平均值为 $\dfrac{80}{31}$.}
\solend
\fi

\end{enumerate}

\end{document}
"
%
% 原始材料：微信公众号“魔都数学加油站”文章，连载“27学年高一98”；
% 原文标题“27学年高一98——2026-2027年上海市市西中学高一上10月月考”；
% 原文链接：https://mp.weixin.qq.com/s/duK4wz-HT3W9tjOeYQaVfg。
% 正文为 8 张 PNG 页（`原图/高一98_市西中学_10月月考/01..08.png`），
% 尺寸混排（第 1—4、6、7 页为 4762×6735，第 5、8 页为 2560×3620），按页序读取。
%
% 篇幅与题号：原卷自第 3 题起编号（第 1、2 题不在原文内）。填空题 3—12（10 题）、
% 选择题 13—16（4 题）、解答题 17—21（5 题），共 19 题。
%
% 讲评版：全部 19 题均印【解析】，其中第 13—16 题的答案在解析末句
% “故选 $B$／$D$／$A$／$D$”（题干括号内未印答案）。全卷无插图。
%
% 副标题：原卷未印副标题。本稿按内容标作“集合与常用逻辑用语、不等式”。
%
% 原卷笔误/存疑（均照录未改，正文对应处另有 % 注）：
%   ① 第 21 题 (1) 求证行、解析“假设”行与“故”行三处把 $x_3$ 误印作 $x_2$
%      （“$x_1$、$x_2$、$x_2$、$x_4$、$x_5$”）；解析同一行后半句
%      “这与 $x_1+x_2+x_3+x_4+x_5=5$ 矛盾”却用 $x_3$，显系笔误，照录；
%   ② 第 6 题解析的判别式行印作“$\Delta=(-4)^2-4a\geqslant0$”，与
%      $f(1)>0$ 并列，照录。
%
% 与原卷的排版差异：
%   ① 原卷每题下直接印【解析】，本稿按全稿惯例将答案与解析置于答案版；
%   ② 原卷未印副标题，本稿补出；填空横线统一排 \blankline[4.4em]；
%   ③ 原卷数集记号排斜体 $R$、$\Z$，本稿按全稿惯例改用 $\R$、$\Z$；
%   ④ 解答题按惯例加 \ansspace；选择题括号排 \mbox{（\quad）}；
%   ⑤ 第 12 题定义的运算 $\odot$ 原卷排实心小圆，本稿排 $\odot$。
%
% 行数核对：逐题数 \step/\stepm 条数，合计 82 条（Q3：1；Q4：2；Q5：1；Q6：4；Q7：3；
% Q8：2；Q9：5；Q10：2；Q11：5；Q12：4；Q13：3；Q14：2；Q15：2；Q16：9；Q17：2；
% Q18：11；Q19：6；Q20：11；Q21：7）。全卷 19 题均印【解析】。
% 原卷解析的换行按语意合并，未逐行强拆。
%
\documentclass[a4paper,11pt]{article}
\usepackage{common}
\begin{document}

\examtitlest[7]{2026--2027 年市西中学高一上 10 月月考}{集合与常用逻辑用语、不等式}

\bigsec{一、填空题}

\begin{enumerate}
\setcounter{enumi}{2}

\item 集合 $A=\{x\mid y=\sqrt{1-x^2}\},B=\{y\mid y=x^2+1,x\in A\}$，则 $A\cap\overline B=\blankline[4.4em]$.

\ifanswer
\solbegin
\step{$1-x^2\geqslant0\Rightarrow A=[-1,1]$，从而 $B=[1,2]$，所以 $A\cap\overline B=[-1,1)$.}
\solend
\fi

\item 已知 $b,c\in R$，关于 $x$ 的不等式 $x^2-bx+c<0$ 的解集为 $(-3,2)$，则 $b+c=\blankline[4.4em]$.

\ifanswer
\solbegin
\step{由题意得 $x=-3$ 与 $x=2$ 是 $x^2-bx+c=0$ 的两根，}
\step{所以 $\begin{cases}-3+2=b\\ -3\times2=c\end{cases}$，解得 $\begin{cases}b=-1\\ c=-6\end{cases}$，所以 $b+c=-7$.}
\solend
\fi

\item 若 $\alpha:\lvert x-m\rvert<1$ 是 $\beta:x^2-4x<0$ 的充分条件，则实数 $m$ 的取值范围是\blankline[4.4em].

\ifanswer
\solbegin
\step{$(m-1,m+1)\subseteq(0,4)\Leftrightarrow m-1\geqslant0$ 且 $m+1\leqslant4\Leftrightarrow m\in[1,3]$.}
\solend
\fi

\item 已知方程 $x^2-4x+a=0$ 的两根都大于 $1$，则 $a$ 的取值范围为\blankline[4.4em].

\ifanswer
\solbegin
\step{设函数 $f(x)=x^2-4x+a$，对称轴为 $x=2$，}
\step{因为方程 $x^2-4x+a=0$ 的两根都大于 $1$，}
\step{所以 $\begin{cases}\Delta=(-4)^2-4a\geqslant0\\ f(1)>0\end{cases}$，即 $\begin{cases}16-4a\geqslant0\\ 1-4+a>0\end{cases}$，解得 $3<a\leqslant4$，}
\step{即 $a$ 的取值范围为 $(3,4]$.}
\solend
\fi

\item 若不等式 $(a-2)x^2+2(a-2)x-4<0$ 对 $x\in R$ 恒成立，则实数 $a$ 的取值范围为\blankline[4.4em].

\ifanswer
\solbegin
\step{①当 $a-2=0$，即 $a=2$ 时，不等式为 $-4<0$ 恒成立，符合题意；}
\step{②当 $a\neq2$ 时，则 $\begin{cases}a-2<0\\ \Delta=4(a-2)^2+16(a-2)<0\end{cases}$，解得 $-2<a<2$.}
\step{综上所述，实数 $a$ 的取值范围是 $(-2,2]$.}
\solend
\fi

\item 已知 $-1\leqslant a+b\leqslant4,2\leqslant a-b\leqslant3$，则 $3a-2b$ 的取值范围为\blankline[4.4em].

\ifanswer
\solbegin
\step{设 $3a-2b=m(a+b)+n(a-b)$，则 $m+n=3,m-n=-2$，}
\step{所以 $m=\dfrac12,n=\dfrac52$，所以 $3a-2b=\dfrac12(a+b)+\dfrac52(a-b)\in\left[\dfrac92,\dfrac{19}2\right]$.}
\solend
\fi

\item 设 $x_1$、$x_2$ 是关于 $x$ 的方程 $2x^2-4mx+2m^2+3m-2=0$ 的两个实数根，则 $\lvert x_1-x_2\rvert$ 的最小值为\blankline[4.4em].

\ifanswer
\solbegin
\step{$x_1$、$x_2$ 是方程 $2x^2-4mx+2m^2+3m-2=0$ 的两个实数根，}
\step{所以 $\Delta=(4m)^2-8(2m^2+3m-2)\geqslant0$，解得 $m\leqslant\dfrac23$，}
\step{由韦达定理得 $x_1+x_2=2m,x_1\cdot x_2=\dfrac{2m^2+3m-2}2$，}
\step{则 $\lvert x_1-x_2\rvert=\sqrt{(x_1+x_2)^2-4x_1\cdot x_2}=\sqrt{4-6m}$，}
\step{当 $m=\dfrac23$ 时，$\lvert x_1-x_2\rvert$ 的最小值为 $0$.}
\solend
\fi

\item 设 $x\in R$，$[x]$ 表示不大于 $x$ 的最大整数，如 $[\pi]=3$，$[-1.2]=-2$，$\left[\dfrac12\right]=0$，则使 $\lvert\lvert x^2-1\rvert\rvert=3$ 成立的 $x$ 的取值范围是\blankline[4.4em].

\ifanswer
\solbegin
\step{若 $\lvert\lvert x^2-1\rvert\rvert=3$，则 $3\leqslant\lvert x^2-1\rvert<4$，即 $4\leqslant x^2<5$，}
\step{解得 $2\leqslant x<\sqrt5$ 或 $-\sqrt5<x\leqslant-2$，故 $x$ 的取值范围是 $(-\sqrt5,-2]\cup[2,\sqrt5)$.}
\solend
\fi

\item 设 $a$ 为实数，若关于 $x$ 的一元一次不等式组 $\begin{cases}2x+a>0\\ 3x-6a<0\end{cases}$ 的解集中有且仅有 $4$ 个整数，则 $a$ 的取值范围是\blankline[4.4em].

\ifanswer
\solbegin
\step{一元一次不等式组 $\begin{cases}2x+a>0\\ 3x-6a<0\end{cases}\Leftrightarrow-\dfrac a2<x<2a$，且 $a>0$，}
\step{关于 $x$ 的一元一次不等式组 $\begin{cases}2x+a>0\\ 3x-6a<0\end{cases}$ 的解集中有且仅有 $4$ 个整数，}
\step{①当 $4$ 个整数为 $0$，$1$，$2$，$3$ 时，则 $\begin{cases}-1\leqslant-\dfrac a2<0\\ 3<2a\leqslant4\end{cases}$，所以 $\dfrac32<a\leqslant2$，}
\step{②当 $4$ 个整数为 $-1$，$0$，$1$，$2$ 时，则 $\begin{cases}-2\leqslant-\dfrac a2<-1\\ 2<2a\leqslant3\end{cases}$，无解，}
\step{综上，$a$ 的取值范围是 $\left(\dfrac32,2\right]$.}
\solend
\fi

\item 集合 $M$、$N$、$S$ 都是非空集合，现规定如下运算：

$M\odot N\odot S=\{x\mid x\in(M\cap N)\cup(N\cap S)\cup(S\cap M)$ 且 $x\notin M\cap N\cap S\}$．假设集合 $A=\{x\mid a<x<b\}$，$B=\{x\mid c<x<d\}$，$C=\{x\mid e<x<f\}$，其中实数 $a$、$b$、$c$、$d$、$e$、$f$ 满足：$a<c<e<0<b<d<f$．计算 $A\odot B\odot C=\blankline[4.4em]$.

\ifanswer
\solbegin
\step{因为集合 $A=\{x\mid a<x<b\}$，$B=\{x\mid c<x<d\}$，$C=\{x\mid e<x<f\}$，}
\step{又因为 $a<c<e<0<b<d<f$，}
\step{所以 $A\cap B=\{x\mid c<x<b\}$，$B\cap C=\{x\mid e<x<d\}$，$C\cap A=\{x\mid e<x<b\}$，}
\step{$A\cap B\cap C=\{x\mid e<x<b\}$，故 $A\odot B\odot C=\{x\mid c<x\leqslant e$ 或 $b\leqslant x<d\}$.}
\solend
\fi

\end{enumerate}

\bigsec{二、选择题}

\begin{enumerate}
\setcounter{enumi}{12}

\item “$x\neq1$”是“$x^2-1\neq0$”的\mbox{（\quad）}条件.

\keepwithprev
A. 充分非必要\qquad B. 必要非充分

C. 充要\qquad D. 既非充分也非必要

\ifanswer
\solbegin
\step{$x^2-1\neq0\Leftrightarrow x\neq1$ 且 $x\neq-1\Rightarrow x\neq1$.}
\step{反之不成立，如 $x=-1$ 时，$x\neq1$，但 $x^2-1=0$.}
\step{因此“$x\neq1$”是“$x^2-1\neq0$”的必要非充分条件，故选 $B$.}
\solend
\fi

\item 若 $\dfrac1a<\dfrac1b<0$，则下列结论中不正确的是\mbox{（\quad）}

\keepwithprev
A. $a+b<0$\qquad B. $a^2<b^2$\qquad C. $ab<b^2$\qquad D. $ab<a^2$

\ifanswer
\solbegin
\step{由 $\dfrac1a<\dfrac1b<0$ 得 $b<a<0$，因此 $a+b<0$，且 $a^2<ab<b^2$，}
\step{所以 A、B、C 正确，D 不正确．故选 $D$.}
\solend
\fi

\item 非空集合 $X=\{x\mid a+1\leqslant x\leqslant3a-5\}$，$Y=\{x\mid1\leqslant x\leqslant16\}$，使得 $X\subseteq(X\cap Y)$ 成立的所有 $a$ 的集合是\mbox{（\quad）}

\keepwithprev
A. $\{a\mid3\leqslant a\leqslant7\}$\qquad B. $\{a\mid0\leqslant a\leqslant7\}$\qquad C. $\{a\mid3<a\leqslant7\}$\qquad D. $\{a\mid a\leqslant7\}$

\ifanswer
\solbegin
\step{因为 $X\subseteq(X\cap Y)$ 成立，所以 $X\subseteq Y$ 成立.}
\step{$X\neq\varnothing$，要使 $X\subseteq Y$，则 $\begin{cases}a\geqslant3\\ a+1\geqslant1\\ 3a-5\leqslant16\end{cases}$，即 $\begin{cases}a\geqslant3\\ a\geqslant0\\ a\leqslant7\end{cases}$，所以 $3\leqslant a\leqslant7$，故选 $A$.}
\solend
\fi

\item 已知 $M$ 是满足下列条件的集合：①$0\in M$，$1\in M$；②若 $x$、$y\in M$，则 $x-y\in M$；③若 $x\in M$ 且 $x\neq0$，则 $\dfrac1x\in M$．则下列说法正确的个数为\mbox{（\quad）}

（1）$\dfrac12\in M$；（2）$x+y\in M$；（3）$x^2\in M$．

\keepwithprev
A. $0$\qquad B. $1$\qquad C. $2$\qquad D. $3$

\ifanswer
\solbegin
\step{（1）$\dfrac12\in M$ 正确，理由如下：}
\step{由①得 $0\in M$，$1\in M$，由②得 $0-1=-1\in M$，$1-(-1)=2\in M$，}
\step{由③得 $\dfrac12\in M$；}
\step{（2）由①得 $0\in M$，由题意 $y\in M$，}
\step{所以由②得 $0-y=-y\in M$，又 $x\in M$，所以 $x-(-y)=x+y\in M$；}
\step{（3）$x\in M$，由②得 $x-1\in M$，由③得 $\dfrac1x\in M$，$\dfrac1{x-1}\in M$，}
\step{所以 $\dfrac1x-\dfrac1{x-1}\in M$，即 $\dfrac1{x(x-1)}\in M$，所以 $x(x-1)\in M$，}
\step{由（2）结论得 $x(x-1)+x\in M$，即 $x^2\in M$，}
\step{所以说法正确的个数为 $3$．故选 $D$.}
\solend
\fi

\end{enumerate}

\bigsec{三、解答题}

\begin{enumerate}
\setcounter{enumi}{16}

\item 已知 $a>b$，用比较法证明：$a^3>b^3$.

\ansspace*[8]

\ifanswer
\solbegin
\step{$a>b$，$a^3-b^3=(a-b)(a^2+ab+b^2)=(a-b)\left[\left(a+\dfrac b2\right)^2+\dfrac34b^2\right]>0$，}
\step{故 $a^3>b^3$.}
\solend
\fi

\item 求下列关于 $x$ 的不等式的解集：

（1）$\dfrac{x^2(x+3)(x^2-x+1)}{x^2-5x+4}>0$；

（2）$\dfrac{ax-1}{x-2}\geqslant0\left(a<\dfrac12\right)$.

\ansspace*[12]

\ifanswer
\solbegin
\stepm{（1）}{分母 $x^2-5x+4=(x-1)(x-4)$，要求 $x\neq1,4$；}
\step{又 $x=0$ 时原不等式不成立，且 $x^2-x+1=\left(x-\dfrac12\right)^2+\dfrac34>0$，}
\step{因此原不等式等价于 $\dfrac{x+3}{(x-1)(x-4)}>0,x\neq0$.}
\step{在数轴上标出 $-3,1,4$，分段判号得 $-3<x<1$ 或 $x>4$，}
\step{再排除 $x=0$，解集为 $(-3,0)\cup(0,1)\cup(4,+\infty)$.}
\stepm{（2）}{原不等式要求 $x\neq2$，按 $a$ 的符号分类讨论.}
\step{①当 $a<0$ 时，$\dfrac1a<2$，原不等式等价于 $\dfrac{x-\dfrac1a}{x-2}\leqslant0\Leftrightarrow\dfrac1a\leqslant x<2$，}
\step{解集为 $\left[\dfrac1a,2\right)$.}
\step{②当 $a=0$ 时，原不等式为 $\dfrac{-1}{x-2}\geqslant0\Leftrightarrow x<2$，解集为 $(-\infty,2)$.}
\step{③当 $0<a<\dfrac12$ 时，$\dfrac1a>2$，}
\step{原不等式等价于 $\dfrac{x-\dfrac1a}{x-2}\geqslant0\Leftrightarrow x<2$ 或 $x\geqslant\dfrac1a$，解集为 $(-\infty,2)\cup\left[\dfrac1a,+\infty\right)$.}
\solend
\fi

\item 集合 $A=\{x\mid-2\leqslant x\leqslant5\}$，$B=\{x\mid m+1\leqslant x\leqslant2m-1\}$.

（1）若 $B\subseteq A$，求实数 $m$ 的取值范围；

（2）若 $A\cap B\neq\varnothing$，求实数 $m$ 的取值范围.

\ansspace*[10]

\ifanswer
\solbegin
\stepm{（1）}{当 $B=\varnothing$ 时，$m+1>2m-1$，解得 $m<2$，}
\step{当 $B\neq\varnothing$ 时，$\begin{cases}m+1\leqslant2m-1\\ m+1\geqslant-2\\ 2m-1\leqslant5\end{cases}$，解得 $2\leqslant m\leqslant3$，}
\step{综上，实数 $m$ 的取值范围是 $(-\infty,3]$；}
\stepm{（2）}{当 $B=\varnothing$ 时，$m+1>2m-1$，解得 $m<2$，}
\step{当 $B\neq\varnothing$ 时，$\begin{cases}m+1\leqslant2m-1\\ 2m-1<-2\end{cases}$ 或 $\begin{cases}m+1\leqslant2m-1\\ m+1>5\end{cases}$，解得 $m\in\varnothing$ 或 $m>4$.}
\step{所以当 $A\cap B\neq\varnothing$ 时，实数 $m$ 的取值范围是 $[2,4]$.}
\solend
\fi

\item 在解决实际问题时，往往会有不同的思路和方法，这些方法有些正确，有些错误；有些简洁，有些复杂.

问题①：设 $a\in R$，集合 $M=\{x\mid(x-1)(x-a)=0\}$，若 $x\in M$ 是 $x\in\{1,2,3\}$ 的充分条件，求：$a$ 的取值集合.

问题②：设 $x\in Z$，$a=x+3$，$b=18-x$，若 $c\in Z$，求证：$c-a$ 和 $c-b$ 至少有一个数是奇数

（1）小明在解决问题①，他认为原问题等价于 $\{1,a\}\subseteq\{1,2,3\}$，解得 $a$ 的取值集合为 $\{2,3\}$，张老师判断小明解题错误，请计算正确的 $a$ 的取值集合；

（2）小红认为既然 $c\in Z$，只需根据 $c$ 是奇数还是偶数，分类讨论即可；小华则认为可以使用反证法解决问题，请你选一种你认为更好的方法并证明.

\ansspace*[12]

\ifanswer
\solbegin
\stepm{（1）}{当 $a=1$ 时，集合 $M=\{x\mid(x-1)^2=0\}=\{1\}$，}
\step{此时 $x\in M\Rightarrow x\in\{1,2,3\}$，即 $x\in M$ 是 $x\in\{1,2,3\}$ 的充分条件.}
\step{当 $a\neq1$ 时，集合 $M=\{x\mid(x-1)(x-a)=0\}=\{1,a\}$，}
\step{由 $x\in M$ 是 $x\in\{1,2,3\}$ 的充分条件，得 $\{1,a\}\subseteq\{1,2,3\}$，}
\step{所以 $a=2$ 或 $a=3$.}
\step{综上，实数 $a$ 的取值集合为 $\{1,2,3\}$；}
\stepm{（2）}{利用反证法证明．证明如下：}
\step{假设 $c-a$、$c-b$ 均为偶数，则 $(c-b)-(c-a)=a-b$ 为偶数.}
\step{因为 $a-b=x+3-18+x=2x-15$ 为奇数，}
\step{与 $b-a$ 为偶数相矛盾，所以原假设不成立.}
\step{故 $c-a$、$c-b$ 至少有一个数是奇数.}
\solend
\fi

\item 已知实数 $x_1$、$x_2$、$x_3$、$x_4$、$x_5$，满足 $x_1+x_2+x_3+x_4+x_5=5$.

（1）求证：$x_1$、$x_2$、$x_2$、$x_4$、$x_5$ 中至少有一个实数不小于 $1$；% 原卷把 x_3 误印作 x_2，照录

（2）设 $x_1$、$x_2$、$x_3$、$x_4$、$x_5$ 这五个实数两两不等，集合 $A=\{x_1,x_2,x_3,x_4,x_5\}$，若 $B\neq\varnothing$ 且 $B\subseteq A$，记 $G(B)$ 是 $B$ 中所有元素之和，对所有的 $B$，求 $G(B)$ 的平均值.

\ansspace*[13]

\ifanswer
\solbegin
\stepm{（1）}{假设 $x_1$、$x_2$、$x_2$、$x_4$、$x_5$ 中的每一个数都小于 $1$，}% 原卷 x_3 误印作 x_2，照录
\step{得 $x_1+x_2+x_2+x_4+x_5<5$，这与 $x_1+x_2+x_3+x_4+x_5=5$ 矛盾，}% 前半句原卷 x_3 误印作 x_2，照录
\step{故 $x_1$、$x_2$、$x_2$、$x_4$、$x_5$ 中至少有一个实数不小于 $1$；}% 原卷 x_3 误印作 x_2，照录
\stepm{（2）}{集合 $A=\{x_1,x_2,x_3,x_4,x_5\}$ 的非空子集个数为 $32-1=31$，}
\step{由于 $G(B)$ 是 $B$ 中所有元素之和，}
\step{得 $G(B)=16(x_1+x_2+x_3+x_4+x_5)=16\times5=80$，}
\step{则 $G(B)$ 的平均值为 $\dfrac{80}{31}$.}
\solend
\fi

\end{enumerate}

\end{document}
"
%
% 原始材料：微信公众号“魔都数学加油站”文章，连载“27学年高一98”；
% 原文标题“27学年高一98——2026-2027年上海市市西中学高一上10月月考”；
% 原文链接：https://mp.weixin.qq.com/s/duK4wz-HT3W9tjOeYQaVfg。
% 正文为 8 张 PNG 页（`原图/高一98_市西中学_10月月考/01..08.png`），
% 尺寸混排（第 1—4、6、7 页为 4762×6735，第 5、8 页为 2560×3620），按页序读取。
%
% 篇幅与题号：原卷自第 3 题起编号（第 1、2 题不在原文内）。填空题 3—12（10 题）、
% 选择题 13—16（4 题）、解答题 17—21（5 题），共 19 题。
%
% 讲评版：全部 19 题均印【解析】，其中第 13—16 题的答案在解析末句
% “故选 $B$／$D$／$A$／$D$”（题干括号内未印答案）。全卷无插图。
%
% 副标题：原卷未印副标题。本稿按内容标作“集合与常用逻辑用语、不等式”。
%
% 原卷笔误/存疑（均照录未改，正文对应处另有 % 注）：
%   ① 第 21 题 (1) 求证行、解析“假设”行与“故”行三处把 $x_3$ 误印作 $x_2$
%      （“$x_1$、$x_2$、$x_2$、$x_4$、$x_5$”）；解析同一行后半句
%      “这与 $x_1+x_2+x_3+x_4+x_5=5$ 矛盾”却用 $x_3$，显系笔误，照录；
%   ② 第 6 题解析的判别式行印作“$\Delta=(-4)^2-4a\geqslant0$”，与
%      $f(1)>0$ 并列，照录。
%
% 与原卷的排版差异：
%   ① 原卷每题下直接印【解析】，本稿按全稿惯例将答案与解析置于答案版；
%   ② 原卷未印副标题，本稿补出；填空横线统一排 \blankline[4.4em]；
%   ③ 原卷数集记号排斜体 $R$、$\Z$，本稿按全稿惯例改用 $\R$、$\Z$；
%   ④ 解答题按惯例加 \ansspace；选择题括号排 \mbox{（\quad）}；
%   ⑤ 第 12 题定义的运算 $\odot$ 原卷排实心小圆，本稿排 $\odot$。
%
% 行数核对：逐题数 \step/\stepm 条数，合计 82 条（Q3：1；Q4：2；Q5：1；Q6：4；Q7：3；
% Q8：2；Q9：5；Q10：2；Q11：5；Q12：4；Q13：3；Q14：2；Q15：2；Q16：9；Q17：2；
% Q18：11；Q19：6；Q20：11；Q21：7）。全卷 19 题均印【解析】。
% 原卷解析的换行按语意合并，未逐行强拆。
%
\documentclass[a4paper,11pt]{article}
\usepackage{common}
\begin{document}

\examtitlest[7]{2026--2027 年市西中学高一上 10 月月考}{集合与常用逻辑用语、不等式}

\bigsec{一、填空题}

\begin{enumerate}
\setcounter{enumi}{2}

\item 集合 $A=\{x\mid y=\sqrt{1-x^2}\},B=\{y\mid y=x^2+1,x\in A\}$，则 $A\cap\overline B=\blankline[4.4em]$.

\ifanswer
\solbegin
\step{$1-x^2\geqslant0\Rightarrow A=[-1,1]$，从而 $B=[1,2]$，所以 $A\cap\overline B=[-1,1)$.}
\solend
\fi

\item 已知 $b,c\in R$，关于 $x$ 的不等式 $x^2-bx+c<0$ 的解集为 $(-3,2)$，则 $b+c=\blankline[4.4em]$.

\ifanswer
\solbegin
\step{由题意得 $x=-3$ 与 $x=2$ 是 $x^2-bx+c=0$ 的两根，}
\step{所以 $\begin{cases}-3+2=b\\ -3\times2=c\end{cases}$，解得 $\begin{cases}b=-1\\ c=-6\end{cases}$，所以 $b+c=-7$.}
\solend
\fi

\item 若 $\alpha:\lvert x-m\rvert<1$ 是 $\beta:x^2-4x<0$ 的充分条件，则实数 $m$ 的取值范围是\blankline[4.4em].

\ifanswer
\solbegin
\step{$(m-1,m+1)\subseteq(0,4)\Leftrightarrow m-1\geqslant0$ 且 $m+1\leqslant4\Leftrightarrow m\in[1,3]$.}
\solend
\fi

\item 已知方程 $x^2-4x+a=0$ 的两根都大于 $1$，则 $a$ 的取值范围为\blankline[4.4em].

\ifanswer
\solbegin
\step{设函数 $f(x)=x^2-4x+a$，对称轴为 $x=2$，}
\step{因为方程 $x^2-4x+a=0$ 的两根都大于 $1$，}
\step{所以 $\begin{cases}\Delta=(-4)^2-4a\geqslant0\\ f(1)>0\end{cases}$，即 $\begin{cases}16-4a\geqslant0\\ 1-4+a>0\end{cases}$，解得 $3<a\leqslant4$，}
\step{即 $a$ 的取值范围为 $(3,4]$.}
\solend
\fi

\item 若不等式 $(a-2)x^2+2(a-2)x-4<0$ 对 $x\in R$ 恒成立，则实数 $a$ 的取值范围为\blankline[4.4em].

\ifanswer
\solbegin
\step{①当 $a-2=0$，即 $a=2$ 时，不等式为 $-4<0$ 恒成立，符合题意；}
\step{②当 $a\neq2$ 时，则 $\begin{cases}a-2<0\\ \Delta=4(a-2)^2+16(a-2)<0\end{cases}$，解得 $-2<a<2$.}
\step{综上所述，实数 $a$ 的取值范围是 $(-2,2]$.}
\solend
\fi

\item 已知 $-1\leqslant a+b\leqslant4,2\leqslant a-b\leqslant3$，则 $3a-2b$ 的取值范围为\blankline[4.4em].

\ifanswer
\solbegin
\step{设 $3a-2b=m(a+b)+n(a-b)$，则 $m+n=3,m-n=-2$，}
\step{所以 $m=\dfrac12,n=\dfrac52$，所以 $3a-2b=\dfrac12(a+b)+\dfrac52(a-b)\in\left[\dfrac92,\dfrac{19}2\right]$.}
\solend
\fi

\item 设 $x_1$、$x_2$ 是关于 $x$ 的方程 $2x^2-4mx+2m^2+3m-2=0$ 的两个实数根，则 $\lvert x_1-x_2\rvert$ 的最小值为\blankline[4.4em].

\ifanswer
\solbegin
\step{$x_1$、$x_2$ 是方程 $2x^2-4mx+2m^2+3m-2=0$ 的两个实数根，}
\step{所以 $\Delta=(4m)^2-8(2m^2+3m-2)\geqslant0$，解得 $m\leqslant\dfrac23$，}
\step{由韦达定理得 $x_1+x_2=2m,x_1\cdot x_2=\dfrac{2m^2+3m-2}2$，}
\step{则 $\lvert x_1-x_2\rvert=\sqrt{(x_1+x_2)^2-4x_1\cdot x_2}=\sqrt{4-6m}$，}
\step{当 $m=\dfrac23$ 时，$\lvert x_1-x_2\rvert$ 的最小值为 $0$.}
\solend
\fi

\item 设 $x\in R$，$[x]$ 表示不大于 $x$ 的最大整数，如 $[\pi]=3$，$[-1.2]=-2$，$\left[\dfrac12\right]=0$，则使 $\lvert\lvert x^2-1\rvert\rvert=3$ 成立的 $x$ 的取值范围是\blankline[4.4em].

\ifanswer
\solbegin
\step{若 $\lvert\lvert x^2-1\rvert\rvert=3$，则 $3\leqslant\lvert x^2-1\rvert<4$，即 $4\leqslant x^2<5$，}
\step{解得 $2\leqslant x<\sqrt5$ 或 $-\sqrt5<x\leqslant-2$，故 $x$ 的取值范围是 $(-\sqrt5,-2]\cup[2,\sqrt5)$.}
\solend
\fi

\item 设 $a$ 为实数，若关于 $x$ 的一元一次不等式组 $\begin{cases}2x+a>0\\ 3x-6a<0\end{cases}$ 的解集中有且仅有 $4$ 个整数，则 $a$ 的取值范围是\blankline[4.4em].

\ifanswer
\solbegin
\step{一元一次不等式组 $\begin{cases}2x+a>0\\ 3x-6a<0\end{cases}\Leftrightarrow-\dfrac a2<x<2a$，且 $a>0$，}
\step{关于 $x$ 的一元一次不等式组 $\begin{cases}2x+a>0\\ 3x-6a<0\end{cases}$ 的解集中有且仅有 $4$ 个整数，}
\step{①当 $4$ 个整数为 $0$，$1$，$2$，$3$ 时，则 $\begin{cases}-1\leqslant-\dfrac a2<0\\ 3<2a\leqslant4\end{cases}$，所以 $\dfrac32<a\leqslant2$，}
\step{②当 $4$ 个整数为 $-1$，$0$，$1$，$2$ 时，则 $\begin{cases}-2\leqslant-\dfrac a2<-1\\ 2<2a\leqslant3\end{cases}$，无解，}
\step{综上，$a$ 的取值范围是 $\left(\dfrac32,2\right]$.}
\solend
\fi

\item 集合 $M$、$N$、$S$ 都是非空集合，现规定如下运算：

$M\odot N\odot S=\{x\mid x\in(M\cap N)\cup(N\cap S)\cup(S\cap M)$ 且 $x\notin M\cap N\cap S\}$．假设集合 $A=\{x\mid a<x<b\}$，$B=\{x\mid c<x<d\}$，$C=\{x\mid e<x<f\}$，其中实数 $a$、$b$、$c$、$d$、$e$、$f$ 满足：$a<c<e<0<b<d<f$．计算 $A\odot B\odot C=\blankline[4.4em]$.

\ifanswer
\solbegin
\step{因为集合 $A=\{x\mid a<x<b\}$，$B=\{x\mid c<x<d\}$，$C=\{x\mid e<x<f\}$，}
\step{又因为 $a<c<e<0<b<d<f$，}
\step{所以 $A\cap B=\{x\mid c<x<b\}$，$B\cap C=\{x\mid e<x<d\}$，$C\cap A=\{x\mid e<x<b\}$，}
\step{$A\cap B\cap C=\{x\mid e<x<b\}$，故 $A\odot B\odot C=\{x\mid c<x\leqslant e$ 或 $b\leqslant x<d\}$.}
\solend
\fi

\end{enumerate}

\bigsec{二、选择题}

\begin{enumerate}
\setcounter{enumi}{12}

\item “$x\neq1$”是“$x^2-1\neq0$”的\mbox{（\quad）}条件.

\keepwithprev
A. 充分非必要\qquad B. 必要非充分

C. 充要\qquad D. 既非充分也非必要

\ifanswer
\solbegin
\step{$x^2-1\neq0\Leftrightarrow x\neq1$ 且 $x\neq-1\Rightarrow x\neq1$.}
\step{反之不成立，如 $x=-1$ 时，$x\neq1$，但 $x^2-1=0$.}
\step{因此“$x\neq1$”是“$x^2-1\neq0$”的必要非充分条件，故选 $B$.}
\solend
\fi

\item 若 $\dfrac1a<\dfrac1b<0$，则下列结论中不正确的是\mbox{（\quad）}

\keepwithprev
A. $a+b<0$\qquad B. $a^2<b^2$\qquad C. $ab<b^2$\qquad D. $ab<a^2$

\ifanswer
\solbegin
\step{由 $\dfrac1a<\dfrac1b<0$ 得 $b<a<0$，因此 $a+b<0$，且 $a^2<ab<b^2$，}
\step{所以 A、B、C 正确，D 不正确．故选 $D$.}
\solend
\fi

\item 非空集合 $X=\{x\mid a+1\leqslant x\leqslant3a-5\}$，$Y=\{x\mid1\leqslant x\leqslant16\}$，使得 $X\subseteq(X\cap Y)$ 成立的所有 $a$ 的集合是\mbox{（\quad）}

\keepwithprev
A. $\{a\mid3\leqslant a\leqslant7\}$\qquad B. $\{a\mid0\leqslant a\leqslant7\}$\qquad C. $\{a\mid3<a\leqslant7\}$\qquad D. $\{a\mid a\leqslant7\}$

\ifanswer
\solbegin
\step{因为 $X\subseteq(X\cap Y)$ 成立，所以 $X\subseteq Y$ 成立.}
\step{$X\neq\varnothing$，要使 $X\subseteq Y$，则 $\begin{cases}a\geqslant3\\ a+1\geqslant1\\ 3a-5\leqslant16\end{cases}$，即 $\begin{cases}a\geqslant3\\ a\geqslant0\\ a\leqslant7\end{cases}$，所以 $3\leqslant a\leqslant7$，故选 $A$.}
\solend
\fi

\item 已知 $M$ 是满足下列条件的集合：①$0\in M$，$1\in M$；②若 $x$、$y\in M$，则 $x-y\in M$；③若 $x\in M$ 且 $x\neq0$，则 $\dfrac1x\in M$．则下列说法正确的个数为\mbox{（\quad）}

（1）$\dfrac12\in M$；（2）$x+y\in M$；（3）$x^2\in M$．

\keepwithprev
A. $0$\qquad B. $1$\qquad C. $2$\qquad D. $3$

\ifanswer
\solbegin
\step{（1）$\dfrac12\in M$ 正确，理由如下：}
\step{由①得 $0\in M$，$1\in M$，由②得 $0-1=-1\in M$，$1-(-1)=2\in M$，}
\step{由③得 $\dfrac12\in M$；}
\step{（2）由①得 $0\in M$，由题意 $y\in M$，}
\step{所以由②得 $0-y=-y\in M$，又 $x\in M$，所以 $x-(-y)=x+y\in M$；}
\step{（3）$x\in M$，由②得 $x-1\in M$，由③得 $\dfrac1x\in M$，$\dfrac1{x-1}\in M$，}
\step{所以 $\dfrac1x-\dfrac1{x-1}\in M$，即 $\dfrac1{x(x-1)}\in M$，所以 $x(x-1)\in M$，}
\step{由（2）结论得 $x(x-1)+x\in M$，即 $x^2\in M$，}
\step{所以说法正确的个数为 $3$．故选 $D$.}
\solend
\fi

\end{enumerate}

\bigsec{三、解答题}

\begin{enumerate}
\setcounter{enumi}{16}

\item 已知 $a>b$，用比较法证明：$a^3>b^3$.

\ansspace*[8]

\ifanswer
\solbegin
\step{$a>b$，$a^3-b^3=(a-b)(a^2+ab+b^2)=(a-b)\left[\left(a+\dfrac b2\right)^2+\dfrac34b^2\right]>0$，}
\step{故 $a^3>b^3$.}
\solend
\fi

\item 求下列关于 $x$ 的不等式的解集：

（1）$\dfrac{x^2(x+3)(x^2-x+1)}{x^2-5x+4}>0$；

（2）$\dfrac{ax-1}{x-2}\geqslant0\left(a<\dfrac12\right)$.

\ansspace*[12]

\ifanswer
\solbegin
\stepm{（1）}{分母 $x^2-5x+4=(x-1)(x-4)$，要求 $x\neq1,4$；}
\step{又 $x=0$ 时原不等式不成立，且 $x^2-x+1=\left(x-\dfrac12\right)^2+\dfrac34>0$，}
\step{因此原不等式等价于 $\dfrac{x+3}{(x-1)(x-4)}>0,x\neq0$.}
\step{在数轴上标出 $-3,1,4$，分段判号得 $-3<x<1$ 或 $x>4$，}
\step{再排除 $x=0$，解集为 $(-3,0)\cup(0,1)\cup(4,+\infty)$.}
\stepm{（2）}{原不等式要求 $x\neq2$，按 $a$ 的符号分类讨论.}
\step{①当 $a<0$ 时，$\dfrac1a<2$，原不等式等价于 $\dfrac{x-\dfrac1a}{x-2}\leqslant0\Leftrightarrow\dfrac1a\leqslant x<2$，}
\step{解集为 $\left[\dfrac1a,2\right)$.}
\step{②当 $a=0$ 时，原不等式为 $\dfrac{-1}{x-2}\geqslant0\Leftrightarrow x<2$，解集为 $(-\infty,2)$.}
\step{③当 $0<a<\dfrac12$ 时，$\dfrac1a>2$，}
\step{原不等式等价于 $\dfrac{x-\dfrac1a}{x-2}\geqslant0\Leftrightarrow x<2$ 或 $x\geqslant\dfrac1a$，解集为 $(-\infty,2)\cup\left[\dfrac1a,+\infty\right)$.}
\solend
\fi

\item 集合 $A=\{x\mid-2\leqslant x\leqslant5\}$，$B=\{x\mid m+1\leqslant x\leqslant2m-1\}$.

（1）若 $B\subseteq A$，求实数 $m$ 的取值范围；

（2）若 $A\cap B\neq\varnothing$，求实数 $m$ 的取值范围.

\ansspace*[10]

\ifanswer
\solbegin
\stepm{（1）}{当 $B=\varnothing$ 时，$m+1>2m-1$，解得 $m<2$，}
\step{当 $B\neq\varnothing$ 时，$\begin{cases}m+1\leqslant2m-1\\ m+1\geqslant-2\\ 2m-1\leqslant5\end{cases}$，解得 $2\leqslant m\leqslant3$，}
\step{综上，实数 $m$ 的取值范围是 $(-\infty,3]$；}
\stepm{（2）}{当 $B=\varnothing$ 时，$m+1>2m-1$，解得 $m<2$，}
\step{当 $B\neq\varnothing$ 时，$\begin{cases}m+1\leqslant2m-1\\ 2m-1<-2\end{cases}$ 或 $\begin{cases}m+1\leqslant2m-1\\ m+1>5\end{cases}$，解得 $m\in\varnothing$ 或 $m>4$.}
\step{所以当 $A\cap B\neq\varnothing$ 时，实数 $m$ 的取值范围是 $[2,4]$.}
\solend
\fi

\item 在解决实际问题时，往往会有不同的思路和方法，这些方法有些正确，有些错误；有些简洁，有些复杂.

问题①：设 $a\in R$，集合 $M=\{x\mid(x-1)(x-a)=0\}$，若 $x\in M$ 是 $x\in\{1,2,3\}$ 的充分条件，求：$a$ 的取值集合.

问题②：设 $x\in Z$，$a=x+3$，$b=18-x$，若 $c\in Z$，求证：$c-a$ 和 $c-b$ 至少有一个数是奇数

（1）小明在解决问题①，他认为原问题等价于 $\{1,a\}\subseteq\{1,2,3\}$，解得 $a$ 的取值集合为 $\{2,3\}$，张老师判断小明解题错误，请计算正确的 $a$ 的取值集合；

（2）小红认为既然 $c\in Z$，只需根据 $c$ 是奇数还是偶数，分类讨论即可；小华则认为可以使用反证法解决问题，请你选一种你认为更好的方法并证明.

\ansspace*[12]

\ifanswer
\solbegin
\stepm{（1）}{当 $a=1$ 时，集合 $M=\{x\mid(x-1)^2=0\}=\{1\}$，}
\step{此时 $x\in M\Rightarrow x\in\{1,2,3\}$，即 $x\in M$ 是 $x\in\{1,2,3\}$ 的充分条件.}
\step{当 $a\neq1$ 时，集合 $M=\{x\mid(x-1)(x-a)=0\}=\{1,a\}$，}
\step{由 $x\in M$ 是 $x\in\{1,2,3\}$ 的充分条件，得 $\{1,a\}\subseteq\{1,2,3\}$，}
\step{所以 $a=2$ 或 $a=3$.}
\step{综上，实数 $a$ 的取值集合为 $\{1,2,3\}$；}
\stepm{（2）}{利用反证法证明．证明如下：}
\step{假设 $c-a$、$c-b$ 均为偶数，则 $(c-b)-(c-a)=a-b$ 为偶数.}
\step{因为 $a-b=x+3-18+x=2x-15$ 为奇数，}
\step{与 $b-a$ 为偶数相矛盾，所以原假设不成立.}
\step{故 $c-a$、$c-b$ 至少有一个数是奇数.}
\solend
\fi

\item 已知实数 $x_1$、$x_2$、$x_3$、$x_4$、$x_5$，满足 $x_1+x_2+x_3+x_4+x_5=5$.

（1）求证：$x_1$、$x_2$、$x_2$、$x_4$、$x_5$ 中至少有一个实数不小于 $1$；% 原卷把 x_3 误印作 x_2，照录

（2）设 $x_1$、$x_2$、$x_3$、$x_4$、$x_5$ 这五个实数两两不等，集合 $A=\{x_1,x_2,x_3,x_4,x_5\}$，若 $B\neq\varnothing$ 且 $B\subseteq A$，记 $G(B)$ 是 $B$ 中所有元素之和，对所有的 $B$，求 $G(B)$ 的平均值.

\ansspace*[13]

\ifanswer
\solbegin
\stepm{（1）}{假设 $x_1$、$x_2$、$x_2$、$x_4$、$x_5$ 中的每一个数都小于 $1$，}% 原卷 x_3 误印作 x_2，照录
\step{得 $x_1+x_2+x_2+x_4+x_5<5$，这与 $x_1+x_2+x_3+x_4+x_5=5$ 矛盾，}% 前半句原卷 x_3 误印作 x_2，照录
\step{故 $x_1$、$x_2$、$x_2$、$x_4$、$x_5$ 中至少有一个实数不小于 $1$；}% 原卷 x_3 误印作 x_2，照录
\stepm{（2）}{集合 $A=\{x_1,x_2,x_3,x_4,x_5\}$ 的非空子集个数为 $32-1=31$，}
\step{由于 $G(B)$ 是 $B$ 中所有元素之和，}
\step{得 $G(B)=16(x_1+x_2+x_3+x_4+x_5)=16\times5=80$，}
\step{则 $G(B)$ 的平均值为 $\dfrac{80}{31}$.}
\solend
\fi

\end{enumerate}

\end{document}
"
%
% 原始材料：微信公众号“魔都数学加油站”文章，连载“27学年高一98”；
% 原文标题“27学年高一98——2026-2027年上海市市西中学高一上10月月考”；
% 原文链接：https://mp.weixin.qq.com/s/duK4wz-HT3W9tjOeYQaVfg。
% 正文为 8 张 PNG 页（`原图/高一98_市西中学_10月月考/01..08.png`），
% 尺寸混排（第 1—4、6、7 页为 4762×6735，第 5、8 页为 2560×3620），按页序读取。
%
% 篇幅与题号：原卷自第 3 题起编号（第 1、2 题不在原文内）。填空题 3—12（10 题）、
% 选择题 13—16（4 题）、解答题 17—21（5 题），共 19 题。
%
% 讲评版：全部 19 题均印【解析】，其中第 13—16 题的答案在解析末句
% “故选 $B$／$D$／$A$／$D$”（题干括号内未印答案）。全卷无插图。
%
% 副标题：原卷未印副标题。本稿按内容标作“集合与常用逻辑用语、不等式”。
%
% 原卷笔误/存疑（均照录未改，正文对应处另有 % 注）：
%   ① 第 21 题 (1) 求证行、解析“假设”行与“故”行三处把 $x_3$ 误印作 $x_2$
%      （“$x_1$、$x_2$、$x_2$、$x_4$、$x_5$”）；解析同一行后半句
%      “这与 $x_1+x_2+x_3+x_4+x_5=5$ 矛盾”却用 $x_3$，显系笔误，照录；
%   ② 第 6 题解析的判别式行印作“$\Delta=(-4)^2-4a\geqslant0$”，与
%      $f(1)>0$ 并列，照录。
%
% 与原卷的排版差异：
%   ① 原卷每题下直接印【解析】，本稿按全稿惯例将答案与解析置于答案版；
%   ② 原卷未印副标题，本稿补出；填空横线统一排 \blankline[4.4em]；
%   ③ 原卷数集记号排斜体 $R$、$\Z$，本稿按全稿惯例改用 $\R$、$\Z$；
%   ④ 解答题按惯例加 \ansspace；选择题括号排 \mbox{（\quad）}；
%   ⑤ 第 12 题定义的运算 $\odot$ 原卷排实心小圆，本稿排 $\odot$。
%
% 行数核对：逐题数 \step/\stepm 条数，合计 82 条（Q3：1；Q4：2；Q5：1；Q6：4；Q7：3；
% Q8：2；Q9：5；Q10：2；Q11：5；Q12：4；Q13：3；Q14：2；Q15：2；Q16：9；Q17：2；
% Q18：11；Q19：6；Q20：11；Q21：7）。全卷 19 题均印【解析】。
% 原卷解析的换行按语意合并，未逐行强拆。
%
\documentclass[a4paper,11pt]{article}
\usepackage{common}
\begin{document}

\examtitlest[7]{2026--2027 年市西中学高一上 10 月月考}{集合与常用逻辑用语、不等式}

\bigsec{一、填空题}

\begin{enumerate}
\setcounter{enumi}{2}

\item 集合 $A=\{x\mid y=\sqrt{1-x^2}\},B=\{y\mid y=x^2+1,x\in A\}$，则 $A\cap\overline B=\blankline[4.4em]$.

\ifanswer
\solbegin
\step{$1-x^2\geqslant0\Rightarrow A=[-1,1]$，从而 $B=[1,2]$，所以 $A\cap\overline B=[-1,1)$.}
\solend
\fi

\item 已知 $b,c\in R$，关于 $x$ 的不等式 $x^2-bx+c<0$ 的解集为 $(-3,2)$，则 $b+c=\blankline[4.4em]$.

\ifanswer
\solbegin
\step{由题意得 $x=-3$ 与 $x=2$ 是 $x^2-bx+c=0$ 的两根，}
\step{所以 $\begin{cases}-3+2=b\\ -3\times2=c\end{cases}$，解得 $\begin{cases}b=-1\\ c=-6\end{cases}$，所以 $b+c=-7$.}
\solend
\fi

\item 若 $\alpha:\lvert x-m\rvert<1$ 是 $\beta:x^2-4x<0$ 的充分条件，则实数 $m$ 的取值范围是\blankline[4.4em].

\ifanswer
\solbegin
\step{$(m-1,m+1)\subseteq(0,4)\Leftrightarrow m-1\geqslant0$ 且 $m+1\leqslant4\Leftrightarrow m\in[1,3]$.}
\solend
\fi

\item 已知方程 $x^2-4x+a=0$ 的两根都大于 $1$，则 $a$ 的取值范围为\blankline[4.4em].

\ifanswer
\solbegin
\step{设函数 $f(x)=x^2-4x+a$，对称轴为 $x=2$，}
\step{因为方程 $x^2-4x+a=0$ 的两根都大于 $1$，}
\step{所以 $\begin{cases}\Delta=(-4)^2-4a\geqslant0\\ f(1)>0\end{cases}$，即 $\begin{cases}16-4a\geqslant0\\ 1-4+a>0\end{cases}$，解得 $3<a\leqslant4$，}
\step{即 $a$ 的取值范围为 $(3,4]$.}
\solend
\fi

\item 若不等式 $(a-2)x^2+2(a-2)x-4<0$ 对 $x\in R$ 恒成立，则实数 $a$ 的取值范围为\blankline[4.4em].

\ifanswer
\solbegin
\step{①当 $a-2=0$，即 $a=2$ 时，不等式为 $-4<0$ 恒成立，符合题意；}
\step{②当 $a\neq2$ 时，则 $\begin{cases}a-2<0\\ \Delta=4(a-2)^2+16(a-2)<0\end{cases}$，解得 $-2<a<2$.}
\step{综上所述，实数 $a$ 的取值范围是 $(-2,2]$.}
\solend
\fi

\item 已知 $-1\leqslant a+b\leqslant4,2\leqslant a-b\leqslant3$，则 $3a-2b$ 的取值范围为\blankline[4.4em].

\ifanswer
\solbegin
\step{设 $3a-2b=m(a+b)+n(a-b)$，则 $m+n=3,m-n=-2$，}
\step{所以 $m=\dfrac12,n=\dfrac52$，所以 $3a-2b=\dfrac12(a+b)+\dfrac52(a-b)\in\left[\dfrac92,\dfrac{19}2\right]$.}
\solend
\fi

\item 设 $x_1$、$x_2$ 是关于 $x$ 的方程 $2x^2-4mx+2m^2+3m-2=0$ 的两个实数根，则 $\lvert x_1-x_2\rvert$ 的最小值为\blankline[4.4em].

\ifanswer
\solbegin
\step{$x_1$、$x_2$ 是方程 $2x^2-4mx+2m^2+3m-2=0$ 的两个实数根，}
\step{所以 $\Delta=(4m)^2-8(2m^2+3m-2)\geqslant0$，解得 $m\leqslant\dfrac23$，}
\step{由韦达定理得 $x_1+x_2=2m,x_1\cdot x_2=\dfrac{2m^2+3m-2}2$，}
\step{则 $\lvert x_1-x_2\rvert=\sqrt{(x_1+x_2)^2-4x_1\cdot x_2}=\sqrt{4-6m}$，}
\step{当 $m=\dfrac23$ 时，$\lvert x_1-x_2\rvert$ 的最小值为 $0$.}
\solend
\fi

\item 设 $x\in R$，$[x]$ 表示不大于 $x$ 的最大整数，如 $[\pi]=3$，$[-1.2]=-2$，$\left[\dfrac12\right]=0$，则使 $\lvert\lvert x^2-1\rvert\rvert=3$ 成立的 $x$ 的取值范围是\blankline[4.4em].

\ifanswer
\solbegin
\step{若 $\lvert\lvert x^2-1\rvert\rvert=3$，则 $3\leqslant\lvert x^2-1\rvert<4$，即 $4\leqslant x^2<5$，}
\step{解得 $2\leqslant x<\sqrt5$ 或 $-\sqrt5<x\leqslant-2$，故 $x$ 的取值范围是 $(-\sqrt5,-2]\cup[2,\sqrt5)$.}
\solend
\fi

\item 设 $a$ 为实数，若关于 $x$ 的一元一次不等式组 $\begin{cases}2x+a>0\\ 3x-6a<0\end{cases}$ 的解集中有且仅有 $4$ 个整数，则 $a$ 的取值范围是\blankline[4.4em].

\ifanswer
\solbegin
\step{一元一次不等式组 $\begin{cases}2x+a>0\\ 3x-6a<0\end{cases}\Leftrightarrow-\dfrac a2<x<2a$，且 $a>0$，}
\step{关于 $x$ 的一元一次不等式组 $\begin{cases}2x+a>0\\ 3x-6a<0\end{cases}$ 的解集中有且仅有 $4$ 个整数，}
\step{①当 $4$ 个整数为 $0$，$1$，$2$，$3$ 时，则 $\begin{cases}-1\leqslant-\dfrac a2<0\\ 3<2a\leqslant4\end{cases}$，所以 $\dfrac32<a\leqslant2$，}
\step{②当 $4$ 个整数为 $-1$，$0$，$1$，$2$ 时，则 $\begin{cases}-2\leqslant-\dfrac a2<-1\\ 2<2a\leqslant3\end{cases}$，无解，}
\step{综上，$a$ 的取值范围是 $\left(\dfrac32,2\right]$.}
\solend
\fi

\item 集合 $M$、$N$、$S$ 都是非空集合，现规定如下运算：

$M\odot N\odot S=\{x\mid x\in(M\cap N)\cup(N\cap S)\cup(S\cap M)$ 且 $x\notin M\cap N\cap S\}$．假设集合 $A=\{x\mid a<x<b\}$，$B=\{x\mid c<x<d\}$，$C=\{x\mid e<x<f\}$，其中实数 $a$、$b$、$c$、$d$、$e$、$f$ 满足：$a<c<e<0<b<d<f$．计算 $A\odot B\odot C=\blankline[4.4em]$.

\ifanswer
\solbegin
\step{因为集合 $A=\{x\mid a<x<b\}$，$B=\{x\mid c<x<d\}$，$C=\{x\mid e<x<f\}$，}
\step{又因为 $a<c<e<0<b<d<f$，}
\step{所以 $A\cap B=\{x\mid c<x<b\}$，$B\cap C=\{x\mid e<x<d\}$，$C\cap A=\{x\mid e<x<b\}$，}
\step{$A\cap B\cap C=\{x\mid e<x<b\}$，故 $A\odot B\odot C=\{x\mid c<x\leqslant e$ 或 $b\leqslant x<d\}$.}
\solend
\fi

\end{enumerate}

\bigsec{二、选择题}

\begin{enumerate}
\setcounter{enumi}{12}

\item “$x\neq1$”是“$x^2-1\neq0$”的\mbox{（\quad）}条件.

\keepwithprev
A. 充分非必要\qquad B. 必要非充分

C. 充要\qquad D. 既非充分也非必要

\ifanswer
\solbegin
\step{$x^2-1\neq0\Leftrightarrow x\neq1$ 且 $x\neq-1\Rightarrow x\neq1$.}
\step{反之不成立，如 $x=-1$ 时，$x\neq1$，但 $x^2-1=0$.}
\step{因此“$x\neq1$”是“$x^2-1\neq0$”的必要非充分条件，故选 $B$.}
\solend
\fi

\item 若 $\dfrac1a<\dfrac1b<0$，则下列结论中不正确的是\mbox{（\quad）}

\keepwithprev
A. $a+b<0$\qquad B. $a^2<b^2$\qquad C. $ab<b^2$\qquad D. $ab<a^2$

\ifanswer
\solbegin
\step{由 $\dfrac1a<\dfrac1b<0$ 得 $b<a<0$，因此 $a+b<0$，且 $a^2<ab<b^2$，}
\step{所以 A、B、C 正确，D 不正确．故选 $D$.}
\solend
\fi

\item 非空集合 $X=\{x\mid a+1\leqslant x\leqslant3a-5\}$，$Y=\{x\mid1\leqslant x\leqslant16\}$，使得 $X\subseteq(X\cap Y)$ 成立的所有 $a$ 的集合是\mbox{（\quad）}

\keepwithprev
A. $\{a\mid3\leqslant a\leqslant7\}$\qquad B. $\{a\mid0\leqslant a\leqslant7\}$\qquad C. $\{a\mid3<a\leqslant7\}$\qquad D. $\{a\mid a\leqslant7\}$

\ifanswer
\solbegin
\step{因为 $X\subseteq(X\cap Y)$ 成立，所以 $X\subseteq Y$ 成立.}
\step{$X\neq\varnothing$，要使 $X\subseteq Y$，则 $\begin{cases}a\geqslant3\\ a+1\geqslant1\\ 3a-5\leqslant16\end{cases}$，即 $\begin{cases}a\geqslant3\\ a\geqslant0\\ a\leqslant7\end{cases}$，所以 $3\leqslant a\leqslant7$，故选 $A$.}
\solend
\fi

\item 已知 $M$ 是满足下列条件的集合：①$0\in M$，$1\in M$；②若 $x$、$y\in M$，则 $x-y\in M$；③若 $x\in M$ 且 $x\neq0$，则 $\dfrac1x\in M$．则下列说法正确的个数为\mbox{（\quad）}

（1）$\dfrac12\in M$；（2）$x+y\in M$；（3）$x^2\in M$．

\keepwithprev
A. $0$\qquad B. $1$\qquad C. $2$\qquad D. $3$

\ifanswer
\solbegin
\step{（1）$\dfrac12\in M$ 正确，理由如下：}
\step{由①得 $0\in M$，$1\in M$，由②得 $0-1=-1\in M$，$1-(-1)=2\in M$，}
\step{由③得 $\dfrac12\in M$；}
\step{（2）由①得 $0\in M$，由题意 $y\in M$，}
\step{所以由②得 $0-y=-y\in M$，又 $x\in M$，所以 $x-(-y)=x+y\in M$；}
\step{（3）$x\in M$，由②得 $x-1\in M$，由③得 $\dfrac1x\in M$，$\dfrac1{x-1}\in M$，}
\step{所以 $\dfrac1x-\dfrac1{x-1}\in M$，即 $\dfrac1{x(x-1)}\in M$，所以 $x(x-1)\in M$，}
\step{由（2）结论得 $x(x-1)+x\in M$，即 $x^2\in M$，}
\step{所以说法正确的个数为 $3$．故选 $D$.}
\solend
\fi

\end{enumerate}

\bigsec{三、解答题}

\begin{enumerate}
\setcounter{enumi}{16}

\item 已知 $a>b$，用比较法证明：$a^3>b^3$.

\ansspace*[8]

\ifanswer
\solbegin
\step{$a>b$，$a^3-b^3=(a-b)(a^2+ab+b^2)=(a-b)\left[\left(a+\dfrac b2\right)^2+\dfrac34b^2\right]>0$，}
\step{故 $a^3>b^3$.}
\solend
\fi

\item 求下列关于 $x$ 的不等式的解集：

（1）$\dfrac{x^2(x+3)(x^2-x+1)}{x^2-5x+4}>0$；

（2）$\dfrac{ax-1}{x-2}\geqslant0\left(a<\dfrac12\right)$.

\ansspace*[12]

\ifanswer
\solbegin
\stepm{（1）}{分母 $x^2-5x+4=(x-1)(x-4)$，要求 $x\neq1,4$；}
\step{又 $x=0$ 时原不等式不成立，且 $x^2-x+1=\left(x-\dfrac12\right)^2+\dfrac34>0$，}
\step{因此原不等式等价于 $\dfrac{x+3}{(x-1)(x-4)}>0,x\neq0$.}
\step{在数轴上标出 $-3,1,4$，分段判号得 $-3<x<1$ 或 $x>4$，}
\step{再排除 $x=0$，解集为 $(-3,0)\cup(0,1)\cup(4,+\infty)$.}
\stepm{（2）}{原不等式要求 $x\neq2$，按 $a$ 的符号分类讨论.}
\step{①当 $a<0$ 时，$\dfrac1a<2$，原不等式等价于 $\dfrac{x-\dfrac1a}{x-2}\leqslant0\Leftrightarrow\dfrac1a\leqslant x<2$，}
\step{解集为 $\left[\dfrac1a,2\right)$.}
\step{②当 $a=0$ 时，原不等式为 $\dfrac{-1}{x-2}\geqslant0\Leftrightarrow x<2$，解集为 $(-\infty,2)$.}
\step{③当 $0<a<\dfrac12$ 时，$\dfrac1a>2$，}
\step{原不等式等价于 $\dfrac{x-\dfrac1a}{x-2}\geqslant0\Leftrightarrow x<2$ 或 $x\geqslant\dfrac1a$，解集为 $(-\infty,2)\cup\left[\dfrac1a,+\infty\right)$.}
\solend
\fi

\item 集合 $A=\{x\mid-2\leqslant x\leqslant5\}$，$B=\{x\mid m+1\leqslant x\leqslant2m-1\}$.

（1）若 $B\subseteq A$，求实数 $m$ 的取值范围；

（2）若 $A\cap B\neq\varnothing$，求实数 $m$ 的取值范围.

\ansspace*[10]

\ifanswer
\solbegin
\stepm{（1）}{当 $B=\varnothing$ 时，$m+1>2m-1$，解得 $m<2$，}
\step{当 $B\neq\varnothing$ 时，$\begin{cases}m+1\leqslant2m-1\\ m+1\geqslant-2\\ 2m-1\leqslant5\end{cases}$，解得 $2\leqslant m\leqslant3$，}
\step{综上，实数 $m$ 的取值范围是 $(-\infty,3]$；}
\stepm{（2）}{当 $B=\varnothing$ 时，$m+1>2m-1$，解得 $m<2$，}
\step{当 $B\neq\varnothing$ 时，$\begin{cases}m+1\leqslant2m-1\\ 2m-1<-2\end{cases}$ 或 $\begin{cases}m+1\leqslant2m-1\\ m+1>5\end{cases}$，解得 $m\in\varnothing$ 或 $m>4$.}
\step{所以当 $A\cap B\neq\varnothing$ 时，实数 $m$ 的取值范围是 $[2,4]$.}
\solend
\fi

\item 在解决实际问题时，往往会有不同的思路和方法，这些方法有些正确，有些错误；有些简洁，有些复杂.

问题①：设 $a\in R$，集合 $M=\{x\mid(x-1)(x-a)=0\}$，若 $x\in M$ 是 $x\in\{1,2,3\}$ 的充分条件，求：$a$ 的取值集合.

问题②：设 $x\in Z$，$a=x+3$，$b=18-x$，若 $c\in Z$，求证：$c-a$ 和 $c-b$ 至少有一个数是奇数

（1）小明在解决问题①，他认为原问题等价于 $\{1,a\}\subseteq\{1,2,3\}$，解得 $a$ 的取值集合为 $\{2,3\}$，张老师判断小明解题错误，请计算正确的 $a$ 的取值集合；

（2）小红认为既然 $c\in Z$，只需根据 $c$ 是奇数还是偶数，分类讨论即可；小华则认为可以使用反证法解决问题，请你选一种你认为更好的方法并证明.

\ansspace*[12]

\ifanswer
\solbegin
\stepm{（1）}{当 $a=1$ 时，集合 $M=\{x\mid(x-1)^2=0\}=\{1\}$，}
\step{此时 $x\in M\Rightarrow x\in\{1,2,3\}$，即 $x\in M$ 是 $x\in\{1,2,3\}$ 的充分条件.}
\step{当 $a\neq1$ 时，集合 $M=\{x\mid(x-1)(x-a)=0\}=\{1,a\}$，}
\step{由 $x\in M$ 是 $x\in\{1,2,3\}$ 的充分条件，得 $\{1,a\}\subseteq\{1,2,3\}$，}
\step{所以 $a=2$ 或 $a=3$.}
\step{综上，实数 $a$ 的取值集合为 $\{1,2,3\}$；}
\stepm{（2）}{利用反证法证明．证明如下：}
\step{假设 $c-a$、$c-b$ 均为偶数，则 $(c-b)-(c-a)=a-b$ 为偶数.}
\step{因为 $a-b=x+3-18+x=2x-15$ 为奇数，}
\step{与 $b-a$ 为偶数相矛盾，所以原假设不成立.}
\step{故 $c-a$、$c-b$ 至少有一个数是奇数.}
\solend
\fi

\item 已知实数 $x_1$、$x_2$、$x_3$、$x_4$、$x_5$，满足 $x_1+x_2+x_3+x_4+x_5=5$.

（1）求证：$x_1$、$x_2$、$x_2$、$x_4$、$x_5$ 中至少有一个实数不小于 $1$；% 原卷把 x_3 误印作 x_2，照录

（2）设 $x_1$、$x_2$、$x_3$、$x_4$、$x_5$ 这五个实数两两不等，集合 $A=\{x_1,x_2,x_3,x_4,x_5\}$，若 $B\neq\varnothing$ 且 $B\subseteq A$，记 $G(B)$ 是 $B$ 中所有元素之和，对所有的 $B$，求 $G(B)$ 的平均值.

\ansspace*[13]

\ifanswer
\solbegin
\stepm{（1）}{假设 $x_1$、$x_2$、$x_2$、$x_4$、$x_5$ 中的每一个数都小于 $1$，}% 原卷 x_3 误印作 x_2，照录
\step{得 $x_1+x_2+x_2+x_4+x_5<5$，这与 $x_1+x_2+x_3+x_4+x_5=5$ 矛盾，}% 前半句原卷 x_3 误印作 x_2，照录
\step{故 $x_1$、$x_2$、$x_2$、$x_4$、$x_5$ 中至少有一个实数不小于 $1$；}% 原卷 x_3 误印作 x_2，照录
\stepm{（2）}{集合 $A=\{x_1,x_2,x_3,x_4,x_5\}$ 的非空子集个数为 $32-1=31$，}
\step{由于 $G(B)$ 是 $B$ 中所有元素之和，}
\step{得 $G(B)=16(x_1+x_2+x_3+x_4+x_5)=16\times5=80$，}
\step{则 $G(B)$ 的平均值为 $\dfrac{80}{31}$.}
\solend
\fi

\end{enumerate}

\end{document}
